% !TeX program = lualatex
% !BIB program = biber
% !TeX TXS-program:bibliography = txs:///biber
\documentclass[a4paper,10pt]{article}

% Options: twocolumn/onecolumn, final/review,
% and nobibliography/bibliography.
\usepackage[onecolumn,final,bibliography]{template}

% !TeX root = paper.tex
% Edit this file for each new paper.

\newcommand{\papertitle}{Small Language Models as the Substrate of Agentic AI: A Survey of Architectures, Compression, Deployment, and Evaluation Practice}
\newcommand{\paperauthorsmetadata}{Mohammed Gamal Ragab, Hitham Alhussian, Said Jadid Abdulkadir, W. Eltayeb Ahmed}
\newcommand{\paperkeywords}{Small language models \sep Agentic AI \sep Model compression \sep Edge AI \sep On-device inference \sep LLM--SLM collaboration}

% ORCID iDs
\newcommand{\orcidauthorA}{0000-0001-5226-4962} % M.G.R.
\newcommand{\orcidauthorB}{0000-0003-0038-3702} % S.J.A.

\title{\papertitle}
\author[1,2*]{Mohammed Gamal Ragab\,\orcidlink{\orcidauthorA}}
\author[1,2*]{Hitham Alhussian}
\author[1,2]{Said Jadid Abdulkadir\,\orcidlink{\orcidauthorB}}
\author[4]{W. Eltayeb Ahmed}

\affil[1]{Department of Computer and Information Sciences, Universiti Teknologi PETRONAS, Seri Iskandar 32610, Perak, Malaysia}
\affil[2]{Center for Research in Data Science (CeRDaS), Universiti Teknologi PETRONAS, Seri Iskandar 32610, Perak, Malaysia}
\affil[4]{Department of Mathematics and Statistics, College of Science, Imam Mohammad Ibn Saud Islamic University (IMSIU), Riyadh, Saudi Arabia}
\affil[*]{Corresponding authors: \href{mailto:mohammed.ragab@utp.edu.my}{mohammed.ragab@utp.edu.my} (M.G.R.); \href{mailto:seddig.alhussian@utp.edu.my}{seddig.alhussian@utp.edu.my} (H.A.)}
\date{}

\hypersetup{
  pdftitle={\papertitle},
  pdfauthor={\paperauthorsmetadata},
  pdfsubject={Journal article manuscript},
  pdfkeywords={\paperkeywords}
}

% abbreviations:

% General AI/ML
\newacronym{ai}{AI}{Artificial Intelligence}
\newacronym{ml}{ML}{Machine Learning}
\newacronym{dl}{DL}{Deep Learning}
\newacronym{nlp}{NLP}{Natural Language Processing}
\newacronym{llm}{LLM}{Large Language Model}
\newacronym{slm}{SLM}{Small Language Model}

% Neural Network Types
\newacronym{nn}{NN}{Neural Network}
\newacronym{dnn}{DNN}{Deep Neural Network}
\newacronym{cnn}{CNN}{Convolutional Neural Network}
\newacronym{rnn}{RNN}{Recurrent Neural Network}
\newacronym{lstm}{LSTM}{Long Short-Term Memory}

% Transformer Architecture
\newacronym{ssm}{SSM}{State Space Model}
\newacronym{mha}{MHA}{Multi-Head Attention}
\newacronym{mqa}{MQA}{Multi-Query Attention}
\newacronym{gqa}{GQA}{Grouped-Query Attention}
\newacronym{swa}{SWA}{Sliding Window Attention}
\newacronym{rope}{RoPE}{Rotary Position Embeddings}
\newacronym{ffn}{FFN}{Feed-Forward Network}

% Training and Fine-tuning
\newacronym{lora}{LoRA}{Low-Rank Adaptation}
\newacronym{qlora}{QLoRA}{Quantized Low-Rank Adaptation}
\newacronym{sft}{SFT}{Supervised Fine-Tuning}
\newacronym{rlhf}{RLHF}{Reinforcement Learning from Human Feedback}
\newacronym{grpo}{GRPO}{Group Relative Policy Optimization}
\newacronym{kd}{KD}{Knowledge Distillation}
\newacronym{sgd}{SGD}{Stochastic Gradient Descent}
\newacronym{peft}{PEFT}{Parameter-Efficient Fine-Tuning}
\newacronym{rag}{RAG}{Retrieval-Augmented Generation}
\newacronym{cot}{CoT}{Chain-of-Thought}

% Quantization
\newacronym{gptq}{GPTQ}{Generative Pre-Trained Transformer Quantization}
\newacronym{awq}{AWQ}{Activation-Aware Weight Quantization}
\newacronym{gguf}{GGUF}{GPT-Generated Unified Format}
\newacronym{nf4}{NF4}{NormalFloat 4-bit}
\newacronym{ptq}{PTQ}{Post-Training Quantization}
\newacronym{qat}{QAT}{Quantization-Aware Training}

% Benchmarks
\newacronym{mmlu}{MMLU}{Massive Multitask Language Understanding}
\newacronym{gsm8k}{GSM8K}{Grade School Math 8K}
\newacronym{arc}{ARC}{AI2 Reasoning Challenge}
\newacronym{piqa}{PIQA}{Physical Intuition Question Answering}
\newacronym{mnli}{MNLI}{Multi-Genre Natural Language Inference}
\newacronym{anli}{ANLI}{Adversarial Natural Language Inference}
\newacronym{mbpp}{MBPP}{Mostly Basic Python Problems}
\newacronym{squad}{SQuAD}{Stanford Question Answering Dataset}

% Metrics
\newacronym{em}{EM}{Exact Match}
\newacronym{ttft}{TTFT}{Time-to-First-Token}
\newacronym{tps}{TPS}{Tokens Per Second}
\newacronym{vram}{VRAM}{Video Random Access Memory}
\newacronym{ram}{RAM}{Random Access Memory}
\newacronym{tdp}{TDP}{Thermal Design Power}

% Hardware/Deployment
\newacronym{gpu}{GPU}{Graphics Processing Unit}
\newacronym{npu}{NPU}{Neural Processing Unit}
\newacronym{cpu}{CPU}{Central Processing Unit}
\newacronym{moe}{MoE}{Mixture of Experts}

% Frameworks
\newacronym{onnx}{ONNX}{Open Neural Network Exchange}

% Other
\newacronym{ci}{CI}{Computational Intelligence}
\newacronym{kl}{KL}{Kullback--Leibler}
\newacronym{stem}{STEM}{Science, Technology, Engineering, and Mathematics}


\setcounter{secnumdepth}{4}
\setcounter{tocdepth}{4}

\begin{document}

\MakeJournalFrontMatter{
  \begin{abstract}
  Large Language Models (LLMs) with hundreds of billions of parameters
  have advanced natural language processing dramatically, but their cost,
  memory demand, and latency make them a poor fit for the fastest-growing
  consumer of language intelligence: the AI agent, which issues many
  narrow, repetitive, well-specified model calls rather than a few
  open-ended ones. Small Language Models (SLMs)---here, decoder-only or
  hybrid transformers of roughly twenty billion parameters or fewer,
  counting active parameters for mixture-of-experts models---have closed
  much of the capability gap through knowledge distillation, synthetic-data
  curation, and post-training quantization, and are now a natural default
  for agentic workloads. This paper is a combined survey and position
  statement. We catalog 126 publicly documented SLM variants (2022--2026),
  including the 2025--2026 wave of hybrid state-space, unified-reasoning,
  and mobile-first multimodal families, and organize them along a
  four-dimensional taxonomy of \emph{architecture}, \emph{training
  strategy}, \emph{modality}, and \emph{application domain}. We review
  compression for the sub-10B regime, where quantization consistently
  outperforms pruning. We then audit the public evaluation record for an
  18-model cohort and report a negative result that bears on how the field
  compares these models: a cross-model accuracy table cannot be assembled
  from current public reporting without misrepresenting its sources,
  because vendors report non-intersecting benchmark sets, shot counts that
  differ per benchmark and per vendor, base and instruction-tuned figures
  that are routinely conflated, and few matched baselines. Building on the
  evidence that does hold, we
  argue that SLMs are the substrate of agentic AI: they are sufficiently
  capable for the sub-tasks agents delegate, roughly 10--30$\times$
  cheaper to serve, cheaper to specialize, and easier to constrain to tool
  interfaces. We survey heterogeneous SLM--LLM routing and the
  LLM-to-SLM conversion pipeline, and close with the open problems---%
  reasoning under parameter constraints, hallucination, hardware
  co-design, fairness under compression, and empirical verification of the
  consolidated evidence---that stand between the present and that future.
  \end{abstract}

  \begin{highlights}
  \item Surveys 126 documented small language model variants (2022--2026), including the 2025--2026 hybrid state-space, unified-reasoning, and mobile-multimodal waves, under a taxonomy of architecture, training strategy, modality, and application domain.
  \item Audits the public evaluation record for 18 SLMs and shows that it does not support the cross-model accuracy comparisons routinely published in surveys, leaderboards, and secondary sources.
  \item Reviews compression for the sub-10B regime and the finding that quantization consistently outperforms pruning below 4B parameters.
  \item Argues that SLMs are the substrate of agentic AI on grounds of capability sufficiency, 10--30$\times$ lower serving cost, cheaper specialization, and tighter behavioral alignment to tool interfaces.
  \item Surveys heterogeneous SLM--LLM routing and the LLM-to-SLM conversion pipeline, and sets an empirical-verification agenda for the consolidated claims.
  \end{highlights}

  \begin{keywords}
  \paperkeywords
  \end{keywords}
}

% acronyms may have been expanded in the abstract; reset so the body re-expands
% each one on its own first use.
\glsresetall

\section{Introduction}
\label{section:introduction}

The landscape of \gls{nlp} has been reshaped by \glspl{llm}. Systems such as GPT-4 \cite{openai2023gpt4}, Claude \cite{anthropic2024claude}, Gemini \cite{gemini2023}, and the LLaMA series \cite{grattafiori2024llama3} demonstrate broad competence across reasoning, code generation, and multilingual tasks. That competence, however, comes at extraordinary cost. GPT-3's 175~billion parameters \cite{brown2020gpt3} inaugurated an era in which scale was treated as synonymous with capability, and scaling laws \cite{kaplan2020scaling,hoffmann2022chinchilla} appeared to promise predictable returns on ever-larger models. Frontier \glspl{llm} now require multi-thousand-GPU clusters to train, tens of gigabytes of accelerator memory to serve, and cloud infrastructure that introduces latency, recurring cost, and privacy exposure. Training GPT-3 is estimated to have consumed 1{,}287~MWh and emitted roughly 552~metric tons of CO$_2$ equivalent \cite{patterson2021carbon}; a single frontier-model query draws on the order of ten times the energy of a conventional web search \cite{devries2023energy}.

\subsection{The Rise of Small Language Models}
\label{subsec:intro_rise}

Running against this trajectory is a counter-trend: \glspl{slm} that achieve disproportionately strong performance for their size. We use \emph{\gls{slm}} to mean a decoder-only or hybrid transformer with roughly twenty billion parameters or fewer---counting \emph{active} parameters for mixture-of-experts models, since deployment cost tracks those---small enough to be plausibly served on a single consumer GPU, a laptop, or a modern phone. The empirical cohort of Section~\ref{sec:reported} is concentrated below 10B, with Phi-4 (14B) retained as an upper-bound reference. Where earlier compact models such as DistilBERT \cite{sanh2019distilbert} and TinyBERT \cite{jiao2020tinybert} were task-specific encoders, the current generation are general-purpose, instruction-following systems. The Chinchilla scaling analysis \cite{hoffmann2022chinchilla} showed that most large models were badly undertrained and that a smaller model given proportionally more data could match them; the Phi family then demonstrated that ``textbook-quality'' synthetic data lets a 3.8B model rival models an order of magnitude larger \cite{abdin2024phi3}, Gemma applied distillation from larger teachers \cite{gemma2024}, and Qwen2.5 pushed sub-7B training to 18~trillion tokens \cite{qwen2024qwen25}. MIT Technology Review named \glspl{slm} one of its 10 Breakthrough Technologies of 2025 \cite{mit2025breakthrough}. By 2026 the strongest 3--4B models post benchmark scores that were state-of-the-art for 13B-class models in 2023, and the 2025--2026 release wave adds hybrid state-space architectures \cite{falconh1_2025}, controllable-reasoning models \cite{qwen3_2025}, and mobile-first multimodal families \cite{gemma4_2026,gunter2025afm}. Figure~\ref{fig:cumulative_releases} shows the resulting acceleration: the cumulative count of documented open \gls{slm} variants has grown from a handful in 2023 to more than one hundred by early 2026.

\begin{figure}[t]
  \centering
  \includegraphics[width=\linewidth]{paper_combined_cumulative_releases.pdf}
  \caption{Cumulative count of documented small (SLM) and large (LLM) language-model variants by release date, from the model inventory of Appendix~\ref{appendix:ecosystem} (snapshot September~2026). SLM releases accelerate sharply from mid-2024.}
  \label{fig:cumulative_releases}
\end{figure}

\subsection{The Agentic Turn}
\label{subsec:intro_agentic}

At the same time, the way language models are consumed has shifted. The fastest-growing workload is not the chat session but the \emph{AI agent}: a control loop in which a model is called repeatedly to choose a tool, fill its arguments, extract a field, classify an intent, or check a result \cite{belcak2025slm_agentic}. These calls are narrow, repetitive, and format-constrained, and they exercise only a thin slice of a frontier model's competence. \citet{belcak2025slm_agentic} argue---and this paper concurs and elaborates---that for such workloads \glspl{slm} are not a compromise but the appropriate default: sufficiently capable for the delegated sub-task, 10--30$\times$ cheaper to serve, cheaper to specialize, and easier to constrain to a code interface, with a strong generalist model reserved for the minority of calls that genuinely need planning or broad context. Surveys of LLM--SLM collaboration \cite{chen2024role_small,chen2025collab_survey} and production systems such as Apple's on-device/server split \cite{gunter2025afm} indicate that this heterogeneous design is already the direction of travel.

\subsection{Scope and Contributions}
\label{subsec:intro_gap}

Existing \gls{slm} surveys \cite{lu2024slmsurvey,vannguyen2024slm_survey,wang2024comprehensive} provide taxonomies and measurements but predate the 2025--2026 model wave and treat agentic deployment only in passing; agentic position pieces \cite{belcak2025slm_agentic} make the argument but do not ground it in a consolidated capability-and-deployment review. This paper joins the two. It is simultaneously a survey of the \gls{slm} landscape and a position statement that \glspl{slm} are the substrate of agentic AI. Its contributions are:

\begin{description}
  \item[Landscape and taxonomy] A catalog of 126 documented \gls{slm} variants (2022--2026), including the recent hybrid, reasoning, and mobile-multimodal families, organized along a four-dimensional taxonomy of architecture, training strategy, modality, and application domain (Section~\ref{sec:taxonomy} and Appendix~\ref{appendix:ecosystem}).
  \item[Compression review] A synthesis of distillation, quantization, pruning, and compound pipelines for the sub-10B regime, including the finding that quantization consistently outperforms pruning below 4B parameters \cite{kurz2025pruning_vs_quant} (Section~\ref{sec:compression}).
  \item[Consolidated reported performance] A comparison of 18 \glspl{slm} across 12 tasks, three evaluation paradigms, six hardware tiers, and three precisions, assembled under a documented grey-literature protocol in which every figure is a \emph{reported} value traced to a dated source rather than an original run (Section~\ref{sec:reported}, Appendix~\ref{app:provenance}).
  \item[Deployment ecosystem] A survey of inference frameworks, quantization formats, and on-device platforms from Raspberry~Pi to Apple silicon (Section~\ref{sec:deployment}).
  \item[The agentic argument] A structured case that \glspl{slm} are the default execution substrate for agents, covering capability sufficiency, economics, behavioral alignment, heterogeneous SLM--LLM routing, and the LLM-to-SLM conversion pipeline (Section~\ref{sec:agentic}).
  \item[Open problems and verification] An agenda of open research problems and targeted empirical-verification opportunities for the consolidated evidence base (Section~\ref{sec:challenges}).
\end{description}

\subsection{Survey Method and Sources}
\label{subsec:intro_method}

This survey draws on three evidence streams, each with its own protocol.
\emph{(1)~The model inventory} (Appendix~\ref{appendix:ecosystem}) is generated
programmatically from \texttt{merged\_llm\_database.json}, a curated snapshot
(September~2026) compiled from vendor technical reports, official model cards, and
dated release announcements; each catalogued variant links to a primary source.
\emph{(2)~The technique literature} on compression, training, and LLM--SLM
collaboration was gathered by keyword search of the ACL Anthology, arXiv (cs.CL,
cs.LG), and the proceedings of NeurIPS, ICML, ICLR, and MLSys for 2019--2026,
seeded from the four prior \gls{slm} surveys
\citep{lu2024slmsurvey,vannguyen2024slm_survey,wang2024comprehensive,belcak2025slm_agentic}
and expanded by forward and backward citation tracing; a work was retained when it
concerned models at or below roughly 20B parameters, or a compression or
deployment method material to that regime. \emph{(3)~The reported-performance
figures} of Section~\ref{sec:reported} follow the grey-literature collection and
reconciliation protocol of Section~\ref{subsec:reported_protocol}. The survey is
narrative rather than systematic in the PRISMA sense: the goal is coverage of the
design space and of the public reported record, not an exhaustive enumeration of
every publication.

\subsection{Paper Organization}
\label{subsec:intro_organization}

Section~\ref{section:background} reviews prior work on the LLM-to-SLM transition, compression, edge inference, and LLM--SLM collaboration. Section~\ref{sec:taxonomy} presents the four-dimensional taxonomy. Section~\ref{sec:compression} reviews compression techniques. Section~\ref{sec:reported} documents the data-collection protocol and consolidates reported performance. Section~\ref{sec:deployment} covers resource-constrained deployment. Section~\ref{sec:agentic} makes the agentic argument. Section~\ref{sec:applications} surveys domain-specific applications. Section~\ref{section:discussion} discusses implications, limitations, and ethics. Section~\ref{sec:challenges} sets out open challenges and an empirical-verification agenda. Section~\ref{section:conclusion} concludes.

The ecosystem inventory referenced throughout is generated directly from \texttt{merged\_llm\_database.json} (snapshot September~2026). Appendix~\ref{appendix:ecosystem} provides the full variant inventory grouped chronologically by release year and family, together with the source-selection protocol.

\section{Related Work and Research Gaps}
\label{section:background}

\subsection{From Scaling to Small Models}
\label{subsec:background_evolution}
\label{subsec:background_techniques}
\label{subsec:background_families}

The Transformer \citep{vaswani2017attention} and the pre-train-then-adapt paradigm \citep{devlin2019bert,brown2020gpt3} led to a scaling race that the Chinchilla study \citep{hoffmann2022chinchilla} partly redirected: a 70B model trained on 1.4~trillion tokens outperformed the 280B Gopher, implying that most large models were undertrained and that smaller, longer-trained models could match them. Subsequent families pushed this further at small scale through curated and synthetic data \citep{gunasekar2023phi1,abdin2024phi3}, distillation from larger teachers \citep{gemma2024,team2025gemma3}, and pruning plus distillation from a larger sibling, as for Llama~3.2's 1B and 3B edge models \citep{meta2024llama32}. Three bodies of technique underlie the current generation: distillation and synthetic-data training \citep{hinton2015distilling,gunasekar2023phi1}; post-training compression and parameter-efficient adaptation \citep{frantar2023gptq,lin2024awq,hu2022lora,dettmers2023qlora}; and on-device runtimes \citep{gerganov2023llamacpp,meta2024executorch,mlc2024}. Sections~\ref{sec:compression} and~\ref{sec:deployment} review them, and \suppsec~\ref{supp:families} and~\ref{supp:background} summarize the principal families and the longer background.

\subsection{LLM--SLM Collaboration and Agentic Systems}
\label{subsec:background_collab}

A line of work studies how small and large models cooperate \citep{chen2024role_small,chen2025collab_survey}. Cascades and routers are the most relevant to agents: FrugalGPT \citep{chen2023frugalgpt} learns a cascade over commercial APIs, Hybrid~LLM \citep{ding2024hybridllm} routes by predicted difficulty, AutoMix \citep{aggarwal2024automix} escalates on self-verification, RouteLLM \citep{ong2024routellm} learns from preference data, RouterBench \citep{hu2024routerbench} benchmarks routers on 405K inference outcomes, and Division-of-Thoughts \citep{division_of_thoughts_2025} routes sub-questions between an on-device and a cloud model. Speculative decoding \citep{leviathan2023speculative,chen2023specsampling} uses a small model to accelerate a large one without delegating answers. On the agentic side, function-calling data pipelines report specialized sub-10B models among the leading entries of the Berkeley Function-Calling Leaderboard \citep{liu2024apigen,liu2025toolace,zhang2024xlam}, and agent evaluation has begun to report cost alongside accuracy \citep{kapoor2024agentsmatter,kapoor2025hal,erol2025costofpass}.

\subsection{Gaps in Existing Literature}
\label{subsec:background_gaps}

General \gls{slm} surveys \citep{lu2024slmsurvey,vannguyen2024slm_survey,wang2024comprehensive} provide technique taxonomies, application coverage, and in the case of \citet{lu2024slmsurvey} a measurement study of 70 models; \citet{gupta2025slmsurvey}, updated in 2026, survey about 160 papers on 1--8B models. None separates active from total parameters or foregrounds agentic deployment. \citet{srivastava2025thinkslm} benchmark 72 small models on 17 reasoning benchmarks with three repeats, including models derived by quantization, pruning, and distillation, and find that reasoning tracks training method and data more than scale and that quantization preserves reasoning while pruning disrupts it; it is the closest thing in the literature to the controlled evaluation whose absence Section~\ref{sec:reported} documents, but covers reasoning only. \citet{feng2025efficient_reasoning_survey} survey efficient reasoning, and \citet{chen2024mas_survey} survey multi-agent systems by application.

Two works state the agentic case directly. \citet{belcak2025slm_agentic} argue it as a position, with a conversion algorithm and estimates of how much agent traffic small models could handle, and \citet{tian2026edgeagents} argue that personal agents must move to the edge. The closest prior work is the survey of \citet{sharma2025slm_agentic}. It reviews 1--12B (occasionally about 20B) models for agentic workloads, proposes an \gls{slm}-default, \gls{llm}-fallback architecture with uncertainty-aware routing, retries, and verifier repair (their Section~VIII), proposes cost per successful task as the deciding metric, and names where large models still win: open-domain synthesis, complex multi-hop reasoning, safety-critical judgment, and long-context work (their Section~XI). We share that design and metric. It does not separate active from total parameters, grade its evidence or report contrary results, examine whether the vendor evaluation record supports comparison, or state a decision rule under which its position would be refuted, and its estimate that cost per successful task falls 10--30$\times$ (their Section~X) is not a reported measurement. Those are the points on which this paper differs: the C1/C2 definition, the verified inventory, the quantified audit of the evaluation record, graded evidence including contrary results, a break-even condition that includes retries, checking, and escalation together with a correctness requirement on the check, and a falsification design with comparators, a decision rule, and a sample size. Table~\ref{tab:survey_compare} summarizes the positioning.

\begin{table}[t]
\centering
\caption{Coverage of this survey relative to prior surveys and position papers (Y~full, P~partial, --~absent), as judged by the present authors from each work's text (grade~\evD). \emph{Active/total}: separates active from total parameters. \emph{Evaluation critique}: analysis of whether the reporting record supports cross-model comparison. \emph{Graded evidence}: an explicit strength-of-evidence grade per claim. \emph{Agentic focus}: small models inside agents as a primary topic. \emph{Decision rule}: a stated condition under which the position is refuted.}
\label{tab:survey_compare}
\footnotesize
\adjustbox{max width=\linewidth}{%
\begin{tabular}{lccccccc}
\toprule
\textbf{Work} & \textbf{Taxonomy} & \textbf{Compression} & \textbf{Active/} & \textbf{Evaluation} & \textbf{Graded} & \textbf{Agentic} & \textbf{Decision} \\
 & & & \textbf{total} & \textbf{critique} & \textbf{evidence} & \textbf{focus} & \textbf{rule} \\
\midrule
\citet{lu2024slmsurvey}            & Y & P & -- & P & -- & -- & -- \\
\citet{vannguyen2024slm_survey}    & Y & Y & -- & -- & -- & -- & -- \\
\citet{wang2024comprehensive}      & Y & Y & -- & -- & -- & P & -- \\
\citet{gupta2025slmsurvey}         & Y & P & -- & -- & -- & -- & -- \\
\citet{srivastava2025thinkslm}     & P & P & -- & Y & -- & -- & -- \\
\citet{belcak2025slm_agentic}      & -- & -- & -- & P & -- & Y & -- \\
\citet{tian2026edgeagents}         & -- & -- & -- & -- & -- & Y & P \\
\citet{sharma2025slm_agentic}      & P & P & -- & -- & -- & Y & -- \\
\textbf{This survey}               & Y & Y & Y & Y & Y & Y & Y \\
\bottomrule
\end{tabular}%
}
\end{table}

% ================================================================
% Section: Taxonomy of Small Language Models
% Merged from main (2).tex — Section III
% ================================================================

\section{Taxonomy of Small Language Models}
\label{sec:taxonomy}

We propose a four-dimensional taxonomy of SLMs organized along orthogonal axes: \emph{architecture}, \emph{training strategy}, \emph{modality}, and \emph{application domain} (Figure~\ref{fig:taxonomy}). A given model may be positioned along all four dimensions simultaneously---Phi-3-vision, for instance, is a dense transformer (architecture), trained from scratch on synthetic data (training strategy), multimodal (modality), and general-purpose (domain).

\begin{figure}[t]
\centering
\footnotesize
\newcommand{\txax}[1]{{\bfseries #1}}
\newcommand{\txcat}[1]{{\scriptsize\bfseries #1}}
\newcommand{\txex}[1]{{\scriptsize\itshape\textcolor{black!60}{#1}}}
\begin{tikzpicture}[
  panel/.style={draw, rounded corners=2.5pt, line width=0.5pt, inner sep=5pt,
                text width=71mm, align=left, fill=black!3},
]
\node[panel] (arch) {%
  \txax{1.\ Architecture}\\[2pt]
  \txcat{Dense transformer} --- GQA, RoPE, SwiGLU\\ \txex{Phi-4-mini, Qwen3, SmolLM3, Gemma 3}\\[1pt]
  \txcat{Mixture-of-experts} --- sparse activation\\ \txex{Mixtral, Qwen3-30B-A3B, Llama 4 Scout}\\[1pt]
  \txcat{Hybrid SSM/attention} --- linear-time\\ \txex{Falcon-H1, Granite 4.0, Nemotron Nano 2}%
};
\node[panel, right=5mm of arch] (train) {%
  \txax{2.\ Training strategy}\\[2pt]
  \txcat{From scratch} --- curated / synthetic data\\ \txex{Phi, Qwen3, SmolLM3}\\[1pt]
  \txcat{Distillation} --- transfer from a larger teacher\\ \txex{Gemma 3/4, DeepSeek-R1-Distill, Phi-4-mini-reasoning}\\[1pt]
  \txcat{Pruning + distillation} --- shrink a large model\\ \txex{Minitron, Sheared-LLaMA}%
};
\node[panel, below=5mm of arch] (mod) {%
  \txax{3.\ Modality}\\[2pt]
  \txcat{Text-only} --- majority of production SLMs\\ \txex{Mistral 7B, SmolLM3, Qwen3}\\[1pt]
  \txcat{Vision--language} --- CLIP/SigLIP encoder\\ \txex{Phi-3-vision, Gemma 3, PaliGemma}\\[1pt]
  \txcat{Text + vision + audio} --- one sub-10B model\\ \txex{Gemma 3n, Gemma 4, Apple AFM}%
};
\node[panel, below=5mm of train] (dom) {%
  \txax{4.\ Application domain}\\[2pt]
  \txcat{General-purpose} --- broad benchmarks\\ \txex{Phi, Gemma, Llama, Qwen, Mistral}\\[1pt]
  \txcat{Healthcare / legal} --- on-premises, compliance\\ \txex{Meerkat, OptimCLM}\\[1pt]
  \txcat{Code / multilingual} --- domain corpora\\ \txex{StarCoder2, DeepSeek-Coder, Aya}%
};
\end{tikzpicture}
\caption{The four-dimensional taxonomy of small language models. The axes are
orthogonal: each model occupies one category on each. Representative models are
italicised. Section~\ref{sec:taxonomy} develops each axis; Table~\ref{tab:taxonomy_summary}
gives the tabular form.}
\label{fig:taxonomy}
\end{figure}

Figure~\ref{fig:releases_by_size} shows how the size distribution of new SLM releases has evolved: the sub-1B and 1--3B bands dominated early releases, but the 3--7B band has become the centre of gravity as capability-per-parameter has improved. Figure~\ref{fig:family_size_range} shows the parameter span of the major families.

\begin{figure}[t]
  \centering
  \includegraphics[width=\linewidth]{paper_slm_releases_by_size.pdf}
  \caption{Documented SLM variants by release year and parameter band. The 3--7B band has grown from a minority of releases in 2023 to the largest single band by 2025.}
  \label{fig:releases_by_size}
\end{figure}

\begin{figure}[t]
  \centering
  \includegraphics[width=\linewidth]{paper_slm_family_size_range.pdf}
  \caption{Parameter span (minimum, median, maximum, in billions, log scale) of the principal SLM families in the inventory.}
  \label{fig:family_size_range}
\end{figure}

% ----------------------------------------------------------------
\subsection{Taxonomy by Architecture}
\label{subsec:taxonomy_architecture}

\subsubsection{Dense Transformer Models}

The majority of production SLMs use a dense transformer decoder—every token attends to every other token (within window) and every parameter is activated per forward pass. This category spans from sub-1B models (SmolLM2-135M \cite{allal2024smollm}) to 7--8B flagships (Mistral 7B \cite{jiang2023mistral}, Llama 3.1 8B \cite{grattafiori2024llama3}). Recent dense families are notable for pushing capability per parameter: SmolLM3-3B \cite{smollm3} uses grouped-query attention with NoPE and is trained on 11.2T tokens with a documented three-stage curriculum, matching or exceeding Llama-3.2-3B and Qwen2.5-3B; Qwen3's dense line \cite{qwen3_2025} spans 0.6B--32B under a single GQA/RoPE/SwiGLU architecture with a unified thinking / non-thinking switch. Key architectural improvements over the vanilla GPT-2 design include:

\begin{itemize}
    \item \textbf{Grouped-Query Attention (GQA):} Introduced in LLaMA-2 \cite{touvron2023llama2} and adopted by Mistral 7B and most post-2023 SLMs. GQA reduces the key-value cache memory by a factor of $G$ (number of query heads divided by number of KV heads), enabling longer context windows without memory explosion.

    \item \textbf{RoPE with Extended Base:} Phi-3, Phi-4, Gemma 3, and LLaMA 3.2 all use Rotary Position Embeddings (RoPE) with extended base frequencies ($\theta = 500{,}000$--$1{,}000{,}000$) to support context windows up to 128K tokens without architectural modification.

    \item \textbf{SwiGLU FFN:} Replacing the standard two-layer FFN with a Gated Linear Unit variant (SwiGLU), adopted across nearly all post-2022 SLMs, providing a roughly 5\% perplexity improvement at equivalent parameter counts.
\end{itemize}

These changes, together with RoPE base extension and positional interpolation, have driven a rapid expansion of context windows in the SLM inventory (Figure~\ref{fig:context_by_year}): median supported context grew from a few thousand tokens in 2023 to over 100K by 2025.

\begin{figure}[t]
  \centering
  \includegraphics[width=\linewidth]{paper_slm_context_by_year.pdf}
  \caption{Supported context window of documented SLM variants by release year (log scale). Each violin is the kernel-density estimate of the distribution; the inner box marks the interquartile range, the dashed line the mean, and the dots individual models. Long context has become standard in the sub-10B regime.}
  \label{fig:context_by_year}
\end{figure}

\subsubsection{Mixture-of-Experts (MoE) SLMs}

MoE architectures \cite{shazeer2017moe,fedus2022switch} decouple \emph{total} parameters from \emph{active} parameters per token: only a subset of ``expert'' feed-forward network (FFN) blocks is activated by a learned router. This allows a model to have a large total parameter count while keeping per-token FLOPs comparable to a much smaller dense model.

Mixtral~$8{\times}7$B \cite{jiang2024mixtral} (47B total, 12.9B active per token) is the canonical example of an MoE SLM-boundary model, competitive with 70B dense models while activating only 25\% of parameters per token. More recent MoE SLMs include DeepSeek-MoE \cite{dai2024deepseekmoe} (finer-grained experts, 2B active from 16B total) and Qwen3-30B-A3B \cite{qwen3_2025}, which activates only 3B of 30B parameters and reportedly matches Qwen2.5's dense base models at one-tenth the active parameter count. Meta's Llama~4 Scout \cite{llama4_2025} (109B total, 17B active, 16 experts) applies the same principle at the small-\emph{active} scale, with Behemoth serving as the distillation teacher. Because deployment cost tracks \emph{active} rather than total parameters, several of these models fall inside the practical SLM envelope for latency and energy even though their storage footprint does not. Dense models still dominate the inventory numerically (Figure~\ref{fig:moe_vs_dense}), but the MoE models that exist carry a much larger median total-parameter count at comparable per-token cost.

\begin{figure}[t]
  \centering
  \includegraphics[width=\linewidth]{paper_slm_moe_vs_dense.pdf}
  \caption{Mixture-of-experts versus dense SLMs in the inventory: variant count (left), median total parameters (centre), and median \emph{active} parameters per token (right). MoE models carry a much larger total-parameter count but a per-token cost comparable to dense models.}
  \label{fig:moe_vs_dense}
\end{figure}

\subsubsection{Hybrid SSM-Transformer Models}

A nascent but rapidly growing class of SLMs combines attention-based transformer layers with State Space Models (SSMs) \cite{gu2023mamba}, which process sequences in linear rather than quadratic time. Mamba \cite{gu2023mamba} introduced hardware-aware selective SSMs; Jamba \cite{lieber2024jamba} interleaves transformer and Mamba blocks in a 1:7 ratio, achieving competitive performance with $2.6\times$ lower KV-cache memory than a pure transformer. In 2025--2026 this design moved from research to production at the small scale. Falcon-H1 \cite{falconh1_2025} spans 0.5B--34B with a parallel Mamba-2/attention block. IBM's Granite~4.0 \cite{granite4_2025} interleaves self-attention and Mamba-2 layers at roughly a 1:9 ratio and pairs it with a fine-grained MoE, reporting more than 70\% memory reduction for long-context and multi-session inference relative to transformer-only models, with ``Tiny'' ($\sim$7B total, $\sim$1B active) and ``Micro'' variants aimed explicitly at edge and local use. NVIDIA's Nemotron Nano~2 \cite{nvidia2025nemotron_nano2} replaces most self-attention in a 9--12B model with Mamba-2 layers (6 attention, 28 Mamba-2, 28 FFN across 62 layers) and reports up to 6$\times$ higher throughput than Qwen3-8B when generating long reasoning traces. Alibaba's Qwen3.5 small series adopts a Gated DeltaNet hybrid at a 3:1 linear-to-full-attention ratio \cite{qwen3_5_2026}. Linear-time inference is most valuable exactly where SLMs are deployed---long-context, streaming, memory-bound edge settings---which is why the hybrid category is growing fastest.

% ----------------------------------------------------------------
\subsection{Taxonomy by Training Strategy}
\label{subsec:taxonomy_training}

\subsubsection{Trained from Scratch}

Most commercially significant SLMs are pre-trained from scratch on carefully curated datasets. The Phi family exemplifies a data-centric philosophy: rather than increasing parameters, Microsoft researchers found that training on ``textbook-quality'' synthetic and curated web data yields disproportionate gains \cite{gunasekar2023phi1,li2023phi2}. Phi-1 (1.3B) \cite{gunasekar2023phi1} trained on $<$8B tokens—two orders of magnitude less than contemporaries—yet achieved competitive Python coding benchmarks. LLaMA-3.2 3B \cite{grattafiori2024llama3} trained on 9T tokens from Meta's curated corpus. Gemma 3 \cite{team2025gemma3} used 14T tokens with a curriculum that progressively increases the ratio of math, code, and multilingual data.

Key training innovations in this category include: (a)~\emph{data quality filtering} using heuristic and classifier-based pipelines; (b)~\emph{synthetic data generation} from larger teacher models \cite{tinystories2023,ho2022orca}; (c)~\emph{curriculum scheduling} where token mixture evolves over training to emphasize harder reasoning data in later stages; and (d)~\emph{long-context extension} via continued pre-training with positional interpolation or extended RoPE.

\subsubsection{Distilled SLMs}

Knowledge distillation \cite{hinton2015distilling} transfers capability from a large teacher model to a smaller student. In the LLM era, distillation takes several forms. \emph{White-box distillation} accesses the teacher's internal representations (intermediate hidden states, attention maps), enabling feature-level alignment. \emph{Black-box distillation} uses only the teacher's outputs (token logits or generated text) as supervision. Orca \cite{ho2022orca} demonstrated that distilling reasoning \emph{traces} from GPT-4—rather than just answers—enables a 13B student to acquire substantially stronger chain-of-thought capabilities. Reasoning distillation became a dominant 2025 strategy: DeepSeek-R1-Distill models \cite{deepseek2025r1} transfer long chain-of-thought from DeepSeek-R1 into Qwen-2.5 students at 1.5B--14B, and Phi-4-mini-reasoning \cite{xu2025phi4mini_reasoning} combines distillation from a large reasoning teacher with reinforcement learning to bring competition-mathematics performance to a 3.8B model. Gemma~3 and Gemma~4 \cite{team2025gemma3,gemma4_2026} distill every model size, including the largest, from a bigger teacher by default.

\subsubsection{Pruned and Compressed SLMs}

Rather than training a small model from scratch, pruning removes parameters from a large pre-trained model. NVIDIA's Minitron pipeline \cite{minitron2024} first prunes LLaMA-3.1-8B by removing attention heads, MLP neurons, and layers guided by activation statistics, then distills from the unpruned teacher, achieving a 40$\times$ reduction in training compute relative to training from scratch while matching or exceeding the performance of same-size models trained from scratch.

% ----------------------------------------------------------------
\subsection{Taxonomy by Modality}
\label{subsec:taxonomy_modality}

\subsubsection{Text-Only SLMs}

Text-only models still dominate the inventory (Figure~\ref{fig:modality}), processing tokenized text as both input and output. This category encompasses all major SLM families: Phi, Gemma, LLaMA, Mistral, Qwen, Yi, InternLM, SmolLM, MiniCPM, OLMo, Pythia, and DeepSeek.

\begin{figure}[t]
  \centering
  \includegraphics[width=\linewidth]{paper_slm_modality_distribution.pdf}
  \caption{Input modalities supported by documented SLM variants (a variant may support several). Vision has become common; audio and video remain a small but growing minority.}
  \label{fig:modality}
\end{figure}

\subsubsection{Multimodal (Vision-Language) SLMs}

A growing subset accepts image inputs alongside text, enabling visual question answering, image captioning, document understanding, and visual reasoning. Notable examples include:

\begin{itemize}
    \item \textbf{PaliGemma 3B} \cite{chen2023paligemma}: Google's 3B vision-language SLM combining a SigLIP-400M image encoder with a Gemma decoder; state-of-the-art on transfer benchmarks.

    \item \textbf{Phi-3-vision (4.2B)}: Microsoft's first vision-capable Phi model, adding a 400M CLIP encoder to Phi-3-mini's text backbone; competitive on DocVQA and ChartQA.

    \item \textbf{Gemma 3 (1B/4B/12B/27B)}: Google's Gemma 3 is the first SLM family where all sizes include image understanding by default, using a SigLIP image encoder \cite{team2025gemma3}.

    \item \textbf{Qwen2.5-VL 3B/7B}: Alibaba's multimodal SLMs with native support for video, charts, and documents.
\end{itemize}

\subsubsection{Audio and Multi-Modal SLMs}

A smaller but rapidly growing category incorporates audio alongside text and vision. Gemma~3n \cite{google2025gemma3n} is built mobile-first: its MatFormer (``Matryoshka'') architecture nests a 2B-effective sub-model inside a 4B-effective model so that one download yields a spectrum of sizes, and Per-Layer Embeddings cut the resident memory footprint of the 5B/8B raw models to roughly that of a 2B/4B model (2--3~GB). Gemma~4 \cite{gemma4_2026} extends the family from 2.3B-effective to 31B with dense and MoE variants, an integrated thinking mode, and, for the 12B model, an encoder-free architecture that ingests raw audio and image patches directly. Apple's on-device Foundation Model \cite{gunter2025afm} is a $\sim$3B multimodal model quantized with 2-bit quantization-aware training and KV-cache sharing. The full text+vision+audio frontier in a single sub-10B model is now a shipping capability rather than a research aspiration.

% ----------------------------------------------------------------
\subsection{Taxonomy by Application Domain}
\label{subsec:taxonomy_domain}

\subsubsection{General-Purpose SLMs}

General-purpose SLMs are trained on diverse data and evaluated on multi-domain benchmarks (MMLU, HellaSwag, ARC, GSM8K). The majority of SLMs in our taxonomy—particularly the flagship families from Microsoft (Phi), Google (Gemma), Meta (LLaMA), Mistral AI, Alibaba (Qwen), and DeepSeek—fall into this category.

\subsubsection{Domain-Specific SLMs}

Domain-specific SLMs sacrifice breadth for depth, outperforming general-purpose models on narrow tasks:

\begin{itemize}
    \item \textbf{Healthcare:} Meerkat (7B) \cite{kim2025meerkat}, fine-tuned on clinical QA datasets, surpasses GPT-4 on the USMLE Step 1 benchmark while running on a single A100 GPU. OptimCLM \cite{hasan2024optimclm} applies combined compression (distillation + quantization + pruning) to clinical language models without sacrificing diagnostic accuracy.

    \item \textbf{Code:} Code LLaMA \cite{roziere2023codellama}, StarCoder2 \cite{lozhkov2024starcoder2}, and DeepSeek-Coder \cite{zhu2024deepseek_coder} represent a mature class of code-focused SLMs with strong HumanEval and MBPP performance.

    \item \textbf{Legal:} Domain-specific legal SLMs trained on corpora of statutory text, case law, and contracts reduce administrative processing time by 30--40\% in pilot deployments while maintaining stricter data privacy than API-based LLMs.

    \item \textbf{Multilingual:} Aya 8B \cite{ustun2024aya} covers 101 languages with strong cross-lingual instruction following. Qwen2.5 7B \cite{qwen2024qwen25} serves as the de facto open-weight SLM for multilingual enterprise deployments.
\end{itemize}

Table~\ref{tab:taxonomy_summary} summarizes the four-dimensional taxonomy with representative models for each combination.

\begin{table}[ht]
\centering
\caption{Summary of the four-dimensional SLM taxonomy with representative models.}
\label{tab:taxonomy_summary}
\footnotesize\adjustbox{max width=\linewidth}{%
\begin{tabular}{llll}
\toprule
\textbf{Dimension} & \textbf{Category} & \textbf{Key Characteristics} & \textbf{Representative Models} \\
\midrule
\multirow{3}{*}{Architecture}
 & Dense Transformer & All params active per token; GQA, RoPE, SwiGLU & Phi-4-mini, Qwen3, SmolLM3, Gemma 3 \\
 & MoE              & Sparse activation; total $\gg$ active params       & Mixtral 8$\times$7B, Qwen3-30B-A3B, Llama 4 Scout \\
 & Hybrid SSM-Transformer & Linear-time; interleaved SSM/attention blocks   & Falcon-H1, Granite 4.0, Nemotron Nano 2 \\
\midrule
\multirow{3}{*}{Training Strategy}
 & From scratch     & Curated/synthetic data; full pre-training          & Phi, Qwen2.5/Qwen3, SmolLM3 \\
 & Distillation     & Knowledge transfer from larger teacher             & Gemma 3/4, DeepSeek-R1-Distill, Phi-4-mini-reasoning \\
 & Pruning-based    & Parameter removal from large model + distillation  & Minitron, Sheared-LLaMA \\
\midrule
\multirow{3}{*}{Modality}
 & Text-only        & Tokenized text I/O; majority of production SLMs   & Mistral 7B, SmolLM3, Qwen3 \\
 & Vision-Language  & Image + text; CLIP/SigLIP encoder added            & Phi-3-vision, Gemma 3, PaliGemma \\
 & Audio/Multimodal & Text + vision + audio in one sub-10B model         & Gemma 3n, Gemma 4, Apple AFM \\
\midrule
\multirow{4}{*}{Domain}
 & General-purpose  & Multi-domain training; evaluated on broad benchmarks & Phi, Gemma, LLaMA, Qwen, Mistral \\
 & Healthcare       & Clinical NLP; HIPAA-compliant on-premises deployment & Meerkat 7B, OptimCLM \\
 & Code             & Code corpora; FIM training; HumanEval/MBPP focus  & Code LLaMA, StarCoder2, DS-Coder \\
 & Legal/Multilingual & Domain corpora; regulatory and cross-lingual tasks & Aya 8B, Qwen2.5 multilingual \\
\bottomrule
\end{tabular}%
}
\end{table}

% ================================================================
% Section: Model Compression Techniques for SLMs
% Merged from main (2).tex — Section IV + sections/background.tex
% ================================================================

\section{Model Compression Techniques for SLMs}
\label{sec:compression}

Model compression is central to the SLM ecosystem: it is the primary pathway for reducing large pre-trained models to sizes deployable on edge hardware, and it provides a complementary route to training-from-scratch. The principal techniques are knowledge distillation, quantization, and pruning; in practice, compound pipelines combining all three have become the state of the art. Table~\ref{tab:compression_compare} summarizes the key characteristics of each approach.

\begin{table}[ht]
\small
\caption{Comparison of model compression techniques for SLMs.}
\label{tab:compression_compare}
\centering
\footnotesize\adjustbox{max width=\linewidth}{%
\begin{tabular}{@{}lp{2.0cm}p{2.0cm}p{2.2cm}p{2.0cm}@{}}
\toprule
\textbf{Method} & \textbf{Requires Retraining} & \textbf{Compression Ratio} & \textbf{Accuracy Retention} & \textbf{Hardware Req.} \\
\midrule
Quantization (8-bit)  & No  & $2\times$          & ${\sim}99\%$   & Low \\
Quantization (4-bit)  & No  & $4\times$          & ${\sim}96\%$   & Low \\
Quantization (2-bit)  & No  & $8\times$          & ${\sim}85\%$   & Low \\
Pruning (50\% sparse) & Partial & $2\times$      & ${\sim}93\%$   & Medium \\
Pruning (75\% sparse) & Partial & $4\times$      & ${\sim}82\%$   & Medium \\
Structured Pruning    & Yes & $2$--$5\times$     & ${\sim}90\%$   & Medium \\
Knowledge Distillation& Yes & $10$--$100\times$ & $85$--$95\%$   & High \\
KD + Pruning + Quant  & Yes & $40\times+$        & ${\sim}90\%$   & High \\
\bottomrule
\end{tabular}%
}
\end{table}

% ----------------------------------------------------------------
\subsection{Knowledge Distillation}
\label{subsec:compression_kd}

Knowledge distillation (KD) \cite{hinton2015distilling} transfers the ``dark knowledge'' embedded in a teacher model's soft output distribution to a smaller student model. The student is trained to minimize a combination of the standard cross-entropy loss on labeled data and a KL-divergence loss with respect to the teacher's output logits at a temperature $T > 1$. This approach was pioneered for BERT compression (DistilBERT \cite{sanh2019distilbert}, TinyBERT \cite{jiao2020tinybert}) and has since been substantially extended for generative SLMs.

\subsubsection{White-Box vs.\ Black-Box KD}

The survey by Yang et al.\ \cite{yang2024kd_survey} classifies LLM distillation into two paradigms:

\textbf{White-box KD} has access to the teacher's intermediate representations. TinyBERT \cite{jiao2020tinybert} pioneered layer-by-layer feature alignment, matching student hidden states, attention matrices, and embedding layers to teacher counterparts via MSE loss. In the decoder-only LLM setting, this requires pairing student and teacher layers—a non-trivial task when sizes differ substantially—and is computationally expensive but yields superior retention of teacher capabilities.

\textbf{Black-box KD} uses only the teacher's generated outputs—either token probabilities or free-form text—as the supervision signal. Orca \cite{ho2022orca} exemplifies text-output black-box KD: GPT-4 generates step-by-step reasoning traces, which a LLaMA-based student learns to replicate. DeepSeek-R1 \cite{deepseek2025r1} extends this to long chain-of-thought reasoning distillation, producing 7B and 14B students that dramatically outperform identically-sized from-scratch models on mathematical reasoning.

\subsubsection{Chain-of-Thought Distillation}

A particularly productive form of black-box KD is chain-of-thought (CoT) distillation \cite{wei2022cot}: the teacher model generates explicit reasoning steps, and the student learns to produce or internalize these steps. Feng et al.\ \cite{feng2024counterfactual} demonstrate counterfactual distillation—augmenting CoT traces with contrastive incorrect reasoning—as a way to sharpen the student's reasoning boundaries. Liao et al.\ \cite{liao2024nesycd} propose Neural-Symbolic Collaborative Distillation, using a neuro-symbolic teacher to provide structured reasoning supervision that is more transferable to small students than purely neural CoT.

\subsubsection{Data Augmentation in KD}

The survey by Xu et al.\ \cite{xu2024kd_survey} highlights data augmentation (DA) as a central paradigm within KD frameworks. When the student lacks sufficient coverage of the teacher's knowledge space, the teacher can generate synthetic training instances targeting specific skill gaps—the mechanism underlying the Phi family's success \cite{gunasekar2023phi1}. The efficient KD approach of Ballout et al.\ \cite{ballout2024efficient_kd} demonstrates that targeted data augmentation can reduce the student training compute by 60\% while preserving 95\% of the teacher's task-specific accuracy.

\subsubsection{Self-Distillation}

An emerging approach eliminates the external teacher entirely. Fu et al.\ \cite{fu2024dynamic_selfd} propose Dynamic Self-Distillation, in which the model uses its own predictions from previous mini-batches as soft targets during fine-tuning, providing a form of implicit regularization that reduces overfitting on small fine-tuning datasets—a critical challenge for SLMs.

% ----------------------------------------------------------------
\subsection{Quantization}
\label{subsec:compression_quantization}

Quantization reduces the numerical precision of model weights and/or activations from 32-bit floats (FP32) or 16-bit (FP16/BF16) to integers or low-bit floats (INT8, INT4, FP4, NF4, INT2). This directly reduces memory footprint and, on hardware with native integer arithmetic (NPUs, modern CPUs), reduces inference latency.

\subsubsection{Post-Training Quantization (PTQ)}

PTQ is performed after training is complete, requiring only a small calibration dataset (typically 128--512 samples) and no gradient computation. It is the dominant quantization paradigm because it is cheap and applies to any model without access to training infrastructure.

\textbf{GPTQ} \cite{frantar2023gptq} applies second-order weight compression: it uses an approximation of the inverse Hessian of the loss with respect to each layer's weights to determine optimal rounding decisions, minimizing the layer-wise squared error introduced by quantization. For SLMs, GPTQ 4-bit quantization of a 7B model reduces VRAM from ${\sim}14$~GB to ${\sim}4$~GB.

\textbf{AWQ (Activation-Aware Weight Quantization)} \cite{lin2024awq} observes that not all weights are equally important for activation output: 1\% of weights (those corresponding to large-magnitude activation channels) account for disproportionate reconstruction error. AWQ identifies these ``salient'' channels using a small calibration set and applies per-channel scaling to protect them before quantization, consistently outperforming GPTQ on complex reasoning tasks and winning the MLSys 2024 Best Paper award.

\textbf{GGUF} is a quantization format popularized by \texttt{llama.cpp} \cite{gerganov2023llamacpp} that packages model weights, vocabulary, and architecture metadata in a single binary file. GGUF supports a range of quantization levels (Q2\_K through Q8\_0), enabling users to trade off quality and memory footprint across a continuous spectrum. Its practical importance for SLM deployment cannot be overstated: GGUF enables LLaMA-3.2-3B to run at 60+ tokens/second on consumer laptops without a GPU.

\textbf{SqueezeLLM} \cite{kim2024squeezellm} introduced sensitivity-based non-uniform quantization with dense-and-sparse decomposition, demonstrating that sensitivity varies dramatically across weight dimensions and enabling superior accuracy at ultra-low bit-widths.

\subsubsection{Quantization-Aware Training (QAT)}

QAT integrates simulated quantization into the training loop, allowing the model to learn weight values that are robust to quantization error. This is particularly valuable for aggressive quantization (2-bit, 3-bit) where PTQ accuracy drops unacceptably. QAT produces models with 2--3\% higher task accuracy than PTQ at the same bit-width but requires access to training data and significant compute.

\subsubsection{Key Finding: Quantization vs.\ Pruning}

Kurz and Zhao \cite{kurz2025pruning_vs_quant} conducted the first systematic comparison of PTQ quantization (GPTQ, AWQ) against post-training pruning (SparseGPT, Wanda) specifically in the SLM regime (0.5B--3.8B parameters, six models, seven languages, seven downstream tasks). Their primary conclusion: \emph{quantization consistently outperforms pruning in preserving model fidelity, multilingual perplexity, and reasoning accuracy}. AWQ at 4-bit outperforms 50\% Wanda pruning on 6 of 7 tasks. Quantization's advantage is most pronounced on language modeling and multilingual tasks; the gap narrows on complex knowledge tasks, where both methods show similar degradation under high compression.

% ----------------------------------------------------------------
\subsection{Pruning}
\label{subsec:compression_pruning}

Pruning removes parameters from a pre-trained model, either individually (unstructured) or in structured groups (entire heads, layers, or channels). Structured pruning yields hardware-friendly sparse models compatible with standard matrix multiplication libraries; unstructured pruning achieves higher compression ratios but requires special sparse hardware or software to realize latency benefits.

\subsubsection{Unstructured Pruning}

\textbf{SparseGPT} \cite{frantar2023sparsegpt} extends the GPTQ framework to weight removal: it prunes weights layer-by-layer using inverse-Hessian information to minimize the reconstruction error introduced by zeroing each weight. SparseGPT can prune 175B GPT-3 to 50\% sparsity in a single forward pass without weight updates, with $<1$\% perplexity increase.

\textbf{Wanda (Pruning by Weights and Activations)} \cite{sun2024wanda} proposes a simpler criterion: prune weights with the smallest product of weight magnitude and corresponding input activation norm. This requires only a forward pass through the calibration data (no second-order information) and runs $10\times$ faster than SparseGPT while achieving competitive performance at 50\% sparsity.

\subsubsection{Structured Pruning}

Structured pruning removes entire attention heads, MLP channels, or transformer layers, yielding immediately deployable models without special sparsity libraries. Layer pruning—removing the least important layers as measured by cosine similarity between adjacent layer outputs—has been shown to reduce LLaMA-3 8B to a 6B model with $<2$\% MMLU degradation. The Minitron pipeline \cite{minitron2024} applies depth and width pruning simultaneously, guided by importance scores derived from activation magnitudes over a calibration set.

% ----------------------------------------------------------------
\subsection{Combined Compression Pipelines}
\label{subsec:compression_combined}

The state of the art in SLM compression combines pruning, distillation, and quantization in a multi-stage pipeline. NVIDIA Minitron \cite{minitron2024} achieves a $40\times$ reduction in training compute relative to training LLaMA-3.1-8B from scratch by: (1)~structured pruning to target size; (2)~knowledge distillation from the original 8B teacher to the pruned student; (3)~post-training quantization for deployment. The resulting 4B Minitron model matches the performance of 4B models trained from scratch on commonsense, code, and math benchmarks.

OptimCLM \cite{hasan2024optimclm} applies a similar pipeline to clinical language models, combining all three compression stages to produce models under 1B that achieve within 5\% of clinical NLP accuracy of their 7B counterparts. These results suggest that the combination of pruning (for architectural reduction) + distillation (for capability recovery) + quantization (for memory reduction) represents the practical frontier of SLM compression.

% ================================================================
% Section: Consolidated Reported Performance
% Replaces the former methods.tex + results.tex.
% Framing follows paper2 (meta-analysis): reported values with
% provenance, NOT original benchmark runs.
% ================================================================

\section{Consolidated Reported Performance}
\label{sec:reported}

This section assembles a comparative picture of SLM capability and deployment
behavior from the public record. In contrast to an original benchmarking study,
the performance values analyzed here are \emph{reported} figures, synthesized
from vendor technical reports, official model cards, community leaderboards, and
third-party evaluations. We first document the protocol used to collect and
reconcile these figures (Section~\ref{subsec:reported_protocol}), then present
the consolidated synthesis (Sections~\ref{subsec:reported_overall}--\ref{subsec:reported_safety}).
Every table entry traces to a dated, citable origin recorded in
Appendix~\ref{app:provenance}; where sources disagree, the reconciliation rule
of Section~\ref{subsec:reported_sources} applies.

\subsection{Data-Collection Protocol}
\label{subsec:reported_protocol}

\subsubsection{Model Selection}
\label{subsec:reported_model_selection}
Model inclusion required: (1)~public availability of weights or a publicly
documented evaluation card; (2)~parameter count below 10~billion, with Phi-4 at
14B retained as an upper-bound reference; (3)~at least one reported score on an
established benchmark used here; and (4)~representation of diverse training
methodologies (synthetic-data-centric, web-scale pretraining, distillation).
Table~\ref{tab:model_specs} lists the resulting 17-model cohort with
architectural specifications as reported by each model's release materials. The
cohort spans four orders of magnitude in parameter count and six developer
organizations, enabling analysis of how training methodology, architecture, and
scale interact to determine reported performance.

\begin{table}[ht]
\centering
\caption{Architectural specifications of reviewed models (as reported by each model's official release materials).}
\label{tab:model_specs}
\footnotesize\adjustbox{max width=\linewidth}{%
\begin{tabular}{lccccccc}
\toprule
\textbf{Model} & \textbf{Params} & \textbf{Layers} & \textbf{Hidden} & \textbf{Attention} & \textbf{Context} & \textbf{Vocab} & \textbf{Train Tokens} \\
\midrule
TinyLlama-1.1B & 1.1B & 22 & 2,048 & MHA-32h & 2K & 32K & 3.0T \\
SmolLM2-1.7B & 1.7B & 24 & 2,048 & MHA-32h & 2K & 49K & 11.0T \\
Qwen2.5-0.5B & 0.5B & 24 & 896 & GQA-14/2 & 32K & 152K & 18.0T \\
Qwen2.5-1.5B & 1.5B & 28 & 1,536 & GQA-12/2 & 32K & 152K & 18.0T \\
Qwen2.5-3B & 3.0B & 36 & 2,048 & GQA-16/2 & 32K & 152K & 18.0T \\
Qwen2.5-7B & 7.0B & 28 & 3,584 & GQA-28/4 & 128K & 152K & 18.0T \\
Llama 3.2-1B & 1.2B & 16 & 2,048 & GQA-32/8 & 128K & 128K & 9.0T \\
Llama 3.2-3B & 3.2B & 28 & 3,072 & GQA-24/8 & 128K & 128K & 9.0T \\
Gemma 2-2B & 2.6B & 26 & 2,304 & GQA-8/4 & 8K & 256K & 2.0T \\
Gemma 2-9B & 9.2B & 42 & 3,584 & GQA-16/8 & 8K & 256K & 8.0T \\
Gemma 3-1B & 1.0B & 26 & 1,536 & GQA-LG & 32K & 262K & 2.0T \\
Gemma 3-4B & 4.3B & 34 & 2,560 & GQA-LG & 128K & 262K & 4.0T \\
Mistral 7B & 7.3B & 32 & 4,096 & GQA-32/8 & 8K & 32K & --- \\
Phi-3-mini & 3.8B & 32 & 3,072 & MHA-32h & 128K & 32K & 3.3T \\
Phi-3.5-mini & 3.8B & 32 & 3,072 & MHA-32h & 128K & 32K & 3.4T \\
Phi-4-mini & 3.8B & 32 & 3,072 & GQA-24/8 & 128K & 200K & 5.0T \\
Phi-4 & 14.0B & 40 & 5,120 & GQA & 16K & 100K & $\sim$10T \\
\bottomrule
\end{tabular}%
}\\[5pt]
{\raggedright \small \textit{Note:} MHA = Multi-Head Attention; GQA = Grouped-Query Attention (Q/KV heads); GQA-LG = GQA with local--global interleaving. Training-token counts are reported or estimated values from each model's official technical report or model card.\par}
\end{table}

\subsubsection{Benchmark Tasks and Metrics}
\label{subsec:reported_benchmarks}
Fifteen benchmark tasks across seven NLP categories were included: MMLU
\citep{hendrycks2021mmlu}, HellaSwag \citep{zellers2019hellaswag}, Winogrande
\citep{sakaguchi2020winogrande}, ARC-Challenge \citep{clark2018arc}, PIQA
\citep{bisk2020piqa}, and BoolQ \citep{clark2019boolq} for commonsense and
language reasoning; GSM8K \citep{cobbe2021gsm8k} and MATH \citep{hendrycks2021math}
for mathematical reasoning; HumanEval \citep{chen2021humaneval} and MBPP
\citep{austin2021mbpp} for code generation; MNLI \citep{williams2018mnli} and ANLI
\citep{nie2020anli} for natural language inference; TriviaQA \citep{joshi2017triviaqa}
and SQuAD~2.0 \citep{rajpurkar2018squad2} for question answering; CNN/DailyMail
\citep{hermann2015cnndm} and XSum \citep{narayan2018xsum} for summarization; and
MT-Bench \citep{zheng2023mtbench} for instruction following. Inclusion of a task
required that at least three reviewed models have publicly reported scores under
compatible settings. Reported figures correspond to each benchmark's standard
metric (accuracy for multiple choice; EM/F1 for QA; pass@1 for code; ROUGE-L for
summarization; GPT-4 judge score for MT-Bench; exact match for GSM8K; symbolic
equivalence for MATH). Efficiency metrics---tokens per second, time-to-first-token,
peak memory, on-disk size, and reported energy---are each annotated in
Appendix~\ref{app:provenance} with the measurement conditions given by their
source, because these numbers are especially sensitive to hardware, batch size,
context length, and backend.

\subsubsection{Hardware and Quantization Scope}
\label{subsec:reported_hardware}
The review covers reported behavior across six configurations spanning the
cloud-to-edge spectrum (Table~\ref{tab:hardware_configs}) and three precisions:
FP16 (reference), INT8, and INT4, with GGUF~Q4\_K\_M, GPTQ group-128, and AWQ
group-128 treated as distinct INT4 subclasses when the source distinguishes them.
Inference-framework attribution (vLLM \citep{kwon2023vllm}, llama.cpp
\citep{gerganov2023llamacpp}, ONNX Runtime \citep{microsoft2019onnx}, ExecuTorch
\citep{meta2024executorch}, MLC-LLM \citep{mlc2024}) follows the stack used by the
source of each figure; cross-framework comparisons are made only when compared
figures share a framework or are explicitly flagged.

\begin{table}[ht]
\centering
\caption{Hardware configurations covered by the review. For each reviewed model, reported performance on a configuration was included where publicly available.}
\label{tab:hardware_configs}
\footnotesize\adjustbox{max width=\linewidth}{%
\begin{tabular}{lllll}
\toprule
\textbf{Configuration} & \textbf{Hardware} & \textbf{RAM/VRAM} & \textbf{TDP} & \textbf{Use Case} \\
\midrule
Cloud GPU & NVIDIA A100 80GB & 80 GB HBM2e & 300W & Training + serving \\
High-end Consumer & NVIDIA RTX 4090 & 24 GB GDDR6X & 450W & Local inference \\
Mid-range Consumer & NVIDIA RTX 3060 & 12 GB GDDR6 & 170W & Budget GPU \\
CPU-only (Laptop) & Intel i7-13700H & 32 GB DDR5 & 45W & Laptop deployment \\
Edge SBC & Raspberry Pi 5 & 8 GB LPDDR4X & 15W & Embedded IoT \\
Mobile SoC & Snapdragon 8 Gen 3 & 12 GB LPDDR5X & $\sim$10W & Smartphone \\
\bottomrule
\end{tabular}%
}
\end{table}

\subsubsection{Source Selection, Extraction, and Reconciliation}
\label{subsec:reported_sources}
Reported values were collected in decreasing order of priority from: (1)~official
vendor technical reports and release papers; (2)~official model cards on Hugging
Face and vendor repositories, including linked evaluation appendices;
(3)~community-operated standardized evaluations (the Open LLM Leaderboard, HELM
\citep{liang2023helm}, and lm-evaluation-harness \citep{gao2023harness}-based
community reports); (4)~third-party independent evaluations (Artificial Analysis,
LMSYS Chatbot Arena, peer-reviewed follow-up studies); and (5)~vendor developer
blogs, used only to fill gaps. All figures were extracted within a single
collection window to preserve time-consistency across models.

When qualifying sources disagreed on the same model--task--condition tuple, the
figure from the source higher in the priority list was retained and the
alternative recorded in Appendix~\ref{app:provenance}; where the discrepancy
exceeded 2.0~percentage points (or the metric equivalent), both figures were
retained as a reported range or the tuple was excluded. A tuple was excluded from
cross-model comparison when: (a)~only a single self-reported value existed and
could not be cross-checked; (b)~the harness, prompt format, or shot count
materially differed from the setting used by at least half the compared models
and could not be normalized; or (c)~the source did not specify base versus
instruction-tuned status.

\subsubsection{Limitations of the Evidence Base}
\label{subsec:reported_limits}
Because the figures are compiled from heterogeneous public sources rather than a
single controlled pipeline, four limitations bound the strength of every
comparison below. \textbf{Prompt-format and harness variation:} different sources
use different templates and answer-extraction heuristics, which can shift reported
accuracy by several points under the same nominal benchmark. \textbf{Self-reporting
bias:} vendor-reported scores are systematically more favorable than third-party
reproductions; where third-party figures exist, comparisons preferentially use
them and flag the divergence. \textbf{Condition-coverage sparsity:} not every
model has reported scores for every task, tier, or precision; sparse tuples are
marked, not imputed. \textbf{Snapshot staleness:} the ecosystem evolves monthly,
and figures reflect the state of public reporting at the collection window. These
limitations motivate the empirical-verification agenda of
Section~\ref{subsec:challenges_verification}.

\subsection{Overall Benchmark Performance}
\label{subsec:reported_overall}

\begin{table}[t]
\centering
\caption{Reported zero-shot FP16 benchmark scores (instruction-tuned variants). \textbf{Bold} = best among SLMs. All values are reported figures; per-figure sources in Appendix~\ref{app:provenance}. MATH and MT-Bench columns and the extended task set appear in Table~\ref{tab:appendix_zero_shot_extended}.}
\label{tab:zero_shot}
\footnotesize\adjustbox{max width=\linewidth}{%
\begin{tabular}{lcccccc}
\toprule
\textbf{Model} & \textbf{Params} & \textbf{ARC-C} & \textbf{HellaSwag} & \textbf{GSM8K} & \textbf{HumanEval} & \textbf{MMLU} \\
\midrule
Phi-4-mini-Inst. & 3.8B & 83.7 & 69.1 & \textbf{88.6} & 72.0 & 67.3 \\
Qwen2.5-7B-Inst. & 7.6B & \textbf{90.1} & 80.0 & 88.7 & \textbf{84.8} & \textbf{72.6} \\
Gemma 2-9B-Inst. & 9.2B & 89.8 & \textbf{87.1} & 84.9 & 68.9 & 71.3 \\
Gemma 3-4B-Inst. & 4.0B & 85.2 & 78.4 & 82.1 & 64.6 & 65.8 \\
LLaMA 3.2-3B-Inst. & 3.2B & 76.1 & 77.2 & 75.6 & 58.5 & 61.8 \\
LLaMA 3.2-1B-Inst. & 1.2B & 62.3 & 63.8 & 47.2 & 35.4 & 48.6 \\
Mistral 7B-Inst. v0.3 & 7.3B & 82.7 & 73.5 & 81.9 & 66.5 & 63.0 \\
Qwen2.5-3B-Inst. & 3.0B & 78.6 & 72.1 & 79.8 & 65.2 & 65.6 \\
Qwen2.5-1.5B-Inst. & 1.5B & 68.4 & 64.2 & 61.8 & 42.8 & 58.3 \\
Qwen2.5-0.5B-Inst. & 0.5B & 52.1 & 48.3 & 38.5 & 22.0 & 47.5 \\
SmolLM2-1.7B-Inst. & 1.7B & 59.8 & 60.2 & 42.6 & 28.7 & 50.1 \\
TinyLlama-1.1B & 1.1B & 47.5 & 53.8 & 18.4 & 12.2 & 36.2 \\
StableLM 2-1.6B & 1.6B & 55.2 & 57.6 & 32.1 & 22.5 & 44.3 \\
\midrule
\textit{GPT-4o (reference)} & \textit{---} & \textit{96.4} & \textit{95.3} & \textit{95.8} & \textit{90.2} & \textit{88.7} \\
\textit{LLaMA 3.1-70B-Inst.} & \textit{70B} & \textit{92.8} & \textit{88.5} & \textit{93.2} & \textit{80.5} & \textit{82.2} \\
\bottomrule
\end{tabular}%
}
\end{table}

Several patterns are consistent across the sources reviewed
(Table~\ref{tab:zero_shot}). No single SLM leads on every task. In the 3--4B band,
Phi-4-mini is reported strongest on mathematical reasoning (GSM8K) and MATH, while
Qwen2.5-3B is close behind with broader consistency. In the 7--10B band,
Qwen2.5-7B is reported ahead of Mistral~7B and Gemma~2-9B on knowledge and code
benchmarks. Phi-4 (14B) serves as an upper-bound reference. The most striking
regularity is the effect of training data: models of near-identical size but
different training corpora---Qwen2.5-1.5B versus SmolLM2-1.7B, or Qwen2.5-0.5B
versus older 500M-class models---differ by 15--25~points on knowledge benchmarks,
a gap far larger than any attributable to architecture at this scale. A Pareto
view of reported MMLU against parameter count places Phi-4-mini and Qwen2.5-3B on
the sub-4B efficient frontier and Qwen2.5-7B and Gemma~2-9B on the 7--10B
frontier.

The extended task set (Appendix~\ref{subsec:appendix_extended_results},
Table~\ref{tab:appendix_zero_shot_extended}) adds ARC-E, TriviaQA, SQuAD~2.0,
MBPP, SST-2, and MNLI and does not change these orderings.

\subsection{Fine-Tuning versus Prompting}
\label{subsec:reported_finetuning}

\begin{table}[ht]
\centering
\caption{Reported zero-shot, 5-shot, and QLoRA fine-tuned accuracy for selected model--task pairs. Values are reported figures.}
\label{tab:finetuning}
\footnotesize\adjustbox{max width=\linewidth}{%
\begin{tabular}{llccccc}
\toprule
\textbf{Model} & \textbf{Setting} & \textbf{SST-2} & \textbf{MNLI} & \textbf{ARC-C} & \textbf{GSM8K} & \textbf{SQuAD 2.0 (F1)} \\
\midrule
Phi-4-mini & Zero-shot & 92.4 & 72.1 & 83.7 & 88.6 & 74.2 \\
 & 5-shot & 94.1 & 75.8 & 85.2 & 89.3 & 77.8 \\
 & QLoRA & \textbf{96.8} & \textbf{82.4} & 86.1 & 89.8 & \textbf{83.6} \\
\midrule
Qwen2.5-7B & Zero-shot & 93.8 & 74.6 & 90.1 & 88.7 & 76.8 \\
 & 5-shot & 95.2 & 77.3 & 91.2 & 89.5 & 79.4 \\
 & QLoRA & \textbf{97.3} & \textbf{84.1} & 91.5 & 90.1 & \textbf{85.2} \\
\midrule
LLaMA 3.2-3B & Zero-shot & 88.2 & 63.5 & 76.1 & 75.6 & 65.4 \\
 & 5-shot & 91.5 & 68.2 & 78.8 & 77.4 & 70.2 \\
 & QLoRA & \textbf{95.4} & \textbf{78.6} & 80.3 & 79.8 & \textbf{78.5} \\
\midrule
SmolLM2-1.7B & Zero-shot & 82.6 & 54.2 & 59.8 & 42.6 & 52.3 \\
 & 5-shot & 86.4 & 59.8 & 63.1 & 44.8 & 57.6 \\
 & QLoRA & \textbf{93.2} & \textbf{72.5} & 65.8 & 48.2 & \textbf{68.4} \\
\bottomrule
\end{tabular}%
}
\end{table}

Across sources that report all three paradigms for a common model--task pair
(Table~\ref{tab:finetuning}), QLoRA adaptation yields the largest gains for the
smallest models and for tasks where the base model is weak: reported sub-3B gains
average roughly 8--10~percentage points over zero-shot, against roughly
4--5~points at 7B. Few-shot prompting is reported to add a smaller, consistent
1.5--4~points but consumes context and raises latency, making it less attractive
for on-device agents. The practical reading, and the one most relevant to
Section~\ref{sec:agentic}, is a specialization threshold near 3B: below it,
task-specific fine-tuning is close to mandatory for delegated agent sub-tasks;
above it, few-shot prompting is often sufficient for well-represented tasks.

\subsection{Quantization Impact}
\label{subsec:reported_quantization}

\begin{table}[ht]
\centering
\caption{Reported effect of INT4 quantization on benchmark accuracy. Subscripts show the reported percentage-point change from the FP16 baseline.}
\label{tab:quantization}
\footnotesize\adjustbox{max width=\linewidth}{%
\begin{tabular}{llccccc}
\toprule
\textbf{Model} & \textbf{Precision} & \textbf{ARC-C} & \textbf{HellaSwag} & \textbf{GSM8K} & \textbf{MMLU} & \textbf{Avg.\ Retention} \\
\midrule
Qwen2.5-7B & FP16 & 90.1 & 80.0 & 88.7 & 72.6 & 100.0\% \\
 & AWQ-4bit & 88.9\textsubscript{$-1.2$} & 79.5\textsubscript{$-0.5$} & 87.2\textsubscript{$-1.5$} & 71.8\textsubscript{$-0.8$} & 98.8\% \\
 & Q4\_K\_M & 88.1\textsubscript{$-2.0$} & 78.8\textsubscript{$-1.2$} & 86.1\textsubscript{$-2.6$} & 70.9\textsubscript{$-1.7$} & 97.7\% \\
\midrule
Phi-4-mini & FP16 & 83.7 & 69.1 & 88.6 & 67.3 & 100.0\% \\
 & AWQ-4bit & 82.6\textsubscript{$-1.1$} & 68.6\textsubscript{$-0.5$} & 86.9\textsubscript{$-1.7$} & 66.5\textsubscript{$-0.8$} & 98.6\% \\
 & Q4\_K\_M & 81.8\textsubscript{$-1.9$} & 67.9\textsubscript{$-1.2$} & 85.8\textsubscript{$-2.8$} & 65.8\textsubscript{$-1.5$} & 97.6\% \\
\midrule
LLaMA 3.2-3B & FP16 & 76.1 & 77.2 & 75.6 & 61.8 & 100.0\% \\
 & Q4\_K\_M & 73.2\textsubscript{$-2.9$} & 75.4\textsubscript{$-1.8$} & 71.8\textsubscript{$-3.8$} & 59.6\textsubscript{$-2.2$} & 96.3\% \\
\midrule
SmolLM2-1.7B & FP16 & 59.8 & 60.2 & 42.6 & 50.1 & 100.0\% \\
 & Q4\_K\_M & 55.4\textsubscript{$-4.4$} & 56.8\textsubscript{$-3.4$} & 37.2\textsubscript{$-5.4$} & 46.5\textsubscript{$-3.6$} & 92.4\% \\
\bottomrule
\end{tabular}%
}
\end{table}

The sources converge on three points (Table~\ref{tab:quantization}). INT8 is
reported as effectively lossless (sub-point degradation on knowledge and math
benchmarks). INT4 costs roughly 1--4~points on average, and the loss is
\emph{not} uniform: smaller models degrade more, consistent with the argument
that each weight in a small model carries more information
\citep{slmquant2025}. Mathematical reasoning is reported to degrade faster
than knowledge recall under INT4, by a factor of roughly 1.7--2.2. Among INT4
methods, activation-aware quantization (AWQ) is reported to retain slightly more
accuracy than GPTQ for SLMs \citep{lin2024awq}, and GGUF's Q4\_K\_M offers the
best size--quality balance for CPU and edge backends. Reported model-size
reductions are close to the theoretical 4$\times$: a 7B model drops from roughly
15~GB to roughly 4.5~GB, and a 3.8B model to roughly 2.4~GB
(Table~\ref{tab:memory_footprint}).

\begin{table}[ht]
\centering
\caption{Reported model-weight memory footprint by precision (GB).}
\label{tab:memory_footprint}
\footnotesize
\begin{tabular}{lcccc}
\toprule
\textbf{Model} & \textbf{FP16} & \textbf{GPTQ-4bit} & \textbf{AWQ-4bit} & \textbf{Q4\_K\_M} \\
\midrule
Qwen2.5-7B & 15.2 & 4.2 & 4.4 & 4.5 \\
Phi-4-mini & 7.6 & 2.2 & 2.3 & 2.4 \\
LLaMA 3.2-3B & 6.4 & 1.9 & 2.0 & 2.1 \\
SmolLM2-1.7B & 3.4 & 1.1 & 1.1 & 1.2 \\
TinyLlama-1.1B & 2.2 & 0.7 & 0.8 & 0.8 \\
\bottomrule
\end{tabular}
\end{table}

\subsection{Edge Deployment}
\label{subsec:reported_edge}

\begin{table}[t]
\centering
\caption{Reported edge-deployment figures (INT4/Q4\_K\_M unless noted). Status labels indicate reported practical interactivity on the tested device.}
\label{tab:edge_deployment}
\footnotesize\adjustbox{max width=\linewidth}{%
\begin{tabular}{llccccl}
\toprule
\textbf{Model} & \textbf{Device} & \textbf{tok/s} & \textbf{TTFT (ms)} & \textbf{Peak RAM (GB)} & \textbf{Load (s)} & \textbf{Status} \\
\midrule
SmolLM2-1.7B & Raspberry Pi 5 & 8.2 & 890 & 1.4 & 4.2 & Good \\
TinyLlama-1.1B & Raspberry Pi 5 & 10.5 & 680 & 1.0 & 3.1 & Good \\
LLaMA 3.2-1B & Raspberry Pi 5 & 9.8 & 720 & 1.1 & 3.4 & Good \\
LLaMA 3.2-3B & Raspberry Pi 5 & 4.8 & 1,450 & 2.3 & 7.8 & Marginal \\
Phi-4-mini & Raspberry Pi 5 & 3.2 & 2,100 & 2.6 & 9.5 & Marginal \\
Qwen2.5-7B & Raspberry Pi 5 & 1.2 & 5,800 & 5.1 & 22.4 & No \\
\midrule
SmolLM2-1.7B & Jetson Orin Nano & 42.5 & 125 & 1.6 & 2.8 & Good \\
TinyLlama-1.1B & Jetson Orin Nano & 54.2 & 98 & 1.2 & 2.1 & Good \\
LLaMA 3.2-3B & Jetson Orin Nano & 28.4 & 210 & 2.5 & 4.5 & Good \\
Phi-4-mini & Jetson Orin Nano & 22.8 & 285 & 2.8 & 5.2 & Good \\
Qwen2.5-7B & Jetson Orin Nano & 11.5 & 620 & 5.3 & 12.8 & Marginal \\
Gemma 2-9B & Jetson Orin Nano & 6.2 & 1,150 & 6.4 & 18.5 & No \\
\midrule
LLaMA 3.2-1B$^*$ & Galaxy S24+ (NPU) & 45.2 & 85 & 1.2 & 3.8 & Good \\
LLaMA 3.2-3B$^*$ & Galaxy S24+ (NPU) & 22.5 & 165 & 2.4 & 6.2 & Good \\
Phi-4-mini & Galaxy S24+ (CPU) & 12.8 & 380 & 3.0 & 8.5 & Good \\
Qwen2.5-7B & Galaxy S24+ (CPU) & 5.4 & 920 & 5.5 & 18.2 & Marginal \\
\bottomrule
\end{tabular}%
}\\[4pt]
{\raggedright \small $^*$Reported with a 4-bit NPU backend (ExecuTorch/SpinQuant).\par}
\end{table}

The reported edge figures (Table~\ref{tab:edge_deployment}) draw clear
boundaries. An 8~GB Raspberry~Pi~5 is reported to run INT4 models up to about 3B
at marginally interactive speeds and sub-2B models comfortably. The Jetson Orin
Nano handles everything up to 9B at useful speeds. Snapdragon-class smartphones
with NPU delegation are reported to run models up to 7B, with interactive use
practical below 4B. Time-to-first-token, dominated by prefill, is the binding
constraint for conversational use: reported TTFT for a 3B INT4 model is around
4~s on the Pi but around 1~s on the Jetson. GQA and local--global attention
substantially reduce the reported KV-cache footprint, which otherwise dominates
memory at long context.

\subsection{Cost and Energy}
\label{subsec:reported_energy}

\begin{table}[ht]
\centering
\caption{Reported energy-per-inference estimates (joules per 1{,}000 completions, 512-token generation).}
\label{tab:energy}
\footnotesize\adjustbox{max width=\linewidth}{%
\begin{tabular}{lccccc}
\toprule
\textbf{Model} & \textbf{A100} & \textbf{RTX 4090} & \textbf{Laptop CPU} & \textbf{Jetson} & \textbf{RPi 5} \\
\midrule
Qwen2.5-0.5B & 180 & 320 & 410 & 280 & 612 \\
Qwen2.5-3B & 420 & 680 & 1,250 & 580 & 2,760 \\
Phi-4-mini & 510 & 780 & 1,480 & 650 & 3,180 \\
Qwen2.5-7B & 720 & 1,100 & 2,340 & 1,020 & --- \\
Phi-4 (14B) & 1,350 & 1,850 & --- & --- & --- \\
\bottomrule
\end{tabular}%
}
\end{table}

Reported energy figures (Table~\ref{tab:energy}) show that a datacenter GPU,
despite its high power draw, achieves the lowest energy \emph{per inference} for
any given model because of its throughput; an edge board consumes more energy per
inference but roughly 20$\times$ less absolute power, which favors it for
always-on, low-duty-cycle workloads. Normalizing accuracy by energy, a small
model on cheap hardware delivers markedly more benchmark points per joule than a
large model on a datacenter GPU, provided its accuracy clears the task threshold.
Table~\ref{tab:latency_compare} in Section~\ref{sec:deployment} translates these
into cost-per-token comparisons.

\subsection{Multilingual and Safety Reporting}
\label{subsec:reported_safety}

Reported multilingual scores (FLORES-200, XQuAD, MGSM) place Qwen2.5-7B ahead of
other SLMs, reflecting its large multilingual training corpus, with a wide gap to
English-centric models on non-English mathematical reasoning. Reported toxicity
evaluations (ToxiGen, RealToxicityPrompts) show substantially higher toxicity for
the smallest models (TinyLlama, sub-1B variants), consistent with the argument
that safety alignment is harder to preserve at reduced scale and is often absent
from research-grade small models. Both observations carry directly into
Section~\ref{sec:agentic}: an agent that routes non-English or safety-sensitive
calls to a small English-centric model inherits these gaps, which is an argument
for capability-aware rather than size-only routing.

% ================================================================
% Section: Deployment on Resource-Constrained Environments
% From main (2).tex — Section VI
% ================================================================

\section{Deployment on Resource-Constrained Environments}
\label{sec:deployment}

% ----------------------------------------------------------------
\subsection{On-Device and Mobile Deployment}
\label{subsec:deployment_mobile}

On-device deployment eliminates the latency, cost, and privacy exposure of cloud API inference. It enables AI functionality in offline and low-bandwidth environments and is a prerequisite for truly private personal assistants. The primary challenges are memory (RAM and storage), sustained inference throughput under thermal constraints, and battery impact.

\subsubsection{Apple Intelligence (AFM)}

Apple's Foundation Models \cite{apple2024intelligence,gunter2025afm} are the most fully documented on-device \gls{slm} deployment. The 2025 report \cite{gunter2025afm} describes a $\sim$3B on-device model, optimized for Apple silicon through KV-cache sharing (dividing the model so the second block reuses the first block's cache, a reported 37.5\% memory saving) and 2-bit quantization-aware training, paired with a Parallel-Track Mixture-of-Experts server model on Private Cloud Compute. The on-device model is reported competitive with Qwen3-4B and Gemma-3-4B on its target tasks; the server model with Llama~4 Scout. The third-generation framework \cite{apple2026afm3} adds a natively multimodal, sparsely activated on-device model and opens the same API surface to third-party cloud models, so that an application calls the local \gls{slm} by default and escalates to a frontier \gls{llm} only on demand. Google's Gemini Nano and Samsung's Gauss \cite{samsung2023gauss} occupy the same design point on Android and Galaxy devices.

\subsubsection{Samsung Gauss}

Samsung Gauss \cite{samsung2023gauss} is a family of language models for Galaxy devices, including a ${\sim}1.7$B on-device model targeting the Gauss1 chip and Exynos 2400 \gls{npu}. Samsung reports that Gauss enables on-device email summarization, translation, and writing assistance without requiring internet connectivity.

\subsubsection{General Mobile Deployment Considerations}

\glspl{slm} targeting mobile deployment must navigate three constraints. (1) \emph{Memory}: Consumer smartphones typically provide 6--12~GB RAM, of which the OS uses 2--4~GB, leaving 4--8~GB for applications. A 3B model at Q4 requires ${\sim}2$~GB, comfortable for most flagship devices. (2)~\emph{Thermal}: Sustained inference on mobile \glspl{npu} is throttled by thermals; realistic sustained throughput is lower than peak throughput benchmarks suggest. (3)~\emph{Battery}: Gemma~3 on Android reportedly consumes ${\sim}0.75\%$ battery per prompt on Pixel~9 Pro, enabling hundreds of interactions per charge. The fundamental constraint in edge deployment is memory bandwidth rather than compute: mobile devices offer 50--90~GB/s of memory bandwidth versus 2--3~TB/s for datacenter GPUs \cite{xu2024device}, making INT4 quantization not merely a storage optimization but a critical performance enabler.

% ----------------------------------------------------------------
\subsection{Edge and IoT Deployment}
\label{subsec:deployment_edge}

Below smartphone-class hardware lies a diverse ecosystem of edge devices: Raspberry Pi single-board computers, NVIDIA Jetson modules, and microcontrollers. \glspl{slm} targeting this tier typically require:

\begin{itemize}
    \item \textbf{Sub-1B or extreme quantization:} A Raspberry Pi~5 (8~GB RAM, Cortex-A76 CPU) can run Q4 quantized 1B models at 10--20 tokens/second. SmolLM2-135M at Q8 runs at $>$30~tok/s on Raspberry Pi~5. Benchmarking shows that llamafile achieves up to $4\times$ higher throughput and 30--40\% lower power consumption than Ollama for equivalent models on single-board computers \cite{costa2025sbc}.

    \item \textbf{Jetson Orin:} NVIDIA's Jetson Orin NX (16~GB) with TensorRT-LLM can serve 7B quantized models at ${\sim}30$~tok/s, enabling real-time local inference for industrial and robotics applications.

    \item \textbf{Microcontrollers (MCU):} Sub-100M models using aggressive quantization (INT4/INT2) can run on ARM Cortex-M series MCUs with 512~KB--2~MB SRAM for specialized tasks like keyword detection and sensor data analysis.
\end{itemize}

% ----------------------------------------------------------------
\subsection{Deployment Frameworks}
\label{subsec:deployment_frameworks}

A mature ecosystem of open-source frameworks enables \gls{slm} deployment across hardware targets. Table~\ref{tab:deployment_frameworks} provides a structured comparison.

\begin{table}[ht]
\caption{Key SLM deployment frameworks, their primary targets, and characteristics.}
\label{tab:deployment_frameworks}
\centering
\footnotesize\adjustbox{max width=\linewidth}{%
\begin{tabular}{@{}lllp{5.5cm}@{}}
\toprule
\textbf{Framework} & \textbf{Primary Target} & \textbf{Quantization} & \textbf{Notes} \\
\midrule
llama.cpp \cite{gerganov2023llamacpp}   & CPU, Metal, CUDA, ROCm & GGUF (Q2--Q8)   & De facto standard; minimal dependencies; ecosystem includes Ollama, llamafile \\
ONNX Runtime \cite{microsoft2019onnx} & GPU, CPU, NPU, WASM & INT4/INT8       & Cross-platform; DirectML/CoreML/CUDA backends; Phi-3-mini at 35~tok/s on Snapdragon NPU \\
MLX \cite{mlx2023}             & Apple Silicon          & 4-bit           & Unified memory; eliminates GPU transfers; 55+ tok/s for LLaMA-3.2-3B on M3~Pro \\
TensorRT-LLM \cite{nvidia_tensorrt2024} & NVIDIA GPUs   & INT4/INT8/FP16  & Continuous batching; kernel fusion; 2--3$\times$ over PyTorch for 7B models \\
ExecuTorch \cite{meta2024executorch} & Mobile (iOS, Android) & INT4          & 50~KB base runtime; 12+ hardware backends; deployed in Instagram, WhatsApp \\
MLC-LLM \cite{mlc2024}        & iOS, Android, web, MCU & 4-bit          & TVM compilation; WebGPU support; universal platform coverage \\
WebLLM \cite{webllm2024}       & Web browsers           & 4-bit WASM     & WebGPU acceleration; 80\% of native performance; OpenAI-compatible API \\
\bottomrule
\end{tabular}%
}
\end{table}

% ----------------------------------------------------------------
\subsection{Latency Comparison: SLMs vs.\ LLMs}
\label{subsec:deployment_latency}

Table~\ref{tab:latency_compare} illustrates the practical latency implications of \gls{slm} vs.\ \gls{llm} deployment for three representative real-world scenarios.

\begin{table}[ht]
\small
\caption{Latency and cost comparison: SLMs vs.\ LLMs for representative use cases.}
\label{tab:latency_compare}
\centering
\footnotesize\adjustbox{max width=\linewidth}{%
\begin{tabular}{@{}lp{2.5cm}p{2.5cm}p{2.5cm}@{}}
\toprule
\textbf{Metric} & \textbf{Real-time Chatbot} & \textbf{Code Completion} & \textbf{Edge IoT} \\
\midrule
\multicolumn{4}{@{}l}{\textit{SLM (3B, Q4, consumer GPU)}} \\
TTFT               & 80--150 ms    & 30--80 ms     & 50--120 ms  \\
Throughput         & 200--400 tok/s & 300--500 tok/s & 100--300 tok/s \\
Memory             & ${\sim}2$~GB  & ${\sim}2$~GB  & ${\sim}1$~GB \\
Cost/1M tokens     & \$0.05--0.2   & \$0.05--0.2   & \$0 (on-device) \\
\midrule
\multicolumn{4}{@{}l}{\textit{LLM (70B, FP16, multi-GPU)}} \\
TTFT               & 500--2000 ms  & 200--800 ms   & Not feasible \\
Throughput         & 30--80 tok/s  & 60--150 tok/s & N/A \\
Memory             & ${\sim}140$~GB & ${\sim}140$~GB & N/A \\
Cost/1M tokens     & \$3--10       & \$3--10        & N/A \\
\midrule
\multicolumn{4}{@{}l}{\textit{Frontier API (GPT-4o, cloud)}} \\
TTFT               & 300--1500 ms  & 200--600 ms   & Requires internet \\
Throughput         & 80--200 tok/s & 80--200 tok/s & N/A \\
Cost/1M tokens     & \$5--15       & \$5--15        & N/A \\
\bottomrule
\end{tabular}%
}
\end{table}

The data in Table~\ref{tab:latency_compare} illustrate a 5--20$\times$ cost advantage and comparable or superior latency for \glspl{slm} in chatbot and code-completion scenarios, with an absolute deployment advantage in offline edge contexts where \glspl{llm} are simply infeasible.

% ================================================================
% Section: SLMs as the Substrate of Agentic AI
% New section (position core of the consolidated paper)
% ================================================================

\section{Small Language Models as the Substrate of Agentic AI}
\label{sec:agentic}

The preceding sections establish three facts about the sub-10B regime: modern SLMs reach capability levels that were state-of-the-art for far larger models only two years earlier (Section~\ref{sec:reported}); they can be compressed to run on commodity and edge hardware with modest accuracy loss (Sections~\ref{sec:compression}--\ref{sec:deployment}); and a mature ecosystem now exists to deploy them (Section~\ref{subsec:deployment_frameworks}). This section argues that these facts, taken together with how agentic systems are actually built, make SLMs the \emph{default} execution substrate for AI agents, with large models reserved for a shrinking set of genuinely hard calls. The argument follows and extends the position of \citet{belcak2025slm_agentic}; we situate it within the broader literature on LLM--SLM collaboration \citep{chen2024role_small,chen2025collab_survey} and connect it to the deployment evidence consolidated in this paper.

\subsection{What Agents Delegate to Language Models}
\label{subsec:agentic_delegation}

An AI agent is a software system in which one or more language models operate under a control loop that also includes tools, memory, and a scaffold of prompts and parsers \citep{belcak2025slm_agentic}. The language model in such a system is rarely asked to hold an open-ended conversation. It is asked to perform bounded sub-tasks: decide which tool to call and with what arguments, extract fields from a document, classify an intent, rewrite a query, summarize a retrieved passage, produce a short code patch, or check whether a step succeeded. These calls share three properties. They are \emph{narrow}---each exercises a small slice of the model's competence. They are \emph{repetitive}---the same handful of call types recur thousands of times in production. And they are \emph{format-constrained}---the output must parse as JSON, match a function signature, or fall within a fixed label set, or the surrounding code breaks \citep{belcak2025slm_agentic}.

None of these properties rewards the breadth of a frontier LLM. A model that can also write a sonnet, debate philosophy, and reason across a 200-page contract is carrying capacity that the agent never invokes on the call in question. \citet{belcak2025slm_agentic} formalize this as the observation that agents use only a narrow subset of language-model functionality, and that agentic systems demand strict behavioral alignment to code interfaces---exactly the setting in which a specialized small model can match or exceed a general large one.

\subsection{The Case for SLMs in Agents}
\label{subsec:agentic_case}

\subsubsection{Capability Sufficiency}
\label{subsubsec:agentic_capability}

The first premise is that SLMs are already good enough for delegated sub-tasks. On instruction following, tool selection, and structured extraction---the operations that dominate agent traces---sub-10B models trained or tuned for the purpose are competitive with frontier models. Salesforce's xLAM family \citep{zhang2024xlam} illustrates the point: its 7B function-calling model ranked at the top of the Berkeley Function-Calling Leaderboard, and its 1B variant was reported to surpass GPT-3.5-Turbo and Claude-3~Haiku on the same task despite being one to two orders of magnitude smaller. Unified-reasoning SLMs extend this to multi-step tasks: Qwen3 \citep{qwen3_2025} exposes a controllable ``thinking budget'' so that an agent can spend inference compute only on the calls that need it, and distilled reasoning models such as Phi-4-mini-reasoning \citep{xu2025phi4mini_reasoning} and the DeepSeek-R1 distillations \citep{deepseek2025r1} bring competition-mathematics performance into the 3--8B range. \citet{belcak2025slm_agentic} argue that objections based on scaling laws (AV1: larger is permanently better) understate SLM capability because scaling-law studies hold architecture fixed, exclude task-specific fine-tuning, and predate inference-time reasoning---all three of which disproportionately help small models.

\subsubsection{Economics}
\label{subsubsec:agentic_economics}

The second premise is cost. \citet{belcak2025slm_agentic} estimate that serving a 7B SLM is 10--30$\times$ cheaper than serving a 70--175B LLM once latency, energy, and FLOPs are accounted for, and that parameter-efficient fine-tuning lets a small model be specialized for a new agent role in GPU-hours rather than GPU-weeks. The deployment figures consolidated in this paper are consistent with that range: Table~\ref{tab:latency_compare} shows a 5--20$\times$ cost advantage and comparable or better latency for a quantized 3B model relative to a 70B model on chatbot and code-completion workloads, and an absolute advantage in offline settings where the large model cannot run at all. Because an agent amortizes model cost over many calls per user task, a constant-factor saving on each call compounds into the difference between a viable and an unviable product margin.

\subsubsection{Behavioral Alignment and Operational Fit}
\label{subsubsec:agentic_alignment}

The third premise is control. An agent needs its model to emit exactly the format the scaffold expects, and to do so reliably across a long tail of inputs. Constraining a small specialized model to a narrow output distribution is easier than suppressing the much larger behavioral range of a frontier model, and fine-tuning a small model on the agent's own logged traces closes the residual gap \citep{belcak2025slm_agentic}. \citet{belcak2025slm_agentic} note a virtuous cycle here (their argument A7): agent interactions generate high-quality, in-distribution training data essentially for free, which is precisely what specialization of a small model requires. Operationally, small models are also easier to version, roll back, run on-premises for data-residency reasons, and co-locate with the tools they call.

\subsection{Heterogeneous SLM--LLM Systems}
\label{subsec:agentic_hetero}

The position is not that LLMs disappear from agents but that the default inverts: SLMs handle the repetitive language errands, and a strong generalist model is invoked only for planning, open-ended decomposition, or calls that genuinely need long context or broad world knowledge \citep{belcak2025slm_agentic}. Figure~\ref{fig:hetero_agent} sketches the arrangement. Realizing it requires a mechanism to decide, per call, which model to use. \citet{chen2025collab_survey} organize the collaboration literature into five interaction patterns---pipeline, routing, auxiliary, distillation, and fusion---of which routing is the most direct fit for agents.

\begin{figure}[t]
\centering
\footnotesize
\begin{tikzpicture}[
  >={Stealth[length=2mm]},
  box/.style={draw, rounded corners=2pt, align=center, inner sep=4pt, line width=0.5pt},
  slm/.style={box, fill=black!4, text width=40mm, align=left},
  llm/.style={box, fill=black!10, text width=40mm, align=left},
  router/.style={box, fill=black!6, text width=24mm, minimum height=10mm},
  lbl/.style={font=\scriptsize\itshape, text=black!55},
]
\node[box, text width=13mm] (task) {User task};
\node[box, right=7mm of task, minimum height=12mm, text width=20mm] (loop) {Agent\\control loop};
\node[box, below=6mm of loop, font=\scriptsize, text width=20mm] (tools) {Tools $\cdot$ Memory};
\node[router, right=9mm of loop] (router) {\textbf{Router}\\[1pt]\scriptsize difficulty $\cdot$ cost $\cdot$ privacy $\cdot$ network};
\node[slm, above right=2mm and 12mm of router] (slm) {\textbf{Specialised SLMs} — default\\[1pt]\scriptsize tool selection, argument filling, field extraction, classification, query rewrite, short code patch, step verification};
\node[llm, below right=2mm and 12mm of router] (llm) {\textbf{Generalist LLM} — escalation\\[1pt]\scriptsize planning, open-ended decomposition, long context, broad world knowledge};
\draw[->] (task) -- (loop);
\draw[->] (tools) -- (loop);
\draw[<->] (loop) -- (router);
\draw[->] (router.east) -- ($(slm.west)+(0,-1mm)$) node[lbl, midway, above] {most calls};
\draw[->] (router.east) -- ($(llm.west)+(0,1mm)$) node[lbl, midway, below] {minority};
\end{tikzpicture}
\caption{Heterogeneous SLM--LLM agent. A per-call router sends the majority of
delegated sub-tasks to specialised small models and escalates only planning-like
or long-context calls to a generalist LLM; tools and memory are shared. After
\citet{belcak2025slm_agentic}.}
\label{fig:hetero_agent}
\end{figure}

\subsubsection{Learned Routing}
\label{subsubsec:agentic_routing}

RouteLLM \citep{ong2024routellm} trains a lightweight router on human-preference data to predict, for each query, whether a weak model's answer would be judged as good as a strong model's, and dispatches accordingly; the authors report cost reductions of more than 2$\times$ on standard benchmarks with no measurable quality loss, and transfer of the trained router to unseen strong/weak model pairs. Division-of-Thoughts \citep{division_of_thoughts_2025} pushes routing below the query level, decomposing a task into sub-questions and assigning each to a local SLM or a cloud LLM according to estimated difficulty, for on-device agents. The open problems are difficulty estimation without an oracle, routing objectives that jointly weigh accuracy, latency, privacy sensitivity, and network availability, and handoff that does not produce a visible seam in output quality \citep{demystifying2025slm}.

\subsubsection{Production Instances}
\label{subsubsec:agentic_production}

Two large deployments already embody the heterogeneous pattern. Apple's Foundation Models framework pairs an on-device model of roughly 3B parameters, quantized with 2-bit quantization-aware training and KV-cache sharing, with a Private Cloud Compute server model built on a Parallel-Track Mixture-of-Experts transformer, and routes to the server only when the on-device model is expected to fall short \citep{gunter2025afm}. The on-device model is reported to be competitive with Qwen3-4B and Gemma-3-4B on its target tasks. Third-party providers can now register cloud models behind the same API surface as the local model, making the SLM the default and the LLM an opt-in escalation \citep{apple2026afm3}. Samsung's Gauss and Google's Gemini Nano occupy the same design point on their respective platforms \citep{samsung2023gauss}.

\subsection{Converting an LLM Agent to an SLM Agent}
\label{subsec:agentic_conversion}

For teams with an existing LLM-only agent, \citet{belcak2025slm_agentic} propose a six-step conversion pipeline that we summarize here because it operationalizes the position. \textbf{S1, instrumentation:} log every non-interface model call (inputs, outputs, latency) from the running agent, with encryption and access control. \textbf{S2, curation:} strip personally identifiable and protected health information and application secrets from the logs. \textbf{S3, clustering:} apply unsupervised clustering to the curated calls to surface the recurring request types---these are the specialization targets. \textbf{S4, selection:} for each cluster, choose a candidate SLM by matching required capability against benchmarks, license, and deployment constraints. \textbf{S5, specialization:} fine-tune the chosen SLM on the cluster's curated data, typically with LoRA or QLoRA. \textbf{S6, iteration:} re-run the loop periodically as the traffic distribution drifts. The pipeline is attractive precisely because the data it needs is a by-product of running the agent, and because each cluster converted removes a recurring frontier-model bill.

\subsection{Barriers to Adoption}
\label{subsec:agentic_barriers}

\citet{belcak2025slm_agentic} identify three reasons the shift has lagged the technical case. First, large centralized-inference investments create sunk-cost inertia and an operational bias toward one big model behind one API (B1). Second, SLMs are still evaluated on generalist leaderboards (MMLU, broad chat benchmarks) that are misaligned with agentic utility, so their fitness for delegated sub-tasks is under-measured (B2)---an evaluation gap this paper's consolidated tables also reflect, since public reporting concentrates on the same generalist benchmarks. Third, SLM releases receive far less attention than frontier launches, so awareness lags capability (B3). None of these is a technical obstacle, which is the point: the barriers are those of infrastructure, measurement, and attention, and each is already eroding as routing frameworks mature, agent-specific benchmarks such as the Berkeley Function-Calling Leaderboard gain traction, and on-device platforms ship SLMs to billions of devices by default.

% ================================================================
% Section: Domain-Specific Applications
% From main (2).tex — Section VII
% ================================================================

\section{Domain-Specific Applications}
\label{sec:applications}

The domains below are where \gls{slm} deployment is most advanced, and each is an
agentic setting in the sense of Section~\ref{sec:agentic}: a workflow of bounded,
repeated language operations---extract a diagnosis code, classify a clause,
complete a function, translate a passage---executed under regulatory or cost
constraints that make a cloud frontier model impractical. Domain fine-tuning of a
small model is the LLM-to-SLM conversion pipeline
(Section~\ref{subsec:agentic_conversion}) applied to a vertical.

\subsection{Healthcare and Clinical NLP}
\label{subsec:applications_healthcare}

Healthcare \gls{nlp} presents unique requirements: high factual accuracy, interpretability, HIPAA/GDPR compliance, and compatibility with clinical workflows. Large general-purpose \glspl{llm} are increasingly restricted from processing identifiable patient data due to data residency regulations, creating a strong pull toward on-premises \glspl{slm}.

The Meerkat family \cite{kim2025meerkat} demonstrates the potential of domain-specific \gls{slm} fine-tuning. Built on a 7B base with instruction fine-tuning on curated clinical QA datasets and chain-of-thought reasoning traces from medical board examinations, Meerkat achieves 85.7\% on USMLE Step~1---surpassing GPT-4 (86.0\% in zero-shot) while being deployable on a single A100 GPU rather than a multi-GPU cloud cluster. Meerkat's clinical reasoning is particularly strong on diagnostic pattern matching, outperforming GPT-4 by 4\% on complex differential diagnosis tasks \cite{kim2025meerkat}.

OptimCLM \cite{hasan2024optimclm} applies the full compression pipeline---knowledge distillation, structured pruning, and 4-bit quantization---to clinical language models, reducing a 7B model to under 1B while preserving 95\% of its accuracy on ICD coding and discharge summary generation tasks. This makes clinical \gls{nlp} feasible on hospital workstations without GPU infrastructure.

The broader \gls{slm} healthcare ecosystem includes BioMedLM (PubMed-trained 2.7B model), clinical note generation models, and radiology-specific models. A key trend is the use of \gls{rag} \cite{lewis2020rag} to anchor \gls{slm} outputs to verified clinical knowledge bases, reducing hallucination rates in patient-facing applications.

\subsection{Legal NLP}
\label{subsec:applications_legal}

Legal document processing---contract review, compliance monitoring, regulatory change analysis, case law search---involves long, precisely structured documents where hallucination is particularly dangerous. \glspl{slm} address legal \gls{nlp} deployment constraints in two ways: (1)~by enabling on-premises deployment to satisfy attorney-client privilege and data security requirements; and (2)~by enabling cost-effective processing of the massive volumes of routine legal work that would be prohibitively expensive with frontier \glspl{llm}.

Pilot deployments of domain-adapted \glspl{slm} (3B--7B range, fine-tuned on contract corpora) have demonstrated 30--40\% reductions in administrative processing time for standard contract review workflows. Models fine-tuned on specific regulatory domains (GDPR, SOC~2, FINRA) consistently outperform general-purpose GPT-4 on domain-specific compliance classification tasks, owing to the depth of in-domain knowledge concentration.

\subsection{Code Generation}
\label{subsec:applications_code}

Code-focused \glspl{slm} have matured into a category distinct from general \glspl{slm}, with specialized pre-training corpora (GitHub, Stack Overflow, documentation), code-aware tokenization, and evaluation on function-level benchmarks (HumanEval, MBPP, SWE-Bench).

\begin{description}
    \item[Code LLaMA \cite{roziere2023codellama}] Meta's family of code models fine-tuned on LLaMA-2 (7B, 13B, 34B). Code LLaMA 7B achieves 33.5\% HumanEval pass@1.

    \item[StarCoder2 \cite{lozhkov2024starcoder2}] BigCode's 3B/7B/15B models trained on The Stack v2 (600+ programming languages). StarCoder2-7B achieves 46.3\% HumanEval pass@1.

    \item[DeepSeek-Coder \cite{zhu2024deepseek_coder}] A 1B--33B family trained from scratch on 2T code tokens. The 6.7B variant achieves 47.6\% HumanEval pass@1, outperforming GPT-3.5 on code completion.

    \item[Codestral \cite{mistral2024codestral}] Mistral's 22B code model trained on 80+ programming languages with fill-in-the-middle (FIM) training for IDE integration.
\end{description}

\subsection{Education and Multilingual NLP}
\label{subsec:applications_multilingual}

Multilingual \glspl{slm} address the linguistic coverage gap in AI assistants: the majority of \gls{llm} pre-training data is English, leaving speakers of lower-resource languages with significantly degraded performance. Aya~8B \cite{ustun2024aya} is fine-tuned on 204 million multilingual instructions spanning 101 languages, with strong performance across African, South Asian, and Southeast Asian languages that are severely underrepresented in standard benchmarks.

Qwen2.5~7B \cite{qwen2024qwen25} is pre-trained on 18T tokens with strong multilingual coverage including Chinese, Japanese, Korean, Arabic, and multiple European languages, serving as the de facto open-weight \gls{slm} for multilingual enterprise deployments. Gemma~3's multilingual extension \cite{team2025gemma3} adds instruction-tuning data for 40+ languages at the 4B and 12B scales.

Educational applications include on-device tutoring systems (leveraging Phi-3 or SmolLM on tablets), interactive language learning assistants with real-time feedback, and low-connectivity educational tools for rural and developing-world deployments. The key requirement for educational \glspl{slm} is not just multilingual capability but also low hallucination rates and age-appropriate safety filtering.

\section{Discussion}
\label{section:discussion}

\subsection{What the Evidence Supports}
\label{subsec:discussion_key_findings}

Read together, the taxonomy (Section~\ref{sec:taxonomy}), the compression review
(Section~\ref{sec:compression}), and the audit of the evaluation record
(Section~\ref{sec:reported}) support four claims.

\textbf{Capability per parameter has risen sharply.} The strongest 3--4B models
of 2026 report scores on mathematics, code, and knowledge benchmarks that were
state of the art for 13B-class systems in 2023. Within families that publish
controlled ablations, the dominant explanatory variable is training data rather
than architecture: corpus composition and curation move results more than any
architectural choice available at this scale. The practical implication for model
builders is that data pipelines dominate architectural novelty as a lever on
\gls{slm} quality.

\textbf{Compression is a first-class capability, not an afterthought.} INT8 is
effectively lossless and INT4 costs a small, non-uniform amount of accuracy, with
smaller models and mathematical reasoning at the higher end. Because
autoregressive inference is memory-bandwidth-bound, 4-bit quantization is
simultaneously a storage optimization and a near-linear speedup on edge hardware.
Quantization preserves fidelity better than pruning below 4B parameters
\citep{kurz2025pruning_vs_quant}, which simplifies the practitioner's decision:
quantize aggressively, and prune only when a structured size reduction is
independently required.

\textbf{Right-sizing beats scaling for bounded tasks.} A quantized 3.8B model
that clears a task's accuracy threshold, running on hardware that already exists
in the user's hand, is preferable on cost, latency, energy, and privacy to a
larger model that also clears it. A 3--4B model at INT4 occupies roughly 2~GB and
runs at interactive speed on a phone or a Jetson-class module. The scaling-era
assumption that a larger model is always the safer default is inverted for any
task whose requirements can be stated in advance.

\textbf{The evaluation record is the weak point, and it is weaker than the field
acknowledges.} Section~\ref{sec:reported} documents why a cross-model accuracy
comparison cannot be built from current public reporting: benchmark sets do not
intersect across vendors, shot counts vary per benchmark and per vendor, base and
instruction-tuned figures are conflated, and matched baselines are rarely
published. This is not a limitation peculiar to this survey. It applies equally
to the comparison tables circulated in secondary sources, community leaderboards,
and prior surveys, which are widely reproduced without these caveats and which
practitioners use to make procurement decisions. Establishing a shared,
independently run evaluation for this model class would be a more valuable
contribution to the field than another survey, and
Section~\ref{subsec:challenges_verification} specifies what it would require.

\subsection{Implications for Agentic Systems}
\label{subsec:discussion_agentic}

The agentic argument of Section~\ref{sec:agentic} rests on these claims rather
than restating them. Capability-per-parameter growth is what makes an \gls{slm}
\emph{sufficient} for a delegated sub-task; compression and edge deployability are
what make it \emph{cheap} and \emph{private}; right-sizing is the design
discipline that turns a heterogeneous SLM--LLM system from an option into the
default. Two cross-cutting observations sharpen the picture. First, routing must
be capability-aware, not size-aware: the reported multilingual and safety gaps
(Section~\ref{subsec:reported_safety}) mean that a naive ``small model unless
hard'' router will silently degrade non-English and safety-sensitive calls.
Second, adaptation matters more as models get smaller, so the economical unit in
an agent is often a \emph{fine-tuned} small model per recurring call type rather
than a single general \gls{slm}, which is exactly what the LLM-to-SLM conversion
pipeline of Section~\ref{subsec:agentic_conversion} produces.

\subsection{Democratization and Sustainability}
\label{subsec:discussion_implications}

Deploying a 3B model that answers grade-school mathematics correctly on an
\$80 single-board computer is a qualitative change in access. Organizations facing
cloud-API cost, data-sovereignty rules, or unreliable connectivity can now run
capable \gls{nlp} locally with no per-query cost and no data egress. The environmental
case is directional rather than precise: reported per-inference energy for a small
model on modest hardware is one to two orders of magnitude below a frontier-model
query in a datacenter, and if even a modest fraction of today's classification,
extraction, and short-generation traffic were served by right-sized \glspl{slm}, the
aggregate saving would be large. These figures are estimates and should be treated
as such.

\subsection{Limitations}
\label{subsec:discussion_limitations}

Beyond the evidence-base limitations already stated
(Section~\ref{subsec:reported_limits}): the review is English-centric in its task
coverage; the ecosystem moves faster than any static snapshot, so models released
after the collection window are discussed only qualitatively; the taxonomy's
boundaries (the 10B threshold, the treatment of \gls{moe} by active parameters) are
pragmatic conventions, not natural kinds; and the agentic argument, while
grounded in a production-scale position paper \citep{belcak2025slm_agentic} and two
large deployments, is a projection about where the field is heading, not a
measured outcome.

\subsection{Ethical Considerations}
\label{subsec:discussion_ethics}

Safety alignment is proportionally harder at small scale: several research-grade
\glspl{slm} lack the safety training of their larger siblings, and fine-tuning can erode
what alignment exists \citep{qi2023finetuning}. Practitioners deploying SLMs---and
especially those wiring them into agents with tool access---should treat
application-level safety filtering and output monitoring as mandatory, not
optional. Compression carries its own risk: pruning and quantization can degrade
performance unevenly across dialects, languages, and demographic groups
\citep{blodgett2020language}, and the reported advantage of quantization over
pruning is largest precisely on multilingual tasks \citep{kurz2025pruning_vs_quant},
which suggests pruning disproportionately harms non-English users. Fairness-aware
compression objectives and per-group performance floors for high-stakes
applications are open requirements. Against these risks, on-device \glspl{slm} advance
privacy (no data leaves the device), accessibility (no connectivity required),
and cost (no subscription), which matters most for the healthcare and
educational applications of Section~\ref{sec:applications} in under-resourced
settings.

% ================================================================
% Section: Challenges and Future Directions
% ================================================================

\section{Challenges and Future Directions}
\label{sec:challenges}

\subsection{Model-Level Challenges}
\label{subsec:challenges_model}
Hallucination is more pronounced at small scale, where fewer facts fit in the parameters, and retrieval augmentation is the main mitigation \citep{zhang2023siren,lewis2020rag}. Multi-step reasoning is where the gap to large models is widest; test-time compute, symbolic tools, and reasoning distillation narrow it. Hardware and model architecture are still designed separately, although \gls{npu} efficiency depends on matching attention and feed-forward shapes to the accelerator. Compression can amplify bias \citep{blodgett2020language,kurz2025pruning_vs_quant}. \suppsec~\ref{supp:challenges} treats each.

\subsection{Open Problems for Agentic SLM Systems}
\label{subsec:challenges_routing}

Several problems stand between hypothesis~H and routine practice. \textbf{Correctness checking:} H4b requires a check that catches valid but wrong outputs at a cost well below a large-model call; executors with observable effects and unit tests exist for code and tool calls, but most agent call types lack them. \textbf{Difficulty estimation without an oracle:} routers use win prediction \citep{ong2024routellm}, predicted difficulty \citep{ding2024hybridllm}, self-verification \citep{aggarwal2024automix}, or decomposition \citep{division_of_thoughts_2025}, but calibration under distribution shift is unsolved. \textbf{Multi-objective routing:} accuracy, latency, energy, cost, privacy sensitivity, and network availability must be weighed jointly, and size-only routing degrades non-English and safety-sensitive calls (Section~\ref{subsec:reported_safety}). \textbf{Handoff without a seam:} escalating mid-task must not produce a visible discontinuity. \textbf{Agent-appropriate, cost-aware evaluation:} function-calling and interaction benchmarks \citep{patil2025bfcl,yao2024taubench,barres2025tau2} and cost-aware agent leaderboards \citep{kapoor2025hal} exist, but report specialized sub-10B models sparsely and do not compare \gls{slm}-first with all-\gls{llm} agents. \textbf{Fleet management:} many fine-tuned specialists raise questions of versioning, drift detection, and retraining cadence that a single model does not.

\subsection{Empirical-Verification Agenda}
\label{subsec:challenges_verification}

Six studies would most reduce the uncertainty documented here. Studies 1--4 use one model set chosen by a stated rule: the instruction-tuned checkpoint of the most downloaded sub-10B variant on the Hugging Face Hub at the time of the study for each of the cohort families (Phi, Qwen, Gemma, Llama) and for the 2025--2026 families (for example Qwen3, Gemma~3n, SmolLM3, Granite~4, LFM2), so that the results bear on both the audited and the current generation. Each fixes the harness version, prompts, extraction rules, and decoding settings in advance, runs sampled tasks with at least three seeds, and reports confidence intervals.

\begin{description}
  \item[1. Harness-consistent reruns (RQ3).] Run the model set on the eight benchmarks of the coverage ledger under one \texttt{lm-evaluation-harness} configuration, separately for pre-trained and instruction-tuned checkpoints, so that every vendor figure in the ledger has an independent counterpart.
  \item[2. Quantization degradation (RQ2).] Measure FP16, INT8, and three INT4 formats (\gls{gguf} Q4\_K\_M, GPTQ, \gls{awq}) on the same benchmarks with matched baselines, to establish per-model degradation and whether it varies with size.
  \item[3. Edge measurement (RQ2).] Measure decode and prefill throughput, time-to-first-token, resident memory, and energy per token on a phone, a single-board computer, a Jetson-class module, and a laptop, reporting batch size, context length, and runtime version.
  \item[4. Agent-specific evaluation (RQ4).] Run the model set on BFCL multi-turn splits and $\tau$-bench, with and without constrained decoding, reporting pass\textsuperscript{$k$} consistency and classifying failures as invalid (caught by H4a) or valid but wrong (caught only by H4b).
  \item[5. Cost per completed task (RQ5, the test of H; Figure~\ref{fig:falsification}).] Fix a workload with a checkable success criterion, identify the call types meeting H1, H3, and H4 by inspection, and record pilot estimates of $N$, $F$, $a$, and $e$ for H2 before the comparison. Run six configurations end to end, each with the same checks and the same retry budget: (i)~a frontier model throughout; (ii)~the same with prompt caching and batch pricing where the provider offers them; (iii)~a hosted general-purpose small model or small-active-parameter \gls{moe} without specialization; (iv)~a hosted small model fine-tuned through the provider's service; (v)~a specialized \gls{slm} with escalation to the frontier model on a failed check, which is the configuration H prescribes; and (vi)~the specialized \gls{slm} without escalation, to measure what escalation buys. Report per configuration task success, calls and retries per completed task, escalation rate, the rate of valid but wrong outputs that pass the checks, latency, and monetary and energy cost per completed task, with $F$ reported separately. The decision rule is a pre-registered one-sided non-inferiority test of (v) against (i) on task success at margin $\delta$, with task-clustered bootstrap intervals, together with a comparison of cost per completed task against every other configuration. For independent arms at a success rate near 0.5, a one-sided test at $\alpha = 0.05$ with power 0.8 needs about $n \approx 3.09/\delta^2$ tasks per arm: about 7{,}700 for $\delta = 2$ percentage points, 1{,}240 for 5, and 310 for 10; a paired design on the same tasks needs fewer. A $\tau$-bench-scale workload of a few hundred tasks can therefore decide only a margin of about ten points, and a tighter margin needs a synthetic or production workload of thousands of tasks. Under the refutation rule of Section~\ref{subsec:intro_thesis}, H is rejected for a workload class if (v) fails non-inferiority or costs more per completed task than any other configuration, and rejected as a general rule if it fails on a majority of at least three pre-registered workloads with distinct call types.
  \item[6. Call-type prevalence (H2).] On logged traces of deployed agents, measure the share of calls, of cost, and of contribution to task success that meets H1--H4, replacing the three estimates of Section~\ref{subsec:agentic_delegation} with measurements.
\end{description}

\begin{figure}[t]
\centering
\footnotesize
\definecolor{fdInk}{HTML}{34404A}
\definecolor{fdBlue}{HTML}{7397B8}
\definecolor{fdSage}{HTML}{79A88E}
\definecolor{fdApr}{HTML}{D8A879}
\definecolor{fdLav}{HTML}{9B8FB5}
\definecolor{fdRose}{HTML}{C9858C}
\definecolor{fdSlate}{HTML}{A2AAB1}
\begin{tikzpicture}[
  >={Stealth[length=1.8mm]},
  col/.style={draw=#1!80!fdInk, fill=#1!10, rounded corners=2pt, align=left, text=fdInk,
              inner sep=3pt, font=\scriptsize},
  cfg/.style={draw=fdBlue!70!fdInk, fill=fdBlue!8, rounded corners=1.5pt, align=left, text=fdInk,
              inner sep=2pt, font=\scriptsize, text width=40mm, minimum height=5mm},
  hd/.style={font=\scriptsize\bfseries, text=fdInk},
  every path/.style={draw=fdInk!75},
]
\node[col=fdSlate, text width=34mm] (fix) {%
  \textbf{Fixed in advance}\\[2pt]
  workload with a checkable\\success criterion\\[1pt]
  call types meeting H1, H3, H4\\[1pt]
  pilot estimates of $N$, $F$, $a$, $e$\\[1pt]
  margin $\delta$ and tasks per arm:\\
  $n \approx 3.09/\delta^2$ ($\delta=$2, 5, 10 pts:\\
  about 7{,}700, 1{,}240, 310)};
\node[cfg, right=8mm of fix.north east, anchor=north west] (c1) {(i) frontier LLM throughout};
\node[cfg, below=1.2mm of c1] (c2) {(ii) frontier, cached and batched};
\node[cfg, below=1.2mm of c2] (c3) {(iii) hosted small model or small-active MoE};
\node[cfg, below=1.2mm of c3] (c4) {(iv) hosted fine-tuned small model};
\node[cfg, below=1.2mm of c4, draw=fdSage!85!fdInk, fill=fdSage!20, line width=0.8pt] (c5) {\textbf{(v) specialized SLM + escalation}\\{\tiny the configuration H prescribes}};
\node[cfg, below=1.2mm of c5] (c6) {(vi) specialized SLM, no escalation};
\node[hd, above=1mm of c1] {same checks and retry budget};
\node[col=fdLav, text width=40mm, right=8mm of c1.north east, anchor=north west] (met) {%
  \textbf{Reported per configuration}\\[2pt]
  task success; calls and retries\\per completed task; escalation rate;\\valid but wrong outputs passing\\the checks; latency; monetary and\\energy cost per completed task; $F$};
\node[col=fdRose, text width=40mm, below=3mm of met] (dec) {%
  \textbf{Decision rule}\\[2pt]
  1. (v) non-inferior to (i) on success\\ \phantom{1. }at margin $\delta$ (one-sided)\\
  2. (v) cheapest per completed task\\[2pt]
  H refuted for the workload if 1 or 2 fails;\\rejected in general if it fails on most\\of at least 3 workloads};
\draw[->] (fix.east |- c3.west) -- (c3.west);
\draw[->] (c3.east) -- (c3.east -| met.west);
\end{tikzpicture}
\caption{The pre-registered test of hypothesis~H (study 5). Call types, pilot estimates for H2, the margin, and the sample size are fixed before running; six configurations are run end to end under the same checks and retry budget; the decision rule combines a non-inferiority test on task success with a comparison of cost per completed task. The sample sizes assume independent arms at a success rate near 0.5, one-sided $\alpha = 0.05$, and power 0.8. No results are implied: the study has not been run.}
\Description{Three columns. Left: items fixed in advance, namely a workload with a checkable success criterion, call types meeting H1, H3, and H4, pilot estimates of N, F, a, and e, and the margin delta with tasks per arm of about 7,700, 1,240, and 310 for margins of 2, 5, and 10 points. Middle: six configurations run with the same checks and retry budget, from frontier model throughout to specialized SLM without escalation, with specialized SLM plus escalation highlighted. Right: metrics reported per configuration and a decision rule: non-inferiority on success and lowest cost per completed task, with refutation for the workload if either fails and general rejection if it fails on most of at least three workloads.}
\label{fig:falsification}
\end{figure}

Studies 1--4 are practical on one modern GPU in weeks; studies 5 and 6 need an agent harness, an API budget, and, for study 6, access to production traces. Study 5 is the one that bears directly on the paper's central hypothesis.

\subsection{Further Directions}
\label{subsec:challenges_future}
Beyond verification, \suppsec~\ref{supp:directions} sets out multilingual evaluation, on-device continual learning, ensembles of specialized \glspl{slm} against a single \gls{llm}, multimodal evaluation, linear-time hybrid architectures, and federated fine-tuning.

\section{Conclusion}
\label{section:conclusion}

This paper set out to do two things at once: survey the small-language-model
landscape as it stands in 2026, and argue a position about where it is going. On
the first, we cataloged 125 documented SLM variants under a four-dimensional taxonomy,
tracked the 2025--2026 shift toward hybrid state-space architectures, unified
reasoning models, and mobile-first multimodal families, reviewed the compression
techniques specific to the sub-10B regime---where quantization consistently
outperforms pruning \cite{kurz2025pruning_vs_quant}---and consolidated publicly
reported performance for a 17-model cohort across 15 tasks, six hardware tiers,
and three precisions, with every figure traced to a dated source rather than an
original run.

On the second: the ``less is more'' thesis is no longer only about efficiency. The
consumer of language intelligence that is growing fastest is the AI agent, and the
agent does not want a generalist. It issues many narrow, repetitive,
format-constrained calls, and for those calls a small specialized model is
sufficiently capable, 10--30$\times$ cheaper to serve, cheaper to adapt, and
easier to hold to a code interface \cite{belcak2025slm_agentic}. The consolidated
evidence in this paper supports each premise of that argument, and two
production-scale deployments---Apple's on-device/server split and the
routing-based systems surveyed in Section~\ref{sec:agentic}---already embody it.
The barriers that remain are those of sunk infrastructure, benchmark
mis-measurement, and attention, not of capability.

The work ahead is concrete. Reasoning still shows a performance cliff below
roughly 7B parameters that inference-time compute only partly closes;
hallucination is worse at small scale; compression can amplify bias; hardware and
model architecture are still designed separately; and the consolidated evidence
base needs controlled verification of the handful of claims that practitioners
most depend on (Section~\ref{sec:challenges}). None of these is a reason to wait.
The future of applied NLP is not one frontier model answering every query through
a cloud API. It is a fabric of small, specialized, mostly on-device models doing
the ordinary work, calling a large model only when the task genuinely demands it.
In that architecture, less is not a limitation. It is the design.

\section*{Declarations}
\label{section:declarations}
\addcontentsline{toc}{section}{Declarations}

\subsection*{Funding}
This work was supported and funded by the Deanship of Scientific Research at Imam Mohammad Ibn Saud Islamic University (IMSIU), grant number IMSIU-DDRSP2602.

\subsection*{Use of Generative AI}
A large language model assisted in compiling 205 of the candidate model records from public listings in March 2026 (Section~\ref{subsec:intro_method}). During revision, Anthropic's Claude was used to assist with editing the manuscript, checking bibliographic metadata against arXiv and Crossref records, writing the scripts that check each record's developer source and apply the inclusion rule, and transcribing the coverage ledger from the documents it cites. Every included record's developer source was checked through the Hugging Face Hub API or against the developer's own page, and the authors reviewed all generated output, verified factual claims and references against their sources, and take full responsibility for the content.

\subsection*{Competing Interests}
The authors declare that they have no competing interests.

\subsection*{Ethics Approval}
Not applicable: this is a review of publicly available documentation and involves no human participants, human data, or animals.

\subsection*{Data and Code Availability}
No benchmark was run for this review; all performance figures are reported values from public sources, whose provenance is given in \suppsec~\ref{app:provenance}. The following artifacts are provided as supplementary files with this submission and will be deposited in a versioned archival repository under a persistent identifier on publication: the variant inventory \file{merged\_llm\_database.json} (version~3.6), one record per variant with its parameter counts, C1/C2 status, inclusion decision and reason, developer source, the result and date of the source check, and a changelog naming every record added, moved, or removed between versions; the source-check log \file{source\_check\_log.csv}; the audit coverage ledger \file{audit\_coverage\_ledger.csv}; and the scripts that correct and verify the records, apply the inclusion rule, build the ledger, and generate every inventory table and figure (\file{fix\_records\_v3\_6.py}, \file{verify\_sources.py}, \file{apply\_inclusion\_rules.py}, \file{build\_audit\_ledger.py}, \file{generate\_model\_tables.py}, \file{generate\_figures\_v2.py}).

\subsection*{Author Contributions}
\textbf{M.G.R.:} Conceptualization, Methodology, Investigation, Data curation, Software, Writing (original draft), Visualization.
\textbf{H.A.:} Conceptualization, Methodology, Writing (review and editing), Supervision, Project administration.
\textbf{S.J.A.:} Investigation, Validation, Writing (review and editing).
\textbf{W.E.A.:} Validation, Writing (review and editing), Funding acquisition.
All authors read and approved the final manuscript.


\appendix
\section{Data Provenance and Supplementary Protocols}
\label{appendix:supplementary}
\label{app:provenance}

This appendix supports the reproducibility of the \emph{review itself}: it documents the sources of the reported figures consolidated in Sections~\ref{sec:reported}--\ref{sec:deployment}, and provides the reference prompt templates, quantization-pipeline specifications, and QLoRA hyperparameter settings under which those figures were reported. The appendix does not report original experimental runs.

\subsection{Per-Figure Provenance Table}
\label{subsec:appendix_provenance_table}
\textit{To be populated at submission:} a per-figure provenance table listing, for every reported performance value appearing in Sections~\ref{sec:reported}--\ref{sec:deployment}, the following fields: (a)~model (name + Hugging~Face commit hash or release version); (b)~benchmark and metric variant; (c)~quantization precision and toolchain; (d)~hardware configuration; (e)~source citation (vendor technical report, model card, leaderboard entry, or third-party study) with a stable identifier (DOI, arXiv ID, or archival URL); (f)~extraction date; and (g)~any conflict-resolution decision applied under the rules of Section~\ref{subsec:reported_sources}.

The remaining subsections of this appendix document the reference conditions---prompt formats, quantization toolchains, adaptation hyperparameters---that the reviewed sources most commonly report; these are preserved here so that readers can interpret any individual reported figure against the conditions it was measured under.

\subsection{Prompt Formats}
\label{subsec:appendix_prompt_templates}
The reviewed source materials commonly preserve each model's own instruction
wrapper while following a shared task-payload structure, which is summarized here
in prose. \textbf{Multiple-choice question answering} presents the question
followed by the enumerated options (A--D) and an ``Answer:'' cue, with the model
expected to emit a single option letter. \textbf{Mathematical reasoning} prepends
an instruction to solve the problem step by step and appends a ``Let's think step
by step'' cue, so that a chain-of-thought precedes the final numeric answer.
\textbf{Extractive question answering} supplies the context passage, then the
question, then an instruction to answer from the context and to respond
``unanswerable'' when the context does not contain the answer. \textbf{Abstractive
summarization} supplies the source article preceded by an instruction to produce a
two- to three-sentence summary. Where a source reports few-shot results, the same
fixed set of exemplars is prepended to the payload for every model on a given
task. Because sources differ in the exact wording of these wrappers and in
answer-extraction heuristics, the consolidated figures are compared only within
tasks and, wherever possible, within a single reporting source.

\subsection{Quantization and Compression Pipeline}
\label{subsec:appendix_quantization_pipeline}
To complement the methodology section, the compression workflows reported across the reviewed sources are summarized here. The reviewed literature reports five precision or compression settings across the consolidated comparison:
\begin{enumerate}
    \item \textbf{FP16 baseline:} standard half-precision inference used as the reference condition.
    \item \textbf{GGUF Q4\_K\_M:} mixed-precision 4-bit quantization via \texttt{llama.cpp}, assigning higher precision to the most sensitive tensors.
    \item \textbf{AWQ 4-bit:} activation-aware weight quantization with group size 128 and calibration on held-out text sequences.
    \item \textbf{GPTQ 4-bit:} layer-wise post-training quantization using second-order curvature approximations.
    \item \textbf{bitsandbytes NF4:} 4-bit NormalFloat quantization with double quantization, matching the QLoRA training regime.
\end{enumerate}

For compression stress tests, the sources report SparseGPT~\cite{frantar2023sparsegpt}-induced 50\% unstructured sparsity applied before 4-bit quantization. Reported figures for this combined pruning-plus-quantization setting are retained in the provenance tables for sensitivity analysis rather than the main comparative tables, because support for sparse inference remains hardware dependent and is not uniformly reported across the reviewed sources.

\subsection{Edge Measurement Protocols Reported in the Source Literature}
\label{subsec:appendix_edge_protocol}
The deployment figures consolidated in Sections~\ref{sec:reported}--\ref{sec:deployment} are drawn from source materials whose stated measurement protocols typically include a multi-minute warm-up to stabilize thermal and frequency behavior before measurement. Across the reviewed edge-deployment sources, the reported per-model, per-device measurements cover:
\begin{enumerate}
    \item \textbf{Time-to-first-token (TTFT):} reported latency from prompt submission to first generated token, typically averaged over 50 prompts with 128-token inputs.
    \item \textbf{Sustained throughput:} reported mean tokens-per-second over 256-token generations.
    \item \textbf{Peak memory:} reported maximum RAM or unified-memory consumption during generation.
    \item \textbf{Load time:} reported elapsed time from process start until the model is ready to serve requests.
    \item \textbf{Energy per inference:} reported device-level power estimates where the source provides counter- or telemetry-based measurements.
\end{enumerate}
Because these reporting conventions differ in small but potentially consequential ways across sources (warm-up duration, prompt-length distribution, power-counter granularity), readers should treat the consolidated figures as indicative rather than precisely comparable; the conflict-resolution rules in Section~\ref{subsec:reported_sources} are applied whenever discrepant protocols yield materially different reported values for the same tuple.

\subsection{Hyperparameter Search Space}
\label{subsec:appendix_hyperparameters}
\begin{table}[ht]
\centering
\caption{QLoRA hyperparameter configuration most commonly reported by sources.}
\label{tab:hyperparam_search}
\footnotesize
\begin{tabular}{lcc}
\toprule
\textbf{Hyperparameter} & \textbf{Search Space} & \textbf{Selected Value} \\
\midrule
LoRA rank $r$ & \{16, 32, 64, 128\} & 64 \\
LoRA scaling $\alpha$ & \{16, 32, 64, 128, 256\} & 128 \\
LoRA dropout & \{0.0, 0.05, 0.1\} & 0.05 \\
Learning rate & \{$10^{-5}$, $5\times10^{-5}$, $10^{-4}$, $2\times10^{-4}$, $5\times10^{-4}$\} & $2\times10^{-4}$ \\
Effective batch size & \{8, 16, 32\} & 16 \\
Max epochs & \{3, 5, 10\} & 5 \\
Warmup steps & \{50, 100, 200\} & 100 \\
Target modules & \{q,v\}, \{q,k,v,o\}, \{all linear\} & \{q,k,v,o\} \\
Weight decay & \{0.0, 0.01, 0.1\} & 0.01 \\
Scheduler & \{cosine, linear\} & cosine \\
\bottomrule
\end{tabular}
\end{table}

The right-hand column records the QLoRA configuration that the reviewed sources most commonly report; where a source used a different rank, target-module set, or schedule, the deviation is noted in the per-figure provenance table so that its reported gain is not compared directly against gains obtained under the common recipe.

\subsection{Extended Results Tables}
\label{subsec:appendix_extended_results}
\begin{table}[ht]
\centering
\caption{Representative five-shot results for selected models.}
\label{tab:appendix_fewshot}
\footnotesize\adjustbox{max width=\linewidth}{%
\begin{tabular}{lcccccc}
\toprule
\textbf{Model} & \textbf{ARC-C} & \textbf{HellaSwag} & \textbf{GSM8K} & \textbf{MMLU} & \textbf{SST-2} & \textbf{MNLI} \\
\midrule
Phi-4-mini-Inst. & 85.2 & 71.4 & 89.3 & 69.8 & 94.1 & 75.8 \\
Qwen2.5-7B-Inst. & 91.2 & 81.5 & 89.5 & 74.2 & 95.2 & 77.3 \\
LLaMA 3.2-3B-Inst. & 78.8 & 79.1 & 77.4 & 64.5 & 91.5 & 68.2 \\
SmolLM2-1.7B-Inst. & 63.1 & 63.5 & 44.8 & 53.4 & 86.4 & 59.8 \\
\bottomrule
\end{tabular}%
}
\end{table}

\begin{landscape}
\begin{table}[ht]
\centering
\caption{Extended zero-shot results across a broader task set (FP16, instruction-tuned variants). All values are reported figures.}
\label{tab:appendix_zero_shot_extended}
\footnotesize\adjustbox{max width=\linewidth}{%
\begin{tabular}{lcccccccccccc}
\toprule
\textbf{Model} & \textbf{ARC-E} & \textbf{ARC-C} & \textbf{HellaSwag} & \textbf{TriviaQA} & \textbf{SQuAD 2.0} & \textbf{GSM8K} & \textbf{MATH} & \textbf{HumanEval} & \textbf{MBPP} & \textbf{MMLU} & \textbf{SST-2} & \textbf{MNLI} \\
\midrule
Phi-4-mini-Inst. & 92.1 & 83.7 & 69.1 & 68.5 & 74.2 & 88.6 & 64.0 & 72.0 & 74.8 & 67.3 & 92.4 & 72.1 \\
Qwen2.5-7B-Inst. & 94.8 & 90.1 & 80.0 & 72.4 & 76.8 & 88.7 & 60.4 & 84.8 & 79.2 & 72.6 & 93.8 & 74.6 \\
Gemma 2-9B-Inst. & 95.2 & 89.8 & 87.1 & 71.8 & 75.5 & 84.9 & 51.3 & 68.9 & 73.5 & 71.3 & 93.5 & 73.8 \\
Gemma 3-4B-Inst. & 92.4 & 85.2 & 78.4 & 64.2 & 70.8 & 82.1 & 53.7 & 64.6 & 69.3 & 65.8 & 91.2 & 70.4 \\
LLaMA 3.2-3B-Inst. & 88.5 & 76.1 & 77.2 & 58.6 & 65.4 & 75.6 & 46.7 & 58.5 & 62.4 & 61.8 & 88.2 & 63.5 \\
LLaMA 3.2-1B-Inst. & 78.2 & 62.3 & 63.8 & 42.5 & 52.8 & 47.2 & 24.1 & 35.4 & 40.2 & 48.6 & 82.5 & 54.2 \\
Mistral 7B-Inst. v0.3 & 91.8 & 82.7 & 73.5 & 65.8 & 68.4 & 81.9 & 41.6 & 66.5 & 68.3 & 63.0 & 90.8 & 68.2 \\
Ministral 8B-Inst. & 93.2 & 85.4 & 76.8 & 67.2 & 71.2 & 83.5 & 45.2 & 69.5 & 71.8 & 65.7 & 92.1 & 71.5 \\
Qwen2.5-3B-Inst. & 90.2 & 78.6 & 72.1 & 58.8 & 66.5 & 79.8 & 50.2 & 65.2 & 67.8 & 65.6 & 89.5 & 66.8 \\
Qwen2.5-0.5B-Inst. & 68.4 & 52.1 & 48.3 & 28.5 & 38.2 & 38.5 & 18.7 & 22.0 & 28.4 & 47.5 & 78.2 & 48.5 \\
SmolLM2-1.7B-Inst. & 74.5 & 59.8 & 60.2 & 35.4 & 52.3 & 42.6 & 19.3 & 28.7 & 35.6 & 50.1 & 82.6 & 54.2 \\
SmolLM2-360M-Inst. & 55.8 & 38.2 & 42.5 & 18.2 & 32.4 & 15.8 & 5.2 & 8.5 & 12.3 & 33.4 & 72.4 & 42.8 \\
TinyLlama 1.1B & 62.4 & 47.5 & 53.8 & 28.8 & 42.5 & 18.4 & 6.8 & 12.2 & 16.8 & 36.2 & 76.5 & 46.2 \\
StableLM 2 1.6B & 72.8 & 55.2 & 57.6 & 32.5 & 48.6 & 32.1 & 12.4 & 22.5 & 28.9 & 44.3 & 80.2 & 52.5 \\
\bottomrule
\end{tabular}%
}
\end{table}
\end{landscape}

\subsection{Detailed Hardware Specifications}
\label{subsec:appendix_hardware_specs}
\begin{table}[ht]
\centering
\caption{Detailed hardware specifications for edge deployment experiments.}
\label{tab:appendix_hardware}
\footnotesize\adjustbox{max width=\linewidth}{%
\begin{tabular}{lcccc}
\toprule
\textbf{Specification} & \textbf{Raspberry Pi 5} & \textbf{Jetson Orin Nano} & \textbf{Galaxy S24+} & \textbf{Laptop} \\
\midrule
CPU & BCM2712 4$\times$A76 @ 2.4GHz & 6$\times$A78AE @ 1.5GHz & Snapdragon 8 Gen 3 & Ryzen 7 7840HS \\
GPU & VideoCore VII & 1024 CUDA (Ampere) & Adreno 750 & Radeon 780M \\
NPU & None & None & Hexagon & Ryzen AI \\
RAM & 8 GB LPDDR4X & 8 GB LPDDR5 (unified) & 12 GB LPDDR5X & 32 GB DDR5-5600 \\
Memory bandwidth & 25.6 GB/s & 68 GB/s & 51.2 GB/s & 89.6 GB/s \\
Storage & microSD / NVMe & NVMe & UFS 4.0 & NVMe SSD \\
TDP & 12W & 25W & $\sim$8W sustained & 35--54W \\
Inference engine & llama.cpp & MLC LLM & ExecuTorch & llama.cpp \\
\bottomrule
\end{tabular}%
}
\end{table}

\subsection{Human-Evaluation Conventions in the Source Literature}
\label{subsec:appendix_human_eval}
For subjective tasks such as summarization quality and instruction following, the reviewed sources that report human judgments typically rate outputs along three axes---\textbf{relevance}, \textbf{coherence}, and \textbf{factual accuracy}---on a 1--5 Likert scale, with model identity hidden from annotators and inter-annotator reliability reported via a Krippendorff-style coefficient. MT-Bench-style results substitute a strong LLM judge for human raters. Because these protocols differ across sources in rubric wording and judge model, the consolidated instruction-following column (MT-Bench) should be read as indicative, and human-rated summarization scores are reported only as ranges.

\section{Broader Small-Model Ecosystem Context}
\label{appendix:ecosystem}

This appendix situates the reviewed \gls{slm} cohort within the wider model ecosystem, which matters because the 2025--2026 release cycle moved compact open models, open-weight reasoning models, and frontier API systems in parallel rather than in sequence. The summary below records the ecosystem context needed to interpret the cohort, followed by the full generated variant inventory.

\subsection{Landscape Snapshot}
\label{subsec:ecosystem_snapshot}
Three ecosystem trends are especially relevant to the present study:
\begin{description}
    \item[Capability compression] Reasoning, long-context handling, and multimodality have moved into the 1B--8B range, narrowing the gap between deployable \glspl{slm} and cloud-scale \glspl{llm}.
    \item[Open-weight acceleration] Families such as Qwen, Phi, Gemma, SmolLM, DeepSeek distills, Granite, and MiniCPM have expanded the design space for self-hosted deployment (Figure~\ref{fig:license_posture}).
    \item[Deployment diversification] Inference targets now span cloud GPUs, consumer GPUs, laptops, smartphones, browsers, and embedded edge boards, making model choice inseparable from hardware choice.
    \item[The small-active-parameter turn] The most distinctive feature of the mid-2026 wave is how much of it activates 3--13B parameters against totals of 30B to 284B: Nemotron 3.5 Lightning (30B/3B), Hy-MT2-30B-A3B, North Mini Code (30B/3B), Laguna XS 2.1 (33B/3B), Ling 3.0 Flash (124B/5.1B) and its domain variants, and DeepSeek V4 Flash (284B/13B). These are the models that make the C1/C2 distinction load-bearing rather than pedantic, since each is cheap per token and none is resident on a phone.
\end{description}

The May--September 2026 additions to the inventory illustrate both ends of the range. At the small end, IBM shipped Granite 4.1 and then Granite 4.2 in 3B, 8B, and 30B dense sizes with a thinking switch, Liquid AI extended LFM2.5 from 230M to an 8.3B/1.5B on-device \gls{moe}, and Tencent released the Hy-MT2 translation family at 1.8B and 7B. At the frontier, OpenAI moved from GPT-5.6 (Sol, Terra, and Luna tiers) to GPT-6 Astra, Anthropic shipped Opus 4.8, Fable 5, Sonnet 5, Opus 5, and Fable 5.1 in a single quarter, Google iterated Gemini Flash from 3.1 to 3.8, and Meta's restructured Superintelligence Labs released the proprietary Muse Spark line alongside Muse Glimmer 30B, its first open-weight model since Llama~4.

Figure~\ref{fig:family_treemap} gives a size-proportional overview of the inventory by developer and family, and Figure~\ref{fig:developer_diversity} tracks the entry of new developers into the open-SLM space.

\begin{figure}[ht]
  \centering
  \includegraphics[width=\linewidth]{paper_slm_family_treemap.png}
  \caption{Documented SLM variants grouped by developer and family; area is proportional to the number of variants. Generated from \file{merged\_llm\_database.json} (snapshot 5~September 2026).}
  \Description{Treemap of documented small language model variants grouped by developer and then by family, with rectangle area proportional to the number of variants in each group.}
  \label{fig:family_treemap}
\end{figure}

\begin{figure}[ht]
  \centering
  \includegraphics[width=\linewidth]{paper_slm_developer_diversity_timeline.png}
  \caption{Cumulative distinct developers releasing documented SLM variants (line) and new developers per quarter (bars).}
  \Description{Combination chart with bars showing the number of new developers releasing their first documented small language model in each quarter, overlaid with a line showing the cumulative count of distinct developers.}
  \label{fig:developer_diversity}
\end{figure}

\begin{figure}[ht]
  \centering
  \includegraphics[width=0.78\linewidth]{paper_slm_license_posture.png}
  \caption{License posture of documented SLM variants. Permissive open licenses (Apache-2.0, MIT) account for the plurality; ``unspecified'' entries are typically research checkpoints without an attached license file.}
  \Description{Bar chart of documented small language model variants by license category. Permissive open licenses, chiefly Apache 2.0 and MIT, form the plurality, followed by custom vendor licenses and unspecified research checkpoints.}
  \label{fig:license_posture}
\end{figure}

\subsection{Representative Family Reference}
\label{subsec:ecosystem_families}
\begin{table}[ht]
\centering
\caption{Representative model families tracked in the inventory.}
\label{tab:ecosystem_families}
\footnotesize\adjustbox{max width=\linewidth}{%
\begin{tabular}{lllll}
\toprule
\textbf{Family} & \textbf{Developer} & \textbf{Representative Versions} & \textbf{Access} & \textbf{Typical Role} \\
\midrule
Qwen2.5 / Qwen3 & Alibaba & 0.5B--14B dense, MoE variants & Open weights & Strong multilingual, math, and code performance \\
Phi-3 / Phi-4 & Microsoft & 3.8B mini variants, 14B reasoning & Open weights & Data-curated reasoning-focused SLMs \\
Gemma 2 / 3 / 3n & Google DeepMind & 1B--12B dense, mobile-first multimodal variants & Open weights & Distilled compact models with strong safety tuning \\
Llama 3.2 / 4 & Meta & 1B--3B edge models, larger MoE frontier models & Open weights & Broad deployment ecosystem from edge to server \\
SmolLM2 / SmolLM3 & Hugging Face & 135M--3B & Open weights & Lightweight, fully open research and edge baselines \\
DeepSeek distills & DeepSeek & 1.5B--14B distilled reasoning models & Open weights & Reasoning transfer into self-hosted sizes \\
Granite 3 / 4 & IBM & 2--8B dense; 4.0 hybrid Mamba-2/MoE (350M--32B, 1--9B active) & Open weights & Enterprise-friendly, tool-use and RAG focus \\
Falcon-H1 & TII & 0.5--34B parallel attention/Mamba-2 hybrid & Open weights & Long-context edge-to-server \\
LFM2 / LFM2.5 & Liquid AI & 230M--8.3B hybrid conv/attention, incl.\ 8B-A1B MoE & Open weights & Fastest-on-CPU on-device models \\
MiniCPM & OpenBMB & Compact multimodal variants & Open weights & Edge multimodal use cases \\
Ling / Nemotron 3.5 / Laguna & inclusionAI / NVIDIA / Poolside & Small-active MoE, 3--13B active & Open weights & SLM-like per-token cost at LLM-like residency \\
GPT-5.6 / GPT-6 / Claude 5 / Gemini 3.x & OpenAI / Anthropic / Google & API frontier models & API only & High-end proprietary baselines and teacher models \\
Grok / Kimi / GLM / Muse / Nova & xAI / Moonshot / Zhipu / Meta / AWS & Large API or MoE systems & Mixed & Frontier alternatives with long context or agent focus \\
\bottomrule
\end{tabular}%
}
\end{table}

\subsection{Classification Conventions}
\label{subsec:ecosystem_classification}
The inventory applies conditions C1 and C2 of Section~\ref{subsubsec:intro_def_operative} mechanically, from the parameter counts each record carries, rather than by editorial judgement:
\begin{itemize}
    \item A record enters the \gls{slm} inventory when its \emph{active} parameter count is at or below roughly 20B (C1). Dense models, for which active equals total, therefore enter at or below 20B total.
    \item Within the inventory, C2 (total at or below 20B) separates \glspl{slm} proper from small-active-parameter \gls{moe} models. Both are listed; only the former carry the paper's storage, resident-memory, and on-device claims.
    \item Where a source states no active-parameter count for a model it describes as mixture-of-experts, the record is marked undisclosed rather than having active silently set equal to total, which would misclassify a small-active model as a large one on no evidence.
    \item Closed API models with undisclosed parameter counts are treated as frontier \gls{llm} baselines, with five exceptions retained on explicit vendor positioning as small models and named in Section~\ref{subsec:ecosystem_generated_inventory}.
\end{itemize}

Applying these conditions mechanically is what changed the headline count between drafts, from 126 to 215; Section~\ref{subsubsec:intro_def_terms} records the reconciliation and the released changelog lists it per model.

\subsection{Generated Variant Inventory}
\label{subsec:ecosystem_generated_inventory}
The following inventory is generated directly from \file{merged\_llm\_database.json} using \file{scripts/generate\_model\_tables.py}. The database covers 215 \gls{slm} variants and 224 \gls{llm} variants across 113 families (439 entries in total, snapshot 5~September 2026); only the \gls{slm} inventory is reproduced here, organized chronologically by release year and then by family, with parameter size as an explicit comparison column. The corresponding \gls{llm} inventory is available in the source database.

Each record in the released database carries the fields \texttt{meets\_c1\_active\_le\_20b}, \texttt{meets\_c2\_total\_le\_20b} and \texttt{slm\_class}, so a reader can recompute the inventory under a different threshold or restrict it to models satisfying both conditions. Of the 215 entries reproduced here, 183 carry \texttt{slm\_class}~$=$~\texttt{slm}, 27 carry \texttt{small-active-moe}, and 5 carry \texttt{undisclosed}. The \texttt{undisclosed} entries are vendor-positioned small models (Amazon Nova Micro and Lite, GPT-4o mini, Voxtral Mini, Nex-N2-Mini) for which no parameter count has been published; they are listed by name in the database changelog rather than assigned an assumed size, and no memory or on-device claim in this paper rests on them.

{\footnotesize\noindent\textbf{Chronology.} \hyperref[tab:model_inventory_slm_by_year]{SLM inventory}: \hyperlink{model-year-slm-2026}{2026}, \hyperlink{model-year-slm-2025}{2025}, \hyperlink{model-year-slm-2024}{2024}, \hyperlink{model-year-slm-2023-and-earlier}{2023 and earlier}.\par}
% Auto-generated by scripts/generate_model_tables.py. Do not edit manually.
\begin{landscape}
{\scriptsize
\setlength{\tabcolsep}{4pt}
\renewcommand{\arraystretch}{1.15}
\begin{longtable}{@{}L{3.3cm}C{1.05cm}L{1.75cm}C{1.15cm}L{2.25cm}L{2.4cm}C{1.6cm}L{3.9cm}@{}}
\caption{Generated chronological small language model inventory from \texttt{merged\_llm\_database.json}, grouped by release year and family.}\label{tab:model_inventory_slm_by_year}\\
\toprule
\textbf{Model} & \textbf{Access} & \textbf{Params} & \textbf{Context} & \textbf{Modalities} & \textbf{License} & \textbf{Release} & \textbf{Notes} \\ \midrule
\endfirsthead
\caption[]{Generated chronological small language model inventory continued.}\\
\toprule
\textbf{Model} & \textbf{Access} & \textbf{Params} & \textbf{Context} & \textbf{Modalities} & \textbf{License} & \textbf{Release} & \textbf{Notes} \\ \midrule
\endhead
\midrule \multicolumn{8}{r@{}}{\footnotesize\itshape continued on next page} \\
\endfoot
\bottomrule
\endlastfoot
\multicolumn{8}{@{}l}{\hypertarget{model-year-slm-2026}{\textbf{2026}}} \\ \midrule
\multicolumn{8}{@{}l}{\textit{Qwen3.5} --- Alibaba (Qwen3.5 (Mar 2026)) \href{https://qwen.ai/blog?id=qwen3.5}{[family page]}} \\*
\href{https://huggingface.co/Qwen/Qwen3.5-0.8B}{Qwen3.5-0.8B} & OW & 0.8B & 256K & Text+Vision+Video & Apache 2.0 & 2026-03-02 & ultra-lightweight multimodal; on-device \\
\href{https://huggingface.co/Qwen/Qwen3.5-2B}{Qwen3.5-2B} & OW & 2B & 256K & Text+Vision+Video & Apache 2.0 & 2026-03-02 & on-device multimodal chat; edge AI \\
\href{https://huggingface.co/Qwen/Qwen3.5-4B}{Qwen3.5-4B} & OW & 4B & 256K & Text+Vision+Video & Apache 2.0 & 2026-03-02 & mobile multimodal assistant; agentic coding \\
\href{https://huggingface.co/Qwen/Qwen3.5-9B}{Qwen3.5-9B} & OW & 9B & 256K & Text+Vision+Video & Apache 2.0 & 2026-03-02 & strong multimodal SLM; coding \\
\href{https://huggingface.co/Qwen/Qwen3.5-35B-A3B}{Qwen3.5-35B-A3B} & OW & 35B (3B act.) & 256K & Text+Vision & Apache 2.0 & 2026-02-16 & efficient MoE for edge/server; multimodal reasoning \\
\href{https://huggingface.co/Qwen/Qwen3.5-122B-A10B}{Qwen3.5-122B-A10B} & OW & 122B (10B act.) & 256K & Text+Vision & Apache 2.0 & 2026-02-16 & strong MoE reasoning; locally runnable flagship \\
\href{https://huggingface.co/Qwen/Qwen3.5-397B-A17B}{Qwen3.5-397B-A17B} & Both & 397B (17B act.) & 256K & Text+Vision & Apache 2.0 & 2026-02-16 & flagship open-weight multimodal; 201 languages \\
\multicolumn{8}{@{}l}{\textit{AFM} --- Apple (AFM) \href{https://machinelearning.apple.com/research/introducing-apple-foundation-models}{[family page]}} \\*
\href{https://machinelearning.apple.com/research/introducing-third-generation-of-apple-foundation-models}{Apple AFM 3 Core Advanced} & API & 20B (1-4B act.) & 64K & Text+Vision & Proprietary & 2026-06-08 & on-device assistant; multimodal understanding \\
\multicolumn{8}{@{}l}{\textit{Command} --- Cohere (Command) \href{https://cohere.com/blog/introducing-embed-v3}{[family page]}} \\*
\href{https://huggingface.co/CohereLabs/North-Mini-Code-1.0}{North Mini Code 1.0} & OW & 30B (3B act.) & 250K & Text & CC-BY-NC (Cohere Labs) & 2026-06-17 & coding; agentic coding \\
\multicolumn{8}{@{}l}{\textit{DeepSeek} --- DeepSeek (DeepSeek) \href{https://github.com/deepseek-ai/DeepSeek-Coder}{[family page]}} \\*
\href{https://huggingface.co/deepseek-ai/DeepSeek-V4-Flash-Vision-Exp}{DeepSeek V4 Flash Vision Exp} & OW & 284B (13B act.) & 1M & Text+Vision & See source & 2026-08-21 & vision; long context \\
\href{https://huggingface.co/deepseek-ai/DeepSeek-V4-Flash-0731}{DeepSeek V4 Flash 0731} & OW & 284B (13B act.) & 1280K & Text & See source & 2026-07-31 & general; long context \\
\multicolumn{8}{@{}l}{\textit{Dots} --- Dots Studio (Dots3-Note Preview (Aug 2026)) \href{https://openrouter.ai/dots-studio/dots-3-note-preview}{[family page]}} \\*
\href{https://openrouter.ai/dots-studio/dots-3-note-preview}{Dots3-Note Preview} & OW & 280B (16B act.) & 500K & Text & See source & 2026-08-14 & note taking; general \\
\multicolumn{8}{@{}l}{\textit{Gemma 4} --- Google DeepMind} \\*
\href{https://huggingface.co/google/gemma-4-E2B}{Gemma 4 E2B} & OW & \textasciitilde{}5B (\textasciitilde{}2B effective act.) & 128K & Text+Vision+Audio & Apache 2.0 & 2026-04-02 & on-device; multimodal \\
\href{https://huggingface.co/google/gemma-4-26B}{Gemma 4 26B (MoE)} & OW & 26B (3.8B act.) & 128K & Text+Vision & Apache 2.0 & 2026-04-02 & reasoning; multimodal \\
\href{https://huggingface.co/google/gemma-4-E4B}{Gemma 4 E4B} & OW & \textasciitilde{}8B (\textasciitilde{}4B effective act.) & 128K & Text+Vision+Audio & Apache 2.0 & 2026-04-02 & multimodal; reasoning \\
\href{https://huggingface.co/google/gemma-4-12b}{Gemma 4 12B} & OW & 12B & 256K & Text+Vision+Audio & Apache 2.0 & 2026-04-02 & multimodal; reasoning \\
\multicolumn{8}{@{}l}{\textit{Granite 4} --- IBM} \\*
\href{https://huggingface.co/ibm-granite/granite-4.2-3b}{Granite 4.2 3B} & OW & 3B & 128K & Text & Apache 2.0 & 2026-08-25 & reasoning; on-device \\
\href{https://huggingface.co/ibm-granite/granite-4.2-8b}{Granite 4.2 8B} & OW & 8B & 128K & Text & Apache 2.0 & 2026-08-25 & reasoning; agentic \\
\href{https://huggingface.co/blog/ibm-granite/granite-4-2}{Granite 4.1 3B} & OW & 3B & 128K & Text & Apache 2.0 & 2026-04-29 & enterprise; on-device \\
\href{https://huggingface.co/blog/ibm-granite/granite-4-2}{Granite 4.1 8B} & OW & 8B & 128K & Text & Apache 2.0 & 2026-04-29 & enterprise; agentic \\
\multicolumn{8}{@{}l}{\textit{LFM2} --- Liquid AI} \\*
\href{https://huggingface.co/LiquidAI/LFM2.5-2.6B}{LFM2.5-2.6B} & Both & 2.69B & 128K & Text & LFM1.0 & 2026-08-04 & agentic workflows; data extraction \\
\href{https://www.liquid.ai/blog/lfm2-5-8b-a1b}{LFM2.5-8B-A1B} & OW & 8.3B (1.5B act.) & 128K & Text & LFM1.0 & 2026-05-28 & on-device; tool calling \\
\href{https://www.liquid.ai/blog/lfm2-5-350m-no-size-left-behind}{LFM2.5-350M} & OW & 350M & 128K & Text & LFM1.0 & 2026 & on-device; edge \\
\href{https://www.liquid.ai/models}{LFM2.5-1.2B} & OW & 1.2B & 128K & Text & LFM1.0 & 2026 & on-device; edge \\
\href{https://www.liquid.ai/models}{LFM2.5-Audio-1.5B} & OW & 1.5B & -- & Text+Audio & LFM1.0 & 2026 & on-device speech; transcription \\
\href{https://www.liquid.ai/models}{LFM2.5-VL-3B} & OW & 3B & -- & Text+Vision & LFM1.0 & 2026 & on-device VLM; document extraction \\
\multicolumn{8}{@{}l}{\textit{Mistral Small} --- Mistral AI (Mistral Small 4 (Mar 2026)) \href{https://mistral.ai/news/mistral-small-4}{[family page]}} \\*
\href{https://huggingface.co/mistralai/Mistral-Small-4-119B-2603}{Mistral Small 4 (119B)} & Both & 119B (6.5B act.) & 256K & Text+Vision & Apache 2.0 & 2026-03-16 & hybrid reasoning; coding \\
\multicolumn{8}{@{}l}{\textit{NVIDIA Nemotron} --- NVIDIA (Nemotron 3 Super (Mar 2026)) \href{https://developer.nvidia.com/blog/introducing-nemotron-3-super-an-open-hybrid-mamba-transformer-moe-for-agentic-reasoning/}{[family page]}} \\*
\href{https://huggingface.co/nvidia/NVIDIA-NemotronLabs-VoiceChat-11B}{NemotronLabs VoiceChat 11B} & OW & 11B & 128K & Text+Audio & NVIDIA Custom & 2026-08-03 & voice interaction; on-device speech \\
\href{https://huggingface.co/nvidia/NVIDIA-Nemotron-3-Super-120B-A12B-BF16}{Nemotron 3 Super} & OW & 120B (12B act.) & 1M & Text & NVIDIA Open Model License & 2026-03 & open hybrid Mamba-Transformer MoE; agentic reasoning \\
\multicolumn{8}{@{}l}{\textit{Nemotron} --- NVIDIA (Nemotron) \href{https://arxiv.org/abs/2402.16819}{[family page]}} \\*
\href{https://huggingface.co/nvidia/NVIDIA-Nemotron-3.5-Lightning-30B-A3B-BF16}{Nemotron 3.5 Lightning 30B-A3B} & OW & 30B (3B act.) & 256K & Text & NVIDIA Open Model License & 2026-08-11 & agentic execution layer; high-volume tool calling \\
\href{https://huggingface.co/nvidia/Nemotron-3.5-Content-Safety}{Nemotron 3.5 Content Safety} & OW & 4B & 128K & Text+Vision & NVIDIA Open Model License & 2026-06-04 & guardrails; content moderation \\
\multicolumn{8}{@{}l}{\textit{Nex-N2} --- Nex AGI (Nex-N2-Mini (Jun 2026)) \href{https://huggingface.co/nex-agi/Nex-N2-Pro}{[family page]}} \\*
\href{https://huggingface.co/nex-agi/Nex-N2-Mini}{Nex-N2-Mini} & OW & Undisclosed & 256K & Text+Vision & See source & 2026-06-24 & agentic; coding \\
\href{https://huggingface.co/nex-agi/Nex-N2-Pro}{Nex-N2-Pro} & OW & 397B (17B act.) & 256K & Text+Vision & See source & 2026-06-08 & agentic; coding \\
\multicolumn{8}{@{}l}{\textit{Laguna} --- Poolside (Laguna S 2.1 (Jul 2026)) \href{https://poolside.ai/}{[family page]}} \\*
\href{https://huggingface.co/poolside/Laguna-S-2.1}{Laguna S 2.1} & OW & 118B (8B act.) & 1M & Text & See source & 2026-07-21 & coding agent \\
\href{https://huggingface.co/poolside/Laguna-XS-2.1}{Laguna XS 2.1} & OW & 33B (3B act.) & 256K & Text & See source & 2026-07-02 & coding agent \\
\multicolumn{8}{@{}l}{\textit{Hunyuan} --- Tencent (Hunyuan) \href{https://huggingface.co/tencent/Tencent-Hunyuan-Large}{[family page]}} \\*
\href{https://huggingface.co/tencent/Hy-MT2-1.8B}{Hy-MT2-1.8B} & OW & 1.8B & 8K & Text & Tencent Hunyuan License & 2026-08-20 & translation; on-device \\
\href{https://huggingface.co/tencent/Hy-MT2-30B-A3B}{Hy-MT2-30B-A3B} & OW & 30B (3B act.) & 8K & Text & Tencent Hunyuan License & 2026-08-20 & translation \\
\href{https://huggingface.co/tencent/Hy-MT2-7B}{Hy-MT2-7B} & OW & 7B & 8K & Text & Tencent Hunyuan License & 2026-08-19 & translation; Lang: 33 language pairs + 5 dialect/minority pairs \\
\multicolumn{8}{@{}l}{\textit{Inkling} --- Thinking Machines Lab (Inkling Small (Jul 2026)) \href{https://huggingface.co/thinkingmachines/Inkling}{[family page]}} \\*
\href{https://huggingface.co/thinkingmachines/Inkling-Small}{Inkling Small} & OW & 276B (12B act.) & 1M & Text+Vision & See source & 2026-07-30 & general; multimodal \\
\multicolumn{8}{@{}l}{\textit{MiMo} --- Xiaomi (MiMo-V2-Pro (Mar 2026)) \href{https://mimo.xiaomi.com/mimo-v2-pro}{[family page]}} \\*
\href{https://huggingface.co/XiaomiMiMo/MiMo-V2-Flash}{MiMo-V2-Flash} & OW & 309B (15B act.) & 256K & Text & Undisclosed & 2026 & high-throughput coding agents; production serving \\
\multicolumn{8}{@{}l}{\textit{Ling} --- inclusionAI (Ling-1T (2026)) \href{https://huggingface.co/01ai/Ling-1T}{[family page]}} \\*
\href{https://openrouter.ai/inclusionai/ling-3.0-flash-sante}{Ling 3.0 Flash Sante} & API & 124B (5.1B act.) & 256K & Text & See source & 2026-09-04 & health; medicine \\
\href{https://openrouter.ai/inclusionai/ling-3.0-flash-fin}{Ling 3.0 Flash Fin} & API & 124B (5.1B act.) & 256K & Text & See source & 2026-08-27 & finance \\
\href{https://huggingface.co/inclusionAI/Ling-3.0-flash}{Ling 3.0 Flash} & OW & 124B (5.1B act.) & 256K & Text & See source & 2026-07-23 & general; agentic \\
\addlinespace[3pt]
\multicolumn{8}{@{}l}{\hypertarget{model-year-slm-2025}{\textbf{2025}}} \\ \midrule
\multicolumn{8}{@{}l}{\textit{Qwen} --- Alibaba (Qwen) \href{https://github.com/QwenLM/Qwen}{[family page]}} \\*
\href{https://huggingface.co/Qwen/Qwen2.5-Omni-7B}{Qwen2.5-Omni-7B} & OW & 7B & -- & Text+Audio+Video+Image & See source & 2025-03-26 & omni-modal; text and audio output \\
\href{https://huggingface.co/Qwen/Qwen2.5-VL-7B-Instruct}{Qwen2.5-VL-7B} & OW & 7B & -- & Text+Video+Image & See source & 2025-01-22 & active \\
\multicolumn{8}{@{}l}{\textit{Qwen3} --- Alibaba (Qwen3 (Apr 2025)) \href{https://github.com/QwenLM/Qwen3}{[family page]}} \\*
\href{https://huggingface.co/Qwen/Qwen3-4B-Instruct-2507}{Qwen3-4B-Instruct-2507} & OW & 4B & 128K & Text & Apache 2.0 & 2025-07 & improved instruction following; non-thinking mode optimised \\
\href{https://huggingface.co/Qwen/Qwen3-0.6B}{Qwen3-0.6B} & OW & 0.6B & 32K & Text & Apache 2.0 & 2025-04-29 & edge inference; IoT \\
\href{https://huggingface.co/Qwen/Qwen3-1.7B}{Qwen3-1.7B} & OW & 1.7B & 32K & Text & Apache 2.0 & 2025-04-29 & mobile inference; Q\&A \\
\href{https://huggingface.co/Qwen/Qwen3-30B-A3B}{Qwen3-30B-A3B} & OW & 30B (3B act.) & 128K & Text & Apache 2.0 & 2025-04-29 & MoE efficiency; broad capability in tiny footprint \\
\href{https://huggingface.co/Qwen/Qwen3-4B-Instruct}{Qwen3-4B-Instruct} & OW & 4B & 128K & Text & Apache 2.0 & 2025-04-29 & on-device chat; coding \\
\href{https://huggingface.co/Qwen/Qwen3-8B-Instruct}{Qwen3-8B-Instruct} & OW & 8B & 128K & Text & Apache 2.0 & 2025-04-29 & single-GPU inference; reasoning \\
\href{https://huggingface.co/Qwen/Qwen3-14B-Instruct}{Qwen3-14B-Instruct} & OW & 14B & 128K & Text & Apache 2.0 & 2025-04-29 & high-quality on-premise chat; coding \\
\multicolumn{8}{@{}l}{\textit{AFM} --- Apple (AFM) \href{https://machinelearning.apple.com/research/introducing-apple-foundation-models}{[family page]}} \\*
\href{https://arxiv.org/abs/2507.13575}{Apple Foundation Model (On-Device, 2nd gen)} & API & \textasciitilde{}3B & 64K & Text+Vision & Proprietary & 2025-06-09 & on-device assistant; writing tools \\
\multicolumn{8}{@{}l}{\textit{DeepSeek R1-Distill} --- DeepSeek (R1-Distill (Jan 2025)) \href{https://huggingface.co/deepseek-ai}{[family page]}} \\*
\href{https://huggingface.co/deepseek-ai/DeepSeek-R1-Distill-Qwen-1.5B}{DeepSeek-R1-Distill-Qwen-1.5B} & OW & 1.5B & 128K & Text & MIT & 2025-01-20 & reasoning at ultra-small scale; math \\
\href{https://huggingface.co/deepseek-ai/DeepSeek-R1-Distill-Qwen-7B}{DeepSeek-R1-Distill-Qwen-7B} & OW & 7B & 128K & Text & MIT & 2025-01-20 & chain-of-thought reasoning; coding \\
\href{https://huggingface.co/deepseek-ai/DeepSeek-R1-Distill-Llama-8B}{DeepSeek-R1-Distill-Llama-8B} & OW & 8B & 128K & Text & MIT & 2025-01-20 & reasoning on Llama 3 base; on-premise \\
\href{https://huggingface.co/deepseek-ai/DeepSeek-R1-Distill-Qwen-14B}{DeepSeek-R1-Distill-Qwen-14B} & OW & 14B & 128K & Text & MIT & 2025-01-20 & strong reasoning 14B; top distill for coding \\
\multicolumn{8}{@{}l}{\textit{Gemma 3} --- Google DeepMind (Gemma 3 (Mar 2025)) \href{https://huggingface.co/blog/gemma3}{[family page]}} \\*
\href{https://huggingface.co/google/gemma-3-270m-it}{Gemma 3 270M-IT} & OW & 270M & 32K & Text & Google Gemma License & 2025-03 & ultra-lightweight prototyping; experiments \\
\href{https://huggingface.co/google/gemma-3-1b-it}{Gemma 3 1B-IT} & OW & 1B & 32K & Text & Google Gemma License & 2025-03 & on-device mobile inference; lightweight Q\&A \\
\href{https://huggingface.co/google/gemma-3-4b-it}{Gemma 3 4B-IT} & OW & 4B & 128K & Text+Vision & Google Gemma License & 2025-03 & multimodal on-device; image understanding \\
\href{https://huggingface.co/google/gemma-3-12b-it}{Gemma 3 12B-IT} & OW & 12B & 128K & Text+Vision & Google Gemma License & 2025-03 & high-quality assistant; 140+ languages \\
\multicolumn{8}{@{}l}{\textit{Gemma 3n} --- Google DeepMind (Gemma-3n (2025)) \href{https://ai.google.dev/gemma/docs/gemma-3n}{[family page]}} \\*
\href{https://huggingface.co/google/gemma-3n-E2B-it}{Gemma-3n-E2B-IT} & OW & \textasciitilde{}5B (\textasciitilde{}2B effective act.) & 32K & Text+Vision+Audio+Video & Google Gemma License & 2025 & mobile-first multimodal; on-device speech/vision/text \\
\href{https://huggingface.co/google/gemma-3n-E4B-it}{Gemma-3n-E4B-IT} & OW & \textasciitilde{}8B (\textasciitilde{}4B effective act.) & 32K & Text+Vision+Audio+Video & Google Gemma License & 2025 & higher-accuracy mobile multimodal; edge devices \\
\multicolumn{8}{@{}l}{\textit{SmolLM} --- Hugging Face (SmolLM3-3B (2025)) \href{https://huggingface.co/HuggingFaceTB/SmolLM3-3B}{[family page]}} \\*
\href{https://huggingface.co/HuggingFaceTB/SmolLM3-3B}{SmolLM3-3B} & OW & 3B & 64K & Text & Apache 2.0 & 2025 & dual-mode reasoning; tool calling \\
\multicolumn{8}{@{}l}{\textit{Granite} --- IBM (Granite) \href{https://github.com/ibm-granite/granite-code-models}{[family page]}} \\*
\href{https://huggingface.co/ibm-granite/granite-3.2-2b-instruct}{Granite 3.2-2B} & OW & 2B & 125K & Text & See source & 2025-02-26 & Thinking mode; 128K context \\
\href{https://huggingface.co/ibm-granite/granite-3.2-8b-instruct}{Granite 3.2-8B} & OW & 8B & 125K & Text & See source & 2025-02-26 & Thinking mode (chain-of-thought); 128K context \\
\multicolumn{8}{@{}l}{\textit{IBM Granite} --- IBM (Granite 3.3 (Apr 2025)) \href{https://huggingface.co/ibm-granite}{[family page]}} \\*
\href{https://huggingface.co/ibm-granite/granite-4.0-nano-350m}{Granite 4.0 Nano 350M} & OW & 350M & 128K & Text & Apache 2.0 & 2025-10-28 & on-device; browser \\
\href{https://huggingface.co/ibm-granite/granite-4.0-nano-1b}{Granite 4.0 Nano 1B} & OW & 1B & 128K & Text & Apache 2.0 & 2025-10-28 & on-device; browser \\
\href{https://huggingface.co/ibm-granite/granite-4.0-h-tiny}{Granite 4.0 H-Tiny} & OW & 7B (1B act.) & 128K & Text & Apache 2.0 & 2025-10-02 & low-latency; agentic \\
\href{https://huggingface.co/ibm-granite/granite-4.0-h-micro}{Granite 4.0 H-Micro} & OW & 3B & 128K & Text & Apache 2.0 & 2025-10-02 & edge; tool use \\
\href{https://huggingface.co/ibm-granite/granite-4.0-micro}{Granite 4.0 Micro} & OW & 3B & 128K & Text & Apache 2.0 & 2025-10-02 & edge; RAG \\
\href{https://huggingface.co/ibm-granite/granite-4.0-h-small}{Granite 4.0 H-Small} & OW & 32B (9B act.) & 128K & Text & Apache 2.0 & 2025-10-02 & enterprise; agentic \\
\href{https://huggingface.co/ibm-granite/granite-3.3-2b-instruct}{Granite 3.3-2B-Instruct} & OW & 2B & 128K & Text & Apache 2.0 & 2025-04 & enterprise NLP; governed AI \\
\href{https://huggingface.co/ibm-granite/granite-3.3-8b-instruct}{Granite 3.3-8B-Instruct} & OW & 8B & 128K & Text & Apache 2.0 & 2025-04 & enterprise reasoning; Fill-in-Middle code \\
\multicolumn{8}{@{}l}{\textit{Jan} --- Jan.ai (Jan-v1-4B (2025)) \href{https://huggingface.co/janhq/Jan-v1-4B}{[family page]}} \\*
\href{https://huggingface.co/janhq/Jan-v1-4B}{Jan-v1-4B} & OW & 4B & 128K & Text & Apache 2.0 & 2025 & agentic reasoning; tool use \\
\multicolumn{8}{@{}l}{\textit{LFM2} --- Liquid AI} \\*
\href{https://huggingface.co/LiquidAI/LFM2-8B-A1.5B}{LFM2-8B-A1.5B (MoE)} & OW & 8.3B (1.5B act.) & 32K & Text & LFM Open License v1.0 & 2025-08 & on-device; agentic \\
\href{https://huggingface.co/LiquidAI/LFM2-2.6B}{LFM2-2.6B} & OW & 2.6B & 32K & Text & LFM Open License v1.0 & 2025-08 & on-device \\
\href{https://huggingface.co/LiquidAI/LFM2-350M}{LFM2-350M} & OW & 350M & 32K & Text & LFM Open License v1.0 & 2025-07-10 & on-device; low latency \\
\href{https://huggingface.co/LiquidAI/LFM2-700M}{LFM2-700M} & OW & 700M & 32K & Text & LFM Open License v1.0 & 2025-07-10 & on-device \\
\href{https://huggingface.co/LiquidAI/LFM2-1.2B}{LFM2-1.2B} & OW & 1.2B & 32K & Text & LFM Open License v1.0 & 2025-07-10 & on-device; agentic \\
\multicolumn{8}{@{}l}{\textit{Llama 4} --- Meta AI (Llama 4 Maverick (Apr 2025)) \href{https://ai.meta.com/blog/llama-4-multimodal-intelligence/}{[family page]}} \\*
\href{https://huggingface.co/meta-llama/Llama-4-Maverick-17B-128E-Instruct}{Llama 4 Maverick} & Both & 400B (17B act.) & 10M & Text+Vision & Llama 4 Community License & 2025-04-05 & image understanding; coding \\
\href{https://huggingface.co/meta-llama/Llama-4-Scout-17B-16E-Instruct}{Llama 4 Scout} & Both & 109B (17B act.) & 10M & Text+Vision & Llama 4 Community License & 2025-04-05 & extreme long-context; single H100 runnable \\
\multicolumn{8}{@{}l}{\textit{Phi-4} --- Microsoft (Phi-4-reasoning-plus (Apr 2025)) \href{https://huggingface.co/microsoft/phi-4}{[family page]}} \\*
\href{https://huggingface.co/microsoft/Phi-4-mini-flash-reasoning}{Phi-4-mini-flash-reasoning} & OW & 3.8B & 128K & Text & MIT & 2025-04 & fast reasoning; edge agents \\
\href{https://huggingface.co/microsoft/Phi-4-mini-reasoning}{Phi-4-mini-reasoning} & OW & 3.8B & 128K & Text & MIT & 2025-04 & chain-of-thought reasoning; math \\
\href{https://huggingface.co/microsoft/Phi-4-reasoning}{Phi-4-reasoning} & OW & 14B & 128K & Text & MIT & 2025-04 & complex reasoning; STEM \\
\href{https://huggingface.co/microsoft/Phi-4-reasoning-plus}{Phi-4-reasoning-plus} & OW & 14B & 128K & Text & MIT & 2025-04 & RL-trained reasoning; math olympiad-level \\
\href{https://huggingface.co/microsoft/Phi-4-mini-instruct}{Phi-4-mini-instruct} & OW & 3.8B & 128K & Text & MIT & 2025-02 & on-device chat; function calling \\
\multicolumn{8}{@{}l}{\textit{Ministral} --- Mistral AI (Ministral-3B (Dec 2025)) \href{https://docs.mistral.ai/getting-started/models}{[family page]}} \\*
\href{https://huggingface.co/mistralai/Ministral-3-14B-Instruct-2512}{Ministral 3 (14B)} & OW & 14B & -- & Text+Vision & See source & 2025-12-02 & edge deployment; multimodal \\
\href{https://huggingface.co/mistralai/Ministral-3B-instruct-2512}{Ministral-3B-Instruct-2512} & OW & 3.8B & 256K & Text+Vision & Apache 2.0 & 2025-12 & tiny multimodal; function calling \\
\multicolumn{8}{@{}l}{\textit{Voxtral} --- Mistral AI (Voxtral) \href{https://mistral.ai/models}{[family page]}} \\*
\href{https://huggingface.co/mistralai/Mistral-Small-3.1-24B-Instruct-2503}{Voxtral Mini} & API & undisclosed & -- & Text+Audio & Proprietary & 2025-09-15 & speech understanding; transcription \\
\multicolumn{8}{@{}l}{\textit{Nemotron} --- NVIDIA (Nemotron) \href{https://arxiv.org/abs/2402.16819}{[family page]}} \\*
\href{https://huggingface.co/nvidia/NVIDIA-Nemotron-Nano-9B-v2}{NVIDIA Nemotron Nano 2} & OW & 9B & 128K & Text & See source & 2025-08-18 & on-device; efficient \\
\multicolumn{8}{@{}l}{\textit{GPT-OSS} --- OpenAI (gpt-oss-120b (2025)) \href{https://openai.com/index/introducing-gpt-oss/}{[family page]}} \\*
\href{https://huggingface.co/openai/gpt-oss-20b}{gpt-oss-20b} & OW & 21B (3.6B act.) & 128K & Text & Apache 2.0 & 2025 & edge reasoning; 16GB device inference \\
\href{https://huggingface.co/openai/gpt-oss-120b}{gpt-oss-120b} & OW & 117B (5.1B act.) & 128K & Text & Apache 2.0 & 2025 & open-weight reasoning; near o4-mini parity \\
\multicolumn{8}{@{}l}{\textit{MiniCPM} --- OpenBMB (MiniCPM-V 4.5 (2025)) \href{https://github.com/OpenBMB/MiniCPM-V}{[family page]}} \\*
\href{https://huggingface.co/openbmb/MiniCPM-V-4\_5}{MiniCPM-V 4.5} & OW & \textasciitilde{}8B & 32K & Text+Vision & Apache 2.0 & 2025 & mobile vision-language; on-device VQA \\
\href{https://github.com/OpenBMB/MiniCPM-o}{MiniCPM-o 3.0} & OW & \textasciitilde{}9B & 32K & Text+Vision+Audio & Apache 2.0 & 2025 & end-to-end multimodal on-device; live streaming \\
\multicolumn{8}{@{}l}{\textit{InternLM} --- Shanghai AI Laboratory (InternLM) \href{https://github.com/InternLM/InternLM}{[family page]}} \\*
\href{https://huggingface.co/internlm/internlm3-8b-instruct}{InternLM3-8B} & OW & 8B & 32K & Text & See source & 2025-01-09 & Reasoning via deep thinking; Knowledge-rich training pipeline \\
\multicolumn{8}{@{}l}{\textit{Falcon-H1} --- Technology Innovation Institute} \\*
\href{https://huggingface.co/tiiuae/Falcon-H1-0.5B-Instruct}{Falcon-H1-0.5B} & OW & 0.5B & 256K & Text & Falcon-LLM License 2.0 & 2025-07-31 & edge; long context \\
\href{https://huggingface.co/tiiuae/Falcon-H1-1.5B-Instruct}{Falcon-H1-1.5B} & OW & 1.5B & 256K & Text & Falcon-LLM License 2.0 & 2025-07-31 & edge; long context \\
\href{https://huggingface.co/tiiuae/Falcon-H1-3B-Instruct}{Falcon-H1-3B} & OW & 3B & 256K & Text & Falcon-LLM License 2.0 & 2025-07-31 & general; long context \\
\href{https://huggingface.co/tiiuae/Falcon-H1-7B-Instruct}{Falcon-H1-7B} & OW & 7B & 256K & Text & Falcon-LLM License 2.0 & 2025-07-31 & general; long context \\
\addlinespace[3pt]
\multicolumn{8}{@{}l}{\hypertarget{model-year-slm-2024}{\textbf{2024}}} \\ \midrule
\multicolumn{8}{@{}l}{\textit{Yi} --- 01.AI (Yi) \href{https://huggingface.co/01-ai/Yi-6B}{[family page]}} \\*
\href{https://huggingface.co/01-ai/Yi-1.5-6B}{Yi-1.5-6B} & OW & 6B & -- & Text & See source & 2024-05-13 & legacy \\
\href{https://huggingface.co/01-ai/Yi-1.5-9B}{Yi-1.5-9B} & OW & 9B & -- & Text & See source & 2024-05-13 & legacy \\
\multicolumn{8}{@{}l}{\textit{Qwen} --- Alibaba (Qwen) \href{https://github.com/QwenLM/Qwen}{[family page]}} \\*
\href{https://huggingface.co/Qwen/Qwen2.5-Coder-1.5B-Instruct}{Qwen2.5-Coder-1.5B} & OW & 1.5B & -- & Text & See source & 2024-11-12 & compact coding model \\
\href{https://huggingface.co/Qwen/Qwen2.5-14B}{Qwen2.5-14B} & OW & 14B & -- & Text & See source & 2024-09-19 & active \\
\href{https://huggingface.co/Qwen/Qwen2-VL-7B-Instruct}{Qwen2-VL-7B} & OW & 7B & -- & Text+Video+Image & See source & 2024-08-28 & vision-language; video understanding \\
\href{https://huggingface.co/Qwen/Qwen2-Audio-7B-Instruct}{Qwen2-Audio-7B} & OW & 7B & -- & Text+Audio & See source & 2024-08-09 & audio understanding; speech recognition \\
\href{https://huggingface.co/Qwen/Qwen2-0.5B}{Qwen2-0.5B} & OW & 0.5B & -- & Text & See source & 2024-06-07 & active \\
\href{https://huggingface.co/Qwen/Qwen2-1.5B}{Qwen2-1.5B} & OW & 1.5B & -- & Text & See source & 2024-06-07 & active \\
\href{https://huggingface.co/Qwen/Qwen2-7B}{Qwen2-7B} & OW & 7B & -- & Text & See source & 2024-06-07 & active \\
\href{https://huggingface.co/Qwen/Qwen1.5-MoE-A2.7B}{Qwen1.5-MoE-A2.7B} & OW & 14.3B & -- & Text & See source & 2024-03-28 & MoE architecture; efficient \\
\href{https://huggingface.co/Qwen/Qwen1.5-0.5B}{Qwen1.5-0.5B} & OW & 0.5B & -- & Text & See source & 2024-02-05 & legacy \\
\href{https://huggingface.co/Qwen/Qwen1.5-7B}{Qwen1.5-7B} & OW & 7B & -- & Text & See source & 2024-02-05 & legacy \\
\href{https://huggingface.co/Qwen/Qwen1.5-14B}{Qwen1.5-14B} & OW & 14B & -- & Text & See source & 2024-02-05 & legacy \\
\multicolumn{8}{@{}l}{\textit{Qwen2.5} --- Alibaba (Qwen2.5 (Sep 2024)) \href{https://huggingface.co/Qwen}{[family page]}} \\*
\href{https://huggingface.co/Qwen/Qwen2.5-0.5B-Instruct}{Qwen2.5-0.5B-Instruct} & OW & 0.5B & 32K & Text & Apache 2.0 & 2024-09-19 & extreme edge/IoT; classification \\
\href{https://huggingface.co/Qwen/Qwen2.5-1.5B-Instruct}{Qwen2.5-1.5B-Instruct} & OW & 1.5B & 32K & Text & Apache 2.0 & 2024-09-19 & edge NLP; lightweight assistants \\
\href{https://huggingface.co/Qwen/Qwen2.5-3B-Instruct}{Qwen2.5-3B-Instruct} & OW & 3B & 32K & Text & Apache 2.0 & 2024-09-19 & mobile chat; summarization \\
\href{https://huggingface.co/Qwen/Qwen2.5-7B-Instruct}{Qwen2.5-7B-Instruct} & OW & 7B & 128K & Text & Apache 2.0 & 2024-09-19 & general assistant; coding \\
\multicolumn{8}{@{}l}{\textit{Qwen2.5-Coder} --- Alibaba (Qwen2.5-Coder (Nov 2024)) \href{https://huggingface.co/Qwen/Qwen2.5-Coder-7B-Instruct}{[family page]}} \\*
\href{https://huggingface.co/Qwen/Qwen2.5-Coder-7B-Instruct}{Qwen2.5-Coder-7B-Instruct} & OW & 7B & 128K & Text+Code & Apache 2.0 & 2024-11 & code generation; debugging \\
\href{https://huggingface.co/Qwen/Qwen2.5-Coder-14B-Instruct}{Qwen2.5-Coder-14B-Instruct} & OW & 14B & 128K & Text+Code & Apache 2.0 & 2024-11 & advanced code assistant; agentic code tasks \\
\multicolumn{8}{@{}l}{\textit{OLMo} --- Allen Institute for AI (Ai2) (OLMo) \href{https://allenai.org/olmo}{[family page]}} \\*
\href{https://huggingface.co/allenai/OLMo-2-1124-7B}{OLMo 2-7B} & OW & 7B & -- & Text & See source & 2024-11-26 & fully open; training code included \\
\href{https://huggingface.co/allenai/OLMo-2-1124-13B}{OLMo 2-13B} & OW & 13B & -- & Text & See source & 2024-11-26 & active \\
\href{https://huggingface.co/allenai/OLMo-1B}{OLMo-1B} & OW & 1B & 2K & Text & See source & 2024-02-01 & Fully open (data, code, weights, logs); Dolma training data \\
\href{https://huggingface.co/allenai/OLMo-7B}{OLMo-7B} & OW & 7B & -- & Text & See source & 2024-02-01 & fully open; training data/code/logs released \\
\multicolumn{8}{@{}l}{\textit{Tülu} --- Allen Institute for AI (Ai2) (Tülu) \href{https://allenai.org/blog/tulu-3-technical}{[family page]}} \\*
\href{https://huggingface.co/allenai/Llama-3.1-Tulu-3-8B}{Tülu 3-8B} & OW & 8B & -- & Text & See source & 2024-11-22 & instruction following; DPO \\
\multicolumn{8}{@{}l}{\textit{Amazon Nova} --- Amazon (Nova Premier (2025)) \href{https://aws.amazon.com/nova/models/}{[family page]}} \\*
\href{https://aws.amazon.com/nova/models/}{Amazon Nova Lite} & API & Undisclosed & 300K & Text+Vision & Proprietary (AWS) & 2024-12-03 & multimodal understanding; cost-efficient AWS inference \\
\href{https://aws.amazon.com/nova/models/}{Amazon Nova Micro} & API & Undisclosed & 128K & Text & Proprietary (AWS) & 2024-12-03 & AWS-native low-latency text tasks; Lambda AI apps \\
\multicolumn{8}{@{}l}{\textit{AFM} --- Apple (AFM) \href{https://machinelearning.apple.com/research/introducing-apple-foundation-models}{[family page]}} \\*
\href{https://machinelearning.apple.com/research/introducing-apple-foundation-models}{Apple Foundation Model (On-Device)} & API & \textasciitilde{}3B & 4K & Text & Proprietary & 2024-06-10 & On-device Apple Intelligence; Privacy-preserving \\
\multicolumn{8}{@{}l}{\textit{OpenELM} --- Apple (OpenELM) \href{https://machinelearning.apple.com/research/openelm}{[family page]}} \\*
\href{https://huggingface.co/apple/OpenELM-3B}{OpenELM} & OW & 3B & -- & Text & See source & 2024-04-24 & layer-wise scaling; on-device \\
\multicolumn{8}{@{}l}{\textit{Aya} --- Cohere (Aya) \href{https://cohere.com/research/aya}{[family page]}} \\*
\href{https://huggingface.co/CohereForAI/aya-expanse-8b}{Aya Expanse-8B} & OW & 8B & 8K & Text & See source & 2024-10-22 & Improved multilingual performance; Data arbitrage and synthetic data strategies \\
\href{https://huggingface.co/CohereForAI/aya-23-8B}{Aya 23-8B} & OW & 8B & 8K & Text & See source & 2024-05-23 & 23 language support; Multilingual instruction following \\
\multicolumn{8}{@{}l}{\textit{Cohere Command} --- Cohere (Command A (Mar 2025)) \href{https://cohere.com/blog/command-a}{[family page]}} \\*
\href{https://huggingface.co/CohereForAI/c4ai-command-r7b-12-2024}{Command R7B} & Both & 7B & 128K & Text & Apache 2.0 & 2024-12 & RAG; multilingual enterprise tasks \\
\multicolumn{8}{@{}l}{\textit{TinyLlama} --- Community (TinyLlama-1.1B (2024)) \href{https://huggingface.co/TinyLlama/TinyLlama-1.1B-Chat-v1.0}{[family page]}} \\*
\href{https://huggingface.co/TinyLlama/TinyLlama-1.1B-Chat-v1.0}{TinyLlama-1.1B-Chat} & OW & 1.1B & 2K & Text & Apache 2.0 & 2024 & ultra-light on-device chat; prototyping \\
\multicolumn{8}{@{}l}{\textit{DeepSeek} --- DeepSeek (DeepSeek) \href{https://github.com/deepseek-ai/DeepSeek-Coder}{[family page]}} \\*
\href{https://huggingface.co/deepseek-ai/deepseek-vl-7b-chat}{DeepSeek-VL-7B} & OW & 7B & 4K & Text+Vision & See source & 2024-03-13 & Vision-language model; Hybrid vision encoder \\
\href{https://huggingface.co/deepseek-ai/deepseek-math-7b-instruct}{DeepSeek-Math-7B} & OW & 7B & 4K & Text & See source & 2024-02-05 & Math-specialized; Group Relative Policy Optimization (GRPO) \\
\multicolumn{8}{@{}l}{\textit{Gemma} --- Google DeepMind (Gemma) \href{https://blog.google/technology/developers/gemma-open-models/}{[family page]}} \\*
\href{https://huggingface.co/google/gemma-2-2b}{Gemma 2-2B} & OW & 2B & 8K & Text & See source & 2024-07-31 & Knowledge distillation from larger model; Sliding window attention \\
\href{https://huggingface.co/google/gemma-2-9b}{Gemma 2-9B} & OW & 9B & -- & Text & See source & 2024-06-27 & improved efficiency; knowledge distillation \\
\href{https://huggingface.co/google/paligemma-3b-pt-224}{PaliGemma} & OW & 3B & 8K & Text+Vision & See source & 2024-05-14 & Vision-language model; SigLIP vision encoder + Gemma 2B \\
\href{https://huggingface.co/google/recurrentgemma-2b}{RecurrentGemma-2B} & OW & 2B & 8K & Text & See source & 2024-04-09 & Griffin architecture with linear recurrence; Efficient inference \\
\href{https://huggingface.co/google/codegemma-7b}{CodeGemma-7B} & OW & 7B & -- & Text & See source & 2024-04-09 & code generation; code completion \\
\href{https://huggingface.co/google/gemma-2b}{Gemma-2B} & OW & 2B & -- & Text & See source & 2024-02-21 & legacy \\
\href{https://huggingface.co/google/gemma-7b}{Gemma-7B} & OW & 7B & -- & Text & See source & 2024-02-21 & built from Gemini technology; efficient \\
\multicolumn{8}{@{}l}{\textit{SmolLM} --- Hugging Face (SmolLM3-3B (2025)) \href{https://huggingface.co/HuggingFaceTB/SmolLM3-3B}{[family page]}} \\*
\href{https://huggingface.co/HuggingFaceTB/SmolLM-135M}{SmolLM-135M} & OW & 135M & 2K & Text & See source & 2024-07-17 & Ultra-small LM; On-device capable \\
\href{https://huggingface.co/HuggingFaceTB/SmolLM-1.7B}{SmolLM} & OW & 1.7B & -- & Text & See source & 2024-07-16 & tiny but capable; on-device \\
\href{https://huggingface.co/HuggingFaceTB/SmolLM2-135M}{SmolLM2-135M} & OW & 135M & 2K & Text & Apache 2.0 & 2024 & ultra-micro inference; browser/WASM \\
\href{https://huggingface.co/HuggingFaceTB/SmolLM2-360M-Instruct}{SmolLM2-360M-Instruct} & OW & 360M & 8K & Text & Apache 2.0 & 2024 & on-device/browser inference \\
\href{https://huggingface.co/HuggingFaceTB/SmolLM2-1.7B-Instruct}{SmolLM2-1.7B-Instruct} & OW & 1.7B & 8K & Text & Apache 2.0 & 2024 & efficient on-device chat; coding assist \\
\multicolumn{8}{@{}l}{\textit{LLaMA} --- Meta AI (LLaMA) \href{https://arxiv.org/abs/2302.13971}{[family page]}} \\*
\href{https://huggingface.co/meta-llama/Meta-Llama-3-8B}{Llama-3-8B} & OW & 8B & -- & Text & See source & 2024-04-18 & improved tokenizer; grouped query attention \\
\multicolumn{8}{@{}l}{\textit{Llama 3.1} --- Meta AI (Llama 3.1 (Jul 2024)) \href{https://huggingface.co/meta-llama}{[family page]}} \\*
\href{https://huggingface.co/meta-llama/Meta-Llama-3.1-8B-Instruct}{Llama 3.1-8B-Instruct} & OW & 8B & 128K & Text & Llama 3.1 Community License & 2024-07-23 & most fine-tuned SLM base globally; RAG \\
\multicolumn{8}{@{}l}{\textit{Llama 3.2} --- Meta AI (Llama 3.2 (Sep 2024)) \href{https://ai.meta.com/blog/llama-3-2-connect-2024-vision-edge-mobile-devices/}{[family page]}} \\*
\href{https://huggingface.co/meta-llama/Llama-3.2-1B-Instruct}{Llama 3.2-1B-Instruct} & OW & 1B & 128K & Text & Llama 3.2 Community License & 2024-09-25 & ultra-small on-device; Android/iOS apps \\
\href{https://huggingface.co/meta-llama/Llama-3.2-3B-Instruct}{Llama 3.2-3B-Instruct} & OW & 3B & 128K & Text & Llama 3.2 Community License & 2024-09-25 & mobile AI; default on-device Qualcomm \\
\href{https://huggingface.co/meta-llama/Llama-3.2-11B-Vision-Instruct}{Llama 3.2-11B-Vision-Instruct} & OW & 11B & 128K & Text+Vision & Llama 3.2 Community License & 2024-09-25 & vision-language tasks; document understanding \\
\multicolumn{8}{@{}l}{\textit{BitNet} --- Microsoft (BitNet) \href{https://arxiv.org/abs/2402.17764}{[family page]}} \\*
\href{https://huggingface.co/microsoft/bitnet-b1.58-2B-4T}{BitNet b1.58} & OW & 3B & -- & Text & See source & 2024-02-27 & 1-bit LLM (ternary weights: -1, 0, 1); Extreme efficiency \\
\multicolumn{8}{@{}l}{\textit{Phi-3} --- Microsoft (Phi-3.5 (Aug 2024)) \href{https://azure.microsoft.com/en-us/products/phi}{[family page]}} \\*
\href{https://huggingface.co/microsoft/Phi-3.5-mini-instruct}{Phi-3.5-mini-instruct} & OW & 3.8B & 128K & Text & MIT & 2024-08 & multilingual instruction following \\
\href{https://huggingface.co/microsoft/Phi-3-mini-128k-instruct}{Phi-3-mini-128k-instruct} & OW & 3.8B & 128K & Text & MIT & 2024-04-23 & on-device chat; instruction following \\
\href{https://huggingface.co/microsoft/Phi-3-small-128k-instruct}{Phi-3-small-128k-instruct} & OW & 7B & 128K & Text & MIT & 2024-04-23 & balanced reasoning \& coding \\
\href{https://huggingface.co/microsoft/Phi-3-medium-128k-instruct}{Phi-3-medium-128k-instruct} & OW & 14B & 128K & Text & MIT & 2024-04-23 & on-premise deployment; complex reasoning \\
\multicolumn{8}{@{}l}{\textit{Phi-4} --- Microsoft (Phi-4-reasoning-plus (Apr 2025)) \href{https://huggingface.co/microsoft/phi-4}{[family page]}} \\*
\href{https://huggingface.co/microsoft/phi-4}{Phi-4} & OW & 14B & 16K & Text & MIT & 2024-12-12 & STEM reasoning; math \\
\multicolumn{8}{@{}l}{\textit{Ministral} --- Mistral AI (Ministral-3B (Dec 2025)) \href{https://docs.mistral.ai/getting-started/models}{[family page]}} \\*
\href{https://huggingface.co/mistralai/Ministral-8B-Instruct-2410}{Ministral-8B-Instruct-2410} & OW & 8B & 128K & Text & Apache 2.0 & 2024-10 & edge-friendly assistant; function calling \\
\multicolumn{8}{@{}l}{\textit{Mistral Small} --- Mistral AI (Mistral Small 4 (Mar 2026)) \href{https://mistral.ai/news/mistral-small-4}{[family page]}} \\*
\href{https://huggingface.co/mistralai/Mistral-Nemo-Instruct-2407}{Mistral NeMo 12B} & OW & 12B & 128K & Text & Apache 2.0 & 2024-07 & enterprise RAG; long-document summarization \\
\multicolumn{8}{@{}l}{\textit{Pixtral} --- Mistral AI (Pixtral) \href{https://mistral.ai/news/pixtral-12b/}{[family page]}} \\*
\href{https://huggingface.co/mistralai/Pixtral-12B-2409}{Pixtral-12B} & OW & 12B & -- & Text+Image & See source & 2024-09-17 & vision-language; variable image resolution \\
\multicolumn{8}{@{}l}{\textit{Nemotron} --- NVIDIA (Nemotron) \href{https://arxiv.org/abs/2402.16819}{[family page]}} \\*
\href{https://huggingface.co/nvidia/Minitron-8B-Base}{Minitron-8B} & OW & 8B & 8K & Text & See source & 2024-07-17 & Width and depth pruning; Knowledge distillation \\
\href{https://huggingface.co/nvidia/Nemotron-4-15B-Base}{Nemotron-4-15B} & OW & 15B & 8K & Text & See source & 2024-03-14 & High-quality training data; Strong downstream task performance \\
\multicolumn{8}{@{}l}{\textit{Hermes} --- Nous Research (Hermes) \href{https://huggingface.co/NousResearch/Nous-Hermes-2-Mixtral-8x7B-DPO}{[family page]}} \\*
\href{https://huggingface.co/NousResearch/Nous-Hermes-2-Mixtral-8x7B-DPO}{Nous Hermes 2 Mixtral} & OW & 46.7B (12.9B act.) & 32K & Text & See source & 2024-01-17 & Fine-tuned Mixtral-8x7B; Function calling \\
\multicolumn{8}{@{}l}{\textit{GPT} --- OpenAI (GPT) \href{https://openai.com/research/instruction-following}{[family page]}} \\*
\href{https://openai.com/index/gpt-4o-mini-advancing-cost-efficient-intelligence/}{GPT-4o mini} & API & undisclosed & -- & Text+Vision & Proprietary & 2024-07-18 & cost-efficient; fast \\
\multicolumn{8}{@{}l}{\textit{Reka} --- Reka AI (Reka) \href{https://reka.ai/news/reka-flash-efficient-and-capable-multimodal-language-models}{[family page]}} \\*
\href{https://huggingface.co/RekaAI/reka-edge-2603}{Reka Edge} & API & 7B & -- & Text+Vision & Proprietary & 2024-02-12 & compact; multimodal \\
\multicolumn{8}{@{}l}{\textit{InternLM} --- Shanghai AI Laboratory (InternLM) \href{https://github.com/InternLM/InternLM}{[family page]}} \\*
\href{https://huggingface.co/internlm/internlm2\_5-7b-chat}{InternLM2.5-7B} & OW & 7B & -- & Text & See source & 2024-07-03 & 1M context; tool use \\
\href{https://huggingface.co/internlm/internlm2\_5-20b-chat}{InternLM2.5-20B} & OW & 20B & 1M & Text & See source & 2024-07-03 & 1M context with LMDeploy; Strong reasoning \\
\href{https://huggingface.co/internlm/internlm2-20b}{InternLM2-20B} & OW & 20B & -- & Text & See source & 2024-01-17 & 200K context; tool use \\
\multicolumn{8}{@{}l}{\textit{StableLM} --- Stability AI (StableLM 2 (Jan 2024)) \href{https://huggingface.co/stabilityai/stablelm-2-1\_6b}{[family page]}} \\*
\href{https://huggingface.co/stabilityai/stablelm-2-12b}{StableLM 2 12B} & OW & 12B & 4K & Text & See source & 2024-02-20 & Multilingual (7 languages); Strong code capability \\
\href{https://huggingface.co/stabilityai/stablelm-2-1\_6b}{StableLM 2 1.6B} & OW & 1.6B & 4K & Text & Stability AI License & 2024-01-19 & lightweight creative writing; basic chat \\
\multicolumn{8}{@{}l}{\textit{Falcon} --- Technology Innovation Institute (TII) (Falcon 180B (Sep 2023)) \href{https://huggingface.co/tiiuae/falcon-180B}{[family page]}} \\*
\href{https://huggingface.co/tiiuae/Falcon3-10B-Instruct}{Falcon 3} & OW & 10B & -- & Text+Audio+Video+Image & See source & 2024-12-12 & multimodal; runs on laptops \\
\href{https://huggingface.co/tiiuae/falcon-mamba-7b}{Falcon Mamba-7B} & OW & 7B & -- & Text & See source & 2024-08-12 & SSM architecture; constant memory usage \\
\href{https://huggingface.co/tiiuae/falcon-11B}{Falcon 2-11B} & OW & 11B & -- & Text & See source & 2024-05-14 & multimodal variant; vision-to-language \\
\multicolumn{8}{@{}l}{\textit{GLM} --- Zhipu AI (GLM-5 (Feb 2026)) \href{https://github.com/THUDM/GLM-4}{[family page]}} \\*
\href{https://huggingface.co/THUDM/glm-4-9b}{GLM-4-9B} & OW & 9B & -- & Text & See source & 2024-06-06 & 128K-1M context; tool use \\
\href{https://huggingface.co/THUDM/cogvlm2-llama3-chat-19B}{CogVLM2-19B} & OW & 19B & 8K & Text+Video+Image & See source & 2024-05-24 & Visual understanding; Video understanding \\
\addlinespace[3pt]
\multicolumn{8}{@{}l}{\hypertarget{model-year-slm-2023-and-earlier}{\textbf{2023 and earlier}}} \\ \midrule
\multicolumn{8}{@{}l}{\textit{Yi} --- 01.AI (Yi) \href{https://huggingface.co/01-ai/Yi-6B}{[family page]}} \\*
\href{https://huggingface.co/01-ai/Yi-6B}{Yi-6B} & OW & 6B & -- & Text & See source & 2023-11-02 & bilingual; strong reasoning \\
\multicolumn{8}{@{}l}{\textit{Fuyu} --- Adept AI (Fuyu) \href{https://www.adept.ai/blog/fuyu-8b}{[family page]}} \\*
\href{https://huggingface.co/adept/fuyu-8b}{Fuyu-8B} & OW & 8B & -- & Text+Vision & See source & 2023-10-17 & Simplified multimodal architecture; No separate vision encoder \\
\multicolumn{8}{@{}l}{\textit{Qwen} --- Alibaba (Qwen) \href{https://github.com/QwenLM/Qwen}{[family page]}} \\*
\href{https://huggingface.co/Qwen/Qwen-Audio}{Qwen-Audio} & OW & 7B & -- & Text+Audio & See source & 2023-11-30 & Audio understanding; Multi-task audio processing \\
\href{https://huggingface.co/Qwen/Qwen-14B}{Qwen-14B} & OW & 14B & -- & Text & See source & 2023-09-25 & legacy \\
\href{https://huggingface.co/Qwen/Qwen-VL}{Qwen-VL} & OW & 9.6B & -- & Text+Image & See source & 2023-08-24 & vision-language; image understanding \\
\href{https://huggingface.co/Qwen/Qwen-7B}{Qwen-7B} & OW & 7B & -- & Text & See source & 2023-08-03 & bilingual; tool use \\
\multicolumn{8}{@{}l}{\textit{Baichuan} --- Baichuan Inc. (Baichuan) \href{https://github.com/baichuan-inc/Baichuan-7B}{[family page]}} \\*
\href{https://huggingface.co/baichuan-inc/Baichuan2-13B-Chat}{Baichuan-2-13B} & OW & 13B & -- & Text & See source & 2023-09-06 & legacy \\
\href{https://huggingface.co/baichuan-inc/Baichuan-13B-Chat}{Baichuan-13B} & OW & 13B & -- & Text & See source & 2023-07-11 & legacy \\
\href{https://huggingface.co/baichuan-inc/Baichuan-7B}{Baichuan-7B} & OW & 7B & -- & Text & See source & 2023-06-15 & bilingual; strong Chinese performance \\
\multicolumn{8}{@{}l}{\textit{Cerebras-GPT} --- Cerebras Systems (Cerebras-GPT) \href{https://huggingface.co/cerebras/Cerebras-GPT-13B}{[family page]}} \\*
\href{https://huggingface.co/cerebras/Cerebras-GPT-13B}{Cerebras-GPT-13B} & OW & 13B & 2K & Text & See source & 2023-03-28 & Chinchilla compute-optimal training; Trained on Cerebras CS-2 hardware \\
\multicolumn{8}{@{}l}{\textit{DeepSeek} --- DeepSeek (DeepSeek) \href{https://github.com/deepseek-ai/DeepSeek-Coder}{[family page]}} \\*
\href{https://huggingface.co/deepseek-ai/deepseek-coder-1.3b-instruct}{DeepSeek-Coder-1.3B} & OW & 1.3B & -- & Text & See source & 2023-11-02 & compact code model; fill-in-the-middle \\
\href{https://huggingface.co/deepseek-ai/deepseek-coder-6.7b-instruct}{DeepSeek-Coder-6.7B} & OW & 6.7B & -- & Text & See source & 2023-11-02 & code generation; fill-in-the-middle \\
\multicolumn{8}{@{}l}{\textit{Pythia} --- EleutherAI (Pythia) \href{https://github.com/EleutherAI/pythia}{[family page]}} \\*
\href{https://huggingface.co/EleutherAI/pythia-12b}{Pythia-12B} & OW & 12B & 2K & Text & See source & 2023-04-03 & 154 intermediate checkpoints; Reproducible training \\
\multicolumn{8}{@{}l}{\textit{T5} --- Google (T5) \href{https://arxiv.org/abs/2210.11416}{[family page]}} \\*
\href{https://huggingface.co/google/flan-t5-xxl}{Flan-T5} & OW & 11B & -- & Text & See source & 2022-10-20 & instruction tuned; encoder-decoder \\
\multicolumn{8}{@{}l}{\textit{Gemini} --- Google DeepMind (Gemini) \href{https://deepmind.google/technologies/gemini/}{[family page]}} \\*
\href{https://deepmind.google/technologies/gemini/}{Gemini 1.0 Nano} & API & 3.25B & -- & Text & Proprietary & 2023-12-06 & on-device; efficient \\
\multicolumn{8}{@{}l}{\textit{StarCoder} --- Hugging Face / BigCode (StarCoder) \href{https://huggingface.co/bigcode/starcoder}{[family page]}} \\*
\href{https://huggingface.co/bigcode/starcoder}{StarCoder} & OW & 15.5B & -- & Text & See source & 2023-05-04 & code generation; 80+ programming languages \\
\multicolumn{8}{@{}l}{\textit{Jais} --- Inception / Core42 / Cerebras (Jais) \href{https://arxiv.org/abs/2308.16149}{[family page]}} \\*
\href{https://huggingface.co/core42/jais-13b-chat}{Jais-13B} & OW & 13B & 2K & Text & See source & 2023-08-30 & Best Arabic LLM at launch; Bilingual Arabic-English \\
\multicolumn{8}{@{}l}{\textit{Vicuna} --- LMSYS (Vicuna) \href{https://lmsys.org/blog/2023-03-30-vicuna/}{[family page]}} \\*
\href{https://huggingface.co/lmsys/vicuna-7b-v1.5}{Vicuna-7B v1.5} & OW & 7B & -- & Text & See source & 2023-07-29 & fine-tuned on ShareGPT; chat assistant \\
\href{https://huggingface.co/lmsys/vicuna-13b-delta-v0}{Vicuna-13B} & OW & 13B & -- & Text & See source & 2023-03-30 & fine-tuned on ShareGPT data; chat assistant \\
\multicolumn{8}{@{}l}{\textit{Code Llama} --- Meta AI (Code Llama) \href{https://ai.meta.com/blog/code-llama-large-language-model-coding/}{[family page]}} \\*
\href{https://huggingface.co/codellama/CodeLlama-7b-hf}{Code Llama-7B} & OW & 7B & -- & Text & See source & 2023-08-24 & code generation; infilling \\
\multicolumn{8}{@{}l}{\textit{LLaMA} --- Meta AI (LLaMA) \href{https://arxiv.org/abs/2302.13971}{[family page]}} \\*
\href{https://huggingface.co/meta-llama/Llama-2-7b-hf}{Llama-2-7B} & OW & 7B & -- & Text & See source & 2023-07-18 & open weights; commercial use allowed \\
\href{https://huggingface.co/meta-llama/Llama-2-13b-hf}{Llama-2-13B} & OW & 13B & -- & Text & See source & 2023-07-18 & legacy \\
\href{https://huggingface.co/huggyllama/llama-7b}{LLaMA-7B} & OW & 7B & -- & Text & Research only & 2023-02-24 & efficient training; strong few-shot performance \\
\href{https://huggingface.co/huggyllama/llama-13b}{LLaMA-13B} & OW & 13B & -- & Text & Research only & 2023-02-24 & legacy \\
\multicolumn{8}{@{}l}{\textit{Phi} --- Microsoft (Phi) \href{https://arxiv.org/abs/2306.11644}{[family page]}} \\*
\href{https://huggingface.co/microsoft/phi-2}{Phi-2} & OW & 2.7B & -- & Text & See source & 2023-12-13 & language comprehension; reasoning \\
\href{https://huggingface.co/microsoft/phi-1\_5}{Phi-1.5} & OW & 1.3B & -- & Text & See source & 2023-09-11 & reasoning; language understanding \\
\href{https://huggingface.co/microsoft/phi-1}{Phi-1} & OW & 1.3B & -- & Text & See source & 2023-06-20 & Python coding; textbook-quality data \\
\multicolumn{8}{@{}l}{\textit{Mistral} --- Mistral AI (Mistral) \href{https://docs.mistral.ai/getting-started/models}{[family page]}} \\*
\href{https://huggingface.co/mistralai/Mistral-7B-Instruct-v0.3}{Mistral Tiny (Mistral-7B-Instruct)} & API & 7B & 32K & Text & Proprietary & 2023-12-11 & API-served Mistral 7B Instruct; Cost-effective API option \\
\multicolumn{8}{@{}l}{\textit{Mistral Small} --- Mistral AI (Mistral Small 4 (Mar 2026)) \href{https://mistral.ai/news/mistral-small-4}{[family page]}} \\*
\href{https://huggingface.co/mistralai/Mistral-7B-Instruct-v0.3}{Mistral 7B v0.3} & OW & 7.3B & 32K & Text & Apache 2.0 & 2023-09-27 & general chat; RAG \\
\multicolumn{8}{@{}l}{\textit{MPT} --- MosaicML (Databricks) (MPT) \href{https://www.mosaicml.com/blog/mpt-7b}{[family page]}} \\*
\href{https://huggingface.co/mosaicml/mpt-7b}{MPT-7B} & OW & 7B & -- & Text & See source & 2023-05-05 & ALiBi positional encoding; long context \\
\multicolumn{8}{@{}l}{\textit{InternLM} --- Shanghai AI Laboratory (InternLM) \href{https://github.com/InternLM/InternLM}{[family page]}} \\*
\href{https://huggingface.co/internlm/internlm-20b}{InternLM-20B} & OW & 20B & -- & Text & See source & 2023-09-20 & legacy \\
\href{https://huggingface.co/internlm/internlm-7b}{InternLM-7B} & OW & 7B & -- & Text & See source & 2023-07-06 & bilingual; long context \\
\multicolumn{8}{@{}l}{\textit{StableLM} --- Stability AI (StableLM 2 (Jan 2024)) \href{https://huggingface.co/stabilityai/stablelm-2-1\_6b}{[family page]}} \\*
\href{https://huggingface.co/stabilityai/stablelm-zephyr-3b}{StableLM-Zephyr-3B} & OW & 3B & -- & Text & See source & 2023-11-29 & chat assistant; compact \\
\href{https://huggingface.co/stabilityai/stablelm-tuned-alpha-7b}{StableLM-Alpha} & OW & 7B & -- & Text & See source & 2023-04-19 & open-source; conversation \\
\multicolumn{8}{@{}l}{\textit{Falcon} --- Technology Innovation Institute (TII) (Falcon 180B (Sep 2023)) \href{https://huggingface.co/tiiuae/falcon-180B}{[family page]}} \\*
\href{https://huggingface.co/tiiuae/falcon-7b}{Falcon-7B} & OW & 7B & -- & Text & See source & 2023-06-05 & RefinedWeb dataset; high-quality pretraining \\
\multicolumn{8}{@{}l}{\textit{RedPajama} --- Together AI (RedPajama) \href{https://www.together.ai/blog/redpajama}{[family page]}} \\*
\href{https://huggingface.co/togethercomputer/RedPajama-INCITE-7B-Instruct}{RedPajama-INCITE-7B} & OW & 7B & 2K & Text & See source & 2023-05-05 & Fully open training data; Reproducible \\
\multicolumn{8}{@{}l}{\textit{SOLAR} --- Upstage (SOLAR) \href{https://arxiv.org/abs/2312.15166}{[family page]}} \\*
\href{https://huggingface.co/upstage/SOLAR-10.7B-Instruct-v1.0}{SOLAR-10.7B} & OW & 10.7B & 4K & Text & See source & 2023-12-13 & Depth Up-Scaling (DUS); Merged and scaled from Llama 2 architecture \\
\multicolumn{8}{@{}l}{\textit{GLM} --- Zhipu AI (GLM-5 (Feb 2026)) \href{https://github.com/THUDM/GLM-4}{[family page]}} \\*
\href{https://huggingface.co/THUDM/chatglm3-6b}{ChatGLM3-6B} & OW & 6B & -- & Text & See source & 2023-10-27 & tool use; code interpreter \\
\href{https://huggingface.co/THUDM/chatglm2-6b}{ChatGLM2-6B} & OW & 6B & -- & Text & See source & 2023-06-25 & legacy \\
\href{https://huggingface.co/THUDM/chatglm-6b}{ChatGLM-6B} & OW & 6B & -- & Text & See source & 2023-03-14 & bilingual; quantized for consumer GPUs \\
\end{longtable}
}
\end{landscape}


\subsection{Monitoring Sources and Maintenance}
\label{subsec:ecosystem_monitoring}
\begin{table}[ht]
\centering
\caption{High-value monitoring sources retained from the inventory files.}
\label{tab:ecosystem_sources}
\footnotesize
\begin{tabularx}{\textwidth}{L{4cm}X}
\toprule
\textbf{Source} & \textbf{Use} \\
\midrule
Hugging Face model hub and leaderboards & New releases, model cards, and benchmark submissions \\
Artificial Analysis & Context windows, pricing, and comparative capability snapshots \\
LLM Stats & Consolidated model metadata and API comparisons \\
LMSYS / Chatbot Arena & Human preference rankings \\
BentoML and vendor engineering blogs & Practical deployment write-ups and model updates \\
Official vendor changelogs & Release confirmation for OpenAI, Anthropic, Google, Meta, Mistral, Alibaba, and DeepSeek \\
\bottomrule
\end{tabularx}
\end{table}

Recommended maintenance queries from the inventory drafts included combinations such as ``best small language model [year]'', ``new open-weight reasoning model'', and vendor-specific searches around context window, model card, and release notes. In practice, the highest-signal workflow is to combine vendor documentation with model cards and a benchmark tracker rather than relying on any single static list.


\printjournalbibliography

\end{document}
